\documentclass[11pt,a4paper]{amsart}

\usepackage{amsmath,amssymb,amsthm,mathtools}
\usepackage[T1]{fontenc}
\usepackage[utf8]{inputenc}
\usepackage{lmodern}
\usepackage{microtype}
\usepackage[hidelinks]{hyperref}
\hypersetup{
  pdftitle={Volterra-Perron Pricing of the Option Manifold},
  pdfauthor={J. Vidal Llaurado},
  pdfsubject={Risk-neutral Volterra-Perron pricing, option manifolds, latent-state recovery},
  pdfkeywords={risk-neutral pricing, Volterra operators, Perron stress, option manifolds, latent-state recovery, residual kernels, synchronization Greeks}
}
\usepackage{orcidlink}
\usepackage{enumitem}
\usepackage{booktabs}
\usepackage[margin=1in]{geometry}
\raggedbottom

\theoremstyle{plain}
\newtheorem{theorem}{Theorem}[section]
\newtheorem{proposition}[theorem]{Proposition}
\newtheorem{lemma}[theorem]{Lemma}
\newtheorem{corollary}[theorem]{Corollary}

\theoremstyle{definition}
\newtheorem{definition}[theorem]{Definition}
\newtheorem{assumption}[theorem]{Assumption}
\newtheorem{condition}[theorem]{Condition}

\theoremstyle{remark}
\newtheorem{remark}[theorem]{Remark}

\DeclareMathOperator{\Var}{Var}
\DeclareMathOperator{\Cov}{Cov}
\DeclareMathOperator{\Span}{span}
\DeclareMathOperator{\Ker}{Ker}
\DeclareMathOperator{\Ran}{Ran}
\DeclareMathOperator{\dist}{dist}
\DeclareMathOperator{\Id}{Id}
\DeclareMathOperator{\Law}{Law}
\DeclareMathOperator{\argmin}{arg\,min}
\newcommand{\R}{\mathbb{R}}
\newcommand{\N}{\mathbb{N}}
\newcommand{\PP}{\mathbb{P}}
\newcommand{\QQ}{\mathbb{Q}}
\newcommand{\EE}{\mathbb{E}}
\newcommand{\cA}{\mathcal{A}}
\newcommand{\cB}{\mathcal{B}}
\newcommand{\cC}{\mathcal{C}}
\newcommand{\cD}{\mathcal{D}}
\newcommand{\cE}{\mathcal{E}}
\newcommand{\cF}{\mathcal{F}}
\newcommand{\cG}{\mathcal{G}}
\newcommand{\cH}{\mathcal{H}}
\newcommand{\cI}{\mathcal{I}}
\newcommand{\cJ}{\mathcal{J}}
\newcommand{\cK}{\mathcal{K}}
\newcommand{\cL}{\mathcal{L}}
\newcommand{\cM}{\mathcal{M}}
\newcommand{\cN}{\mathcal{N}}
\newcommand{\cO}{\mathcal{O}}
\newcommand{\cP}{\mathcal{P}}
\newcommand{\cQ}{\mathcal{Q}}
\newcommand{\cR}{\mathcal{R}}
\newcommand{\cS}{\mathcal{S}}
\newcommand{\cT}{\mathcal{T}}
\newcommand{\cU}{\mathcal{U}}
\newcommand{\cV}{\mathcal{V}}
\newcommand{\cW}{\mathcal{W}}
\newcommand{\cX}{\mathcal{X}}
\newcommand{\cY}{\mathcal{Y}}
\newcommand{\cZ}{\mathcal{Z}}
\newcommand{\one}{\mathbf{1}}
\newcommand{\dd}{\,d}
\newcommand{\eps}{\varepsilon}
\newcommand{\norm}[1]{\left\lVert #1\right\rVert}
\newcommand{\ip}[2]{\left\langle #1,#2\right\rangle}
\newcommand{\perpH}{^{\perp,\cH}}

% Preload AMS symbol fonts so microtype expansion is configured before first page use.
\AtBeginDocument{\begingroup\setbox0=\hbox{$\square\mathbb{R}$}\endgroup}

\begin{document}

\title[Volterra-Perron Pricing]{%
Volterra--Perron Pricing of the Option Manifold}
\author[J.~Vidal Llaurad\'o]{J.~Vidal Llaurad\'o\,\orcidlink{0009-0002-1918-5458}}
\thanks{DOI: \href{https://doi.org/10.2139/ssrn.6744119}{10.2139/ssrn.6744119}.}
\email{vidal.llaurado@gmail.com}
\subjclass[2020]{60G22, 91G20, 91G80, 47A75, 47N30}
\keywords{risk-neutral pricing, option manifolds, Volterra operators, Perron stress, latent-state recovery, residual kernels, synchronization risk, incomplete markets}
\date{May 2026}


\begin{abstract}
Public option surfaces can be calibrated without determining the latent risk-neutral \(\QQ\)-coordinates
that move non-public stress claims.  This paper treats the option manifold as the public image of a
latent Volterra--Perron stress state and formulates recovery as a quotient problem.  At date \(t\),
the canonical state is
\[
        \Xi_t^\QQ
        =
        (G_t^\QQ,\mathfrak{H}_t^\QQ,\rho_t^\QQ,e_t^\QQ,\ell_t^\QQ,\varpi_t^\QQ),
\]
with a directed Volterra contagion operator, a rough maturity-scaling state, a Perron spectral
triple, and a synchronized stress coordinate.  A residual-density selector is appended only after
recovery.  Public calibration does not fix this selected pricing coordinate.

The governing theorem identifies the public-observable quotient.  A Volterra--Perron source state is
locally recoverable, up to canonical scale, sign, basis, and Perron-normalization gauges, exactly
when its source quotient embeds into that public quotient.  Stable recovery is positivity of the
quotient observability modulus.  Without source-family restrictions, public-equivalent density
families preserve the full public option manifold while moving latent-stress claims.  With a
gauge-normalized Volterra source chart, a simple well-conditioned Perron eigenvalue,
nondegenerate carrier Gram matrix, maturity-response rank, and spanning index/single-name
synchronization responses, the canonical risk-neutral state is locally and stably recovered.

Recovery leaves pricing and hedging completion as separate conditions.  Residual price intervals
collapse only at public payoff-module completion; otherwise the public-annihilator fiber remains.  On recovered state
paths, discounted option price functionals satisfy martingale drift restrictions, phase transitions
add drift and covariance charges, and selected prices admit native sensitivities to Perron stress,
spectral radius, carrier direction, roughness, directed Volterra kernels, and the residual
selector.  The synchronization relative-value identity decomposes index option value relative to a
carrier-replicated single-name book into synchronized-stress premium, residual compression charge,
and tradeable basis error.  A finite Bachelier--Volterra--Perron benchmark gives a closed option
formula, a recovered-state pricing equation, explicit native Greeks, and a finite public hedge
residual.
\end{abstract}

\maketitle


\section{Introduction}\label{sec:intro}

Modern option pricing has become increasingly good at fitting the public surface.  Dynamic
surface models describe the surface itself \cite{contDaFonseca2002}; rough volatility explains the
fractional short-maturity signature in skews and realized volatility paths
\cite{bayer2016pricing,gatheralJaissonRosenbaum2018}; affine Volterra and rough Heston models give
tractable non-Markovian pricing equations \cite{elEuchRosenbaum2019,abijaber2019affine,abiJaberElEuch2019};
and neural or surrogate calibration maps interpolate option prices at scales unavailable to early
parametric models \cite{horvathMuguruzaTomas2021}.  These advances improve the public fit.  They
do not settle which latent risk-neutral coordinates the surface has identified.

The distinction is sharp on an index/single-name option manifold.  Single-name surfaces determine
marginal pricing laws.  They do not determine directed covariance routes, carrier absorption, or
synchronized index stress.  Index and basket options add joint information, but only through the
public response coordinates admitted by the carrier design.  In incomplete markets, exact
calibration of that public manifold still need not price a claim with nonzero loading on a latent
stress direction; unspanned risk appears inside the option manifold itself
\cite{collinDufresneGoldstein2002,follmerSondermann1986,schweizer1992meanVariance}.

The companion compression papers isolate the same obstruction from the pricing side.
Risk-neutral compression identifies finite \(\QQ\)-compressions of latent Perron stress from public
option data \cite{vidal2026compression}.  Residual pricing shows that public option prices can fix
\(C_{\cH}^{\QQ}(L_X^\QQ)\) while leaving the full latent loading \(L_X^\QQ\) movable along a
public-annihilator fiber \cite{vidal2026pricing}.  The present paper supplies the risk-neutral
option-manifold geometry that generates those compressions.  The option surface is a measurement
of a latent Volterra--Perron state, not the state itself.

At date \(t\), the public object is an option manifold \(\cM_t^{\mathrm{opt}}\) containing admitted
index, sector, carrier, and single-name option prices.  The latent source state is
\[
        \Xi_t^\QQ
        =
        (G_t^\QQ,\mathfrak{H}_t^\QQ,\rho_t^\QQ,e_t^\QQ,\ell_t^\QQ,\varpi_t^\QQ).
\]
Here \(G_t^\QQ\) is a directed Volterra contagion operator, \(\mathfrak{H}_t^\QQ\) records rough
maturity scaling, \((\rho_t^\QQ,e_t^\QQ,\ell_t^\QQ)\) is the Perron spectral triple of the screened
synchronization operator, and \(\varpi_t^\QQ\) is the synchronized latent stress coordinate.  A
selected pricing state \(\Theta_t^\QQ=(\Xi_t^\QQ,\theta_t)\) appends a residual density selector
inside the public-equivalent family.  Public prices are generated by the adapted response map
\[
        \mathfrak P_t:\Theta_t^\QQ\longmapsto \cM_t^{\mathrm{opt}} .
\]
The finite Volterra source chart gives this notation a specified representation: rough kernels
produce integrated covariance routes, basket and carrier prices observe finite linear functionals
of those routes, and residual selectors are bounded post-recovery pricing rules built from the
recovered canonical coordinates.

Calibration is therefore not a single inverse.  The public fit maps a state to observed quote
coordinates.  Source recovery asks whether the gauge quotient of the Volterra--Perron source state
embeds into the public-observable quotient.  Residual selection chooses a point in a
public-equivalent pricing fiber.  Tradeable completion projects terminal payoffs onto the admitted
instrument span.  These four maps may collapse in a complete finite model, but they are distinct
on the option manifold considered here.

The central statement is the Volterra--Perron option-manifold equivalence theorem.  Public option
prices identify exactly the public quotient \(\mathfrak X_t^{\mathrm{pub}}\).  A source family
\(\mathfrak S_t\) is locally recoverable only through the induced map
\[
        \mathfrak S_t/\!\sim_g\longrightarrow \mathfrak X_t^{\mathrm{pub}} .
\]
In a differentiable source chart the infinitesimal object is the quotient observability operator
\[
        \overline{\mathcal O}_t^\QQ
        =
        D(\mathcal R_t\circ\psi)(u_0)\quad\text{on the gauge quotient}.
\]
Recovery holds when this operator has no source kernel and a positive inverse modulus.  In the
block chart, the condition decomposes into rough maturity, directed-Volterra, carrier, and Perron
synchronization blocks.  Single-name vanilla surfaces identify marginal laws; they do not identify
these joint blocks by themselves.

Pointwise inverses are then assembled through time.  Uniform recovery windows give \(\cH_t\)-adapted
recovered processes
\[
        \widehat\Xi_t^\QQ=\mathcal I_t(\cM_t^{\mathrm{opt}})
\]
and pathwise stability under quote perturbations.  The same modulus controls regularized finite
panels and gives lower bounds: no reconstruction rule using the same quote data can beat the
inverse-modulus scale.  Ill-conditioning at carrier-thin, Perron-gap, maturity-rank, or
public-completion seams is an information limit, not a defect of a particular numerical inverse.

No-arbitrage enters twice.  Static public manifolds must satisfy the convex-order restrictions
implied by discounted martingales.  After recovery, option prices written as functionals of the
canonical state must have discounted local-martingale dynamics.  In a finite recovered chart this
gives
\[
        \partial_t\Phi_a
        +
        D\Phi_a[b]
        +
        \frac12\operatorname{Tr}(\Sigma^\ast D^2\Phi_a\Sigma)
        -
        r_t\Phi_a
        =
        0
\]
for each admitted option coordinate.  Smooth phase allocations add transition drift and covariance
terms because the phase weights themselves move with the recovered \(\QQ\)-state.

Residual non-identification remains after state recovery.  A dynamic public-equivalent density
family preserves every admitted public option payoff at date \(t\) and across maturities \(S>t\)
while moving a latent-stress claim whenever that claim has nonzero loading on the
public-annihilator fiber.  A selector \(\theta_t\), measurable
with respect to public and recovered coordinates, selects a price inside the residual interval.
This gives the dynamic form of the residual-pricing geometry of \cite{vidal2026pricing}, with no
second calibration of the public surface.

The recovered quotient also carries model-native sensitivities and a relative-value identity.
Native Greeks differentiate selected prices with respect to Perron stress, spectral radius, carrier
direction, roughness, directed Volterra kernels, and the residual selector.  The synchronization
identity decomposes the value of an index option relative to a carrier-replicated single-name book
into synchronized-stress premium, residual compression charge, and basis error.  The finite
Bachelier--Volterra--Perron benchmark closes these objects in a tractable chart with a closed option
formula, recovered-state pricing equation, explicit Greeks, and a finite hedge residual.  Its role is
the quotient-theory chart; universal smile modelling is outside the claim.

\subsection*{Organization}
Section~\ref{sec:manifold} defines the dynamic option manifold and information structure.
Section~\ref{sec:state} introduces the risk-neutral Volterra--Perron state and its gauge
equivalence.  Section~\ref{sec:functional-forms} gives a Volterra--Perron source chart.
Section~\ref{sec:phases} defines risk-neutral market phases and transition surfaces.
Section~\ref{sec:nonidentification} defines the public-observable quotient and proves the public
non-identification and public-completion results.
Section~\ref{sec:observability} develops risk-neutral observability, vanilla insufficiency, the
minimal option-manifold hierarchy, and finite-quote stability.  Section~\ref{sec:recovery} gives
the quotient recovery, finite-panel, inverse-modulus, and regularized reconstruction theorems.
Section~\ref{sec:dynamic-recovery} turns pointwise recovery into an adapted recovered state
process.  Section~\ref{sec:martingale}
derives the martingale-consistency drift restrictions for recovered risk-neutral state dynamics.
Section~\ref{sec:finite-benchmark} gives the canonical finite Bachelier--Volterra--Perron
benchmark with its formula, pricing equation, Greeks, and hedge residual.
Section~\ref{sec:gaussian} verifies the recovery conditions in a reduced Gaussian--Volterra
class.  Sections~\ref{sec:residual}--\ref{sec:rv} develop residual kernel selection, native Greeks,
public-neutral hedging, and the synchronization relative-value identity.
Section~\ref{sec:emm} gives one equivalent-martingale-measure specialization, and
Section~\ref{sec:examples} records reduced examples.

\section{Dynamic option manifolds}\label{sec:manifold}

Fix a filtered probability space
\[
        (\Omega,\cF,(\cF_t)_{t\ge0},\QQ)
\]
supporting all risk-neutral objects.  Conditional-expectation and martingale conventions follow
\cite{williams1991}.  The public market information is
\((\cH_t)_{t\ge0}\), the latent stress information is \((\cK_t)_{t\ge0}\), and
\[
        \cH_t\subseteq\cK_t\subseteq\cF_t.
\]
The tradeable public span \(\cS_t\subseteq L^2(\QQ,\cH_t)\) is the closed Hilbert span generated by
publicly tradeable instruments available at date \(t\).  It may be finite dimensional in
applications; the projection statements require only closedness.  This span is a public-feature
projection space, not the full dynamic self-financing gain space.  Section~\ref{sec:rv} separates
recoverability, tradeability, and hedgeability when terminal payoffs and dynamic strategies enter.

Let
\[
        B_t:=\exp\!\left(\int_0^t r_s\,\dd s\right)
\]
be the money-market account.  Unless an undiscounted payoff is explicitly displayed, all payoff
variables in Sections~\ref{sec:manifold}--\ref{sec:rv} are understood in date-\(t\) money-market
units:
\[
        \widetilde X_a:=\frac{B_t}{B_{T(a)}}X_a^{\mathrm{cash}}.
\]
Thus \(X_a\) below denotes \(\widetilde X_a\), and public prices are conditional expectations of
discounted payoffs.  Section~\ref{sec:emm} writes the discount factor explicitly again when the
underlying cash dynamics are introduced.

\begin{definition}[Public option manifold]\label{def:option-manifold}
Let \(\cI_t\) be the date-\(t\) list of underliers, strikes, maturities, and option types admitted
to the public surface.  A public option manifold is the \(\cH_t\)-measurable element
\[
        \cM_t^{\mathrm{opt}}
        =
        \bigl(C_t(a)\bigr)_{a\in\cI_t}
        \in \cY_t,
\]
where \(C_t(a)\) is the arbitrage-free risk-neutral price of payoff \(X_a\) and \(\cY_t\) is the
corresponding finite- or infinite-dimensional quote space.
\end{definition}

\begin{definition}[Quote-space norm]\label{def:quote-space-norm}
The quote space \(\cY_t\) is equipped with a fixed norm.  For a finite panel this may be a
weighted Euclidean norm, for example bid--ask or reliability weighted.  For a continuum maturity
and moneyness manifold it is a weighted \(C^k\), Sobolev, or Banach-supremum norm on the admitted
coordinate domain.  Dynamic stability statements use a deterministic ambient quote space, or a
Banach bundle with uniformly equivalent local norms, over the recovery window.
\end{definition}

\begin{definition}[Observable response coordinates]\label{def:observable-response}
Write a quote coordinate as
\[
        a=(u,\tau,m,\kappa,\chi),
\]
where \(u\) is an admitted index, sector, or single-name underlier, \(\tau>0\) is maturity,
\(m\) is the moneyness coordinate, \(\kappa\) records call, put, variance, or
surface-functional type, and \(\chi\in\{\mathrm B,\mathrm N\}\) records the pricing chart.
Here \(\mathrm B\) denotes the Black, or lognormal-forward, chart and \(\mathrm N\) denotes the
Bachelier, or normal-forward, chart.  In the Black chart one may take
\(m^{\mathrm B}=\log(K/F_{u,t,T})\); in the normal chart one may take
\(m^{\mathrm N}=K-F_{u,t,T}\), or its deterministically rescaled version.  For variance and
surface-functional coordinates, \(\chi\) records the convention in which the reported total
variance is expressed.  For a latent state \(\Theta\), the observable response map is
\[
        \mathcal R_t(\Theta)(u,\tau,m,\kappa,\chi)
        :=
        C_t^\Theta(u,\tau,m,\kappa,\chi).
\]
The option manifold \(\cM_t^{\mathrm{opt}}\) is the realized public value of this response map on
the admitted coordinate set.  Continuous maturity/moneyness manifolds and finite quoted panels
are treated as two observation regimes for the same map.
\end{definition}

\begin{assumption}[Arbitrage-free public manifold]\label{ass:arbitrage-free-manifold}
For each \(t\), the model works under a selected equivalent martingale measure \(\QQ\) on the
public payoff span, so the prices \(C_t(a)\) are represented by \(\QQ\)-conditional
expectations.  Equivalently, the public manifold is a no-arbitrage price system on its admitted
public payoff family, and \(\QQ\) is a selected representative of the equivalent pricing class
\cite{harrisonKreps1979,harrisonPliska1981,delbaen1994,karatzasShreve1998}.
Residual selectors introduced below are conditional density changes around this representative
that preserve the admitted public prices.
\end{assumption}

\begin{definition}[Public no-arbitrage cone]\label{def:no-arbitrage-cone}
For a finite public payoff panel, write \(X=(X_1,\dots,X_M)\) and let
\[
        \cA_t
        :=
        \left\{p\in\R^M:
        h\cdot X\ge0\ \QQ\text{-a.s. and }h\cdot X>0
        \text{ with positive probability imply }h\cdot p>0
        \right\}
\]
be the strict no-arbitrage cone of public prices.  Its tangent cone at a price vector
\(p_0\in\cA_t\) is denoted \(T_{p_0}\cA_t\).
\end{definition}

\begin{proposition}[No-arbitrage tangent condition]\label{prop:no-arbitrage-tangent}
Let \(F:U\to\R^M\) be a finite-panel option response map whose image lies in the
no-arbitrage cone \(\cA_t\), and let \(p_0=F(u_0)\).  Then every source tangent direction
that is visible in public quotes satisfies
\[
        DF(u_0)h\in T_{p_0}\cA_t .
\]
If \(p_0\) is in the relative interior of \(\cA_t\), this imposes no first-order inequality.
At the boundary of the cone, it removes source directions whose first-order public response
would create a static arbitrage at first order.
\end{proposition}

\begin{proof}
For every sufficiently small \(s>0\) with \(u_0+sh\in U\), the response \(F(u_0+sh)\) lies in
\(\cA_t\).  Therefore
\[
        \frac{F(u_0+sh)-p_0}{s}
        \longrightarrow DF(u_0)h
\]
is a tangent direction to \(\cA_t\) at \(p_0\).  If \(p_0\) lies in the relative interior of the
cone, the tangent cone is the linear span of \(\cA_t-p_0\), so no additional first-order inequality
is imposed.  At a boundary point, the tangent cone is a proper cone and excludes first-order
responses that point outside the public no-arbitrage region.
\end{proof}

\begin{definition}[Conditional convex order]\label{def:conditional-convex-order}
Let \(\mu_{t,T_1}\) and \(\mu_{t,T_2}\) be conditional laws, given \(\cH_t\), of a discounted
underlying vector at maturities \(T_1<T_2\).  Write
\[
        \mu_{t,T_1}\preceq_{\mathrm{cx}}\mu_{t,T_2}
\]
if, for every convex function \(\varphi\) for which both conditional expectations are finite,
\[
        \int \varphi(x)\,\mu_{t,T_1}(\dd x)
        \le
        \int \varphi(x)\,\mu_{t,T_2}(\dd x)
\]
pathwise, conditionally on \(\cH_t\).
\end{definition}

\begin{theorem}[Calendar convex-order restriction]\label{thm:calendar-convex-order}
If the discounted underlying vector is a \(\QQ\)-martingale and \(\mu_{t,T}\) denotes its
conditional law given \(\cH_t\), then, for \(t<T_1<T_2\),
\[
        \mu_{t,T_1}\preceq_{\mathrm{cx}}\mu_{t,T_2}.
\]
For a scalar underlying this implies the usual calendar inequalities for call prices,
\[
        C_t(K,T_2)\ge C_t(K,T_1),
\]
after the common-forward convention is imposed.  For a vector underlying, every linear projection
must satisfy the corresponding one-dimensional convex-order restriction.
\end{theorem}

\begin{proof}
For convex \(\varphi\), Jensen's inequality and the martingale property give
\[
        \EE^\QQ[\varphi(\widetilde S_{T_2})\mid\cF_{T_1}]
        \ge
        \varphi(\EE^\QQ[\widetilde S_{T_2}\mid\cF_{T_1}])
        =
        \varphi(\widetilde S_{T_1}).
\]
Taking conditional expectation given \(\cH_t\) gives the convex-order statement.  The scalar call
form follows from the representation of convex functions by affine functions and positive
mixtures of calls.  Linear projections of vector martingales are martingales, so the same
argument applies to every projection.
\end{proof}

\begin{corollary}[Dynamic falsifiability before source fitting]\label{cor:convex-order-falsifiability}
A multi-maturity public option manifold that violates conditional convex order cannot be
generated by any equivalent martingale measure, regardless of the Volterra--Perron source
family.  Conversely, a source-family fit that matches maturity slices separately but violates
the convex-order monotonicity across maturities is dynamically inadmissible.
\end{corollary}

\begin{proof}
The conditional convex-order inequalities of Theorem~\ref{thm:calendar-convex-order} are necessary
for every equivalent martingale measure generating the discounted underlier.  They are imposed on
the public manifold before any Volterra--Perron source family is chosen.  A violation therefore
rules out the manifold at the martingale level.  A fit that matches separate maturity slices but
breaks those inequalities produces the same obstruction, now inside the fitted source family.
\end{proof}


\begin{definition}[Latent loading and public compression]\label{def:latent-loading}
For \(X\in L^2(\QQ,\cF_T)\), define the latent loading at date \(t\) by
\[
        L_{t,T}^\QQ(X):=\EE^\QQ[X\mid\cK_t],
\]
and its public compression by
\[
        C_{\cH_t}^\QQ L_{t,T}^\QQ(X)
        :=
        \EE^\QQ[L_{t,T}^\QQ(X)\mid\cH_t].
\]
The tradeable public projection is \(C_{\cS_t}^\QQ X\), the orthogonal projection of \(X\) onto
\(\cS_t\) in \(L^2(\QQ)\).
\end{definition}

\begin{proposition}[Compression factorization of public prices]\label{prop:compression-factor}
If an option payoff \(X_a\) is priced at date \(t\) by conditional expectation under \(\QQ\), then
its public price depends on the latent loading only through its public compression:
\[
        C_t(a)
        =
        \EE^\QQ[X_a\mid\cH_t]
        =
        C_{\cH_t}^\QQ L_{t,T(a)}^\QQ(X_a).
\]
\end{proposition}

\begin{proof}
The latent loading is \(L_{t,T(a)}^\QQ(X_a)=\EE^\QQ[X_a\mid\cK_t]\).  Since
\(\cH_t\subseteq\cK_t\), the tower property gives
\[
        \EE^\QQ[X_a\mid\cH_t]
        =
        \EE^\QQ[\EE^\QQ[X_a\mid\cK_t]\mid\cH_t]
        =
        C_{\cH_t}^\QQ L_{t,T(a)}^\QQ(X_a).
\]
Thus the public price is the public compression of the latent loading.
\end{proof}

The full-basket and multi-date results identify the information frontier of ideal complete claim
families.  They do not assert that listed-option datasets contain such manifolds.  Empirical
panels usually provide finite, noisy submanifolds.  Later sections therefore keep finite-panel
stability, source-family restrictions, and residual selection separate from full state recovery.

\begin{definition}[Full terminal basket manifold]\label{def:full-terminal-basket-manifold}
Fix \(t<T\), and let \(S_T=(S_T^1,\dots,S_T^N)\in\R_+^N\) be the terminal underlier vector.
The full nonnegative basket-call manifold at \((t,T)\) is
\[
        C_t(w,K;T)
        :=
        \EE^\QQ\!\left[
        \frac{B_t}{B_T}(w^\top S_T-K)^+
        \,\middle|\,\cH_t
        \right],
        \qquad w\in\R_+^N,\ K\ge0.
\]
If signed spread baskets are admitted, \(w\) may range over \(\R^N\) and \(K\) over \(\R\).
\end{definition}

\begin{theorem}[Static basket-law recovery]\label{thm:static-basket-law-recovery}
Assume the conditional pricing law of \(S_T\) is supported on \(\R_+^N\), and assume the public
discount factor \(P(t,T)\) is known at date \(t\), equivalently the zero-coupon payoff at
maturity \(T\) is included in the admitted public span.  Write
\[
        D_{t,T}:=\frac{B_t}{B_T},
        \qquad
        P(t,T):=\EE^\QQ[D_{t,T}\mid\cH_t].
\]
The full nonnegative basket-call manifold, augmented by this public discount coordinate,
determines the \(T\)-forward conditional pricing law
\[
        \mu^T_{t,T}(\omega,A)
        :=
        \frac{
        \EE^\QQ[D_{t,T}\mathbf 1_{\{S_T\in A\}}\mid\cH_t](\omega)
        }{P(t,T)(\omega)}
\]
up to \(\QQ\)-null sets.  Consequently it determines the date-\(t\) price of every bounded
Borel European payoff \(f(S_T)\) by integration against \(\mu^T_{t,T}\).  If \(D_{t,T}\) is
\(\cH_t\)-measurable, for example under deterministic discounting conditional on \(\cH_t\), then
\(\mu^T_{t,T}\) coincides with the raw conditional law
\(\Law_\QQ(S_T\mid\cH_t)\).
\end{theorem}

\begin{proof}
The public discount coordinate supplies \(P(t,T)\).  After division by \(P(t,T)\), each fixed
\(w\) gives the full call surface of the \(T\)-forward pricing distribution of \(w^\top S_T\).
The Breeden--Litzenberger distributional derivative recovers the law of each nonnegative scalar
projection under \(\mu^T_{t,T}\) \cite{breedenLitzenberger1978}.  Consequently the manifold
determines, pathwise in the conditioning variable, the \(T\)-forward conditional Laplace transform
\[
        L^T_{t,T}(\xi)
        =
        \int_{\R_+^N}e^{-\xi^\top x}\,\mu^T_{t,T}(\dd x),
        \qquad \xi\in\R_+^N,
\]
since \(\xi=sw\) for some \(s\ge0\) and \(w\in\R_+^N\).  The multivariate Laplace transform on
the positive orthant uniquely determines a probability law on \(\R_+^N\).  Prices of bounded
European payoffs then follow from
\[
        \EE^\QQ[D_{t,T}f(S_T)\mid\cH_t]
        =
        P(t,T)\int f(x)\,\mu^T_{t,T}(\dd x).
\]
\end{proof}

\begin{corollary}[Terminal-law maximality]\label{cor:terminal-law-maximality}
Single-name vanilla surfaces recover only the one-dimensional \(T\)-forward pricing marginals of
\(S_T\).  A finite family of basket surfaces recovers only the pricing laws of finitely many
projections \(w_a^\top S_T\).  The full terminal basket manifold is maximal for terminal
European payoffs at date \(T\).  Path-law information lies outside that terminal object.  Knowing
the terminal pricing law for each \(T\) does not by itself identify lead-lag structure, directed Volterra kernels,
quadratic covariation paths, or the joint pricing law of
\((S_{T_1},\dots,S_{T_m})\) across future dates.
\end{corollary}

\begin{proof}
Theorem~\ref{thm:static-basket-law-recovery} identifies the full terminal pricing law from the full
basket manifold.  Restricting the manifold to single-name vanillas leaves only coordinate
marginals; restricting it to finitely many baskets leaves only finitely many projection laws.  None
of these terminal objects specifies a coupling across future dates.  Directed kernels,
lead-lag effects, quadratic-covariation paths, and multi-date pricing laws are functions of that
missing coupling, and therefore require the multi-date manifold of Theorem~\ref{thm:path-law-recovery}
or additional source restrictions.
\end{proof}


\begin{definition}[Multi-date linear option manifold]\label{def:multi-date-linear-manifold}
Fix \(t<s_1<\cdots<s_m\le T_\ast\).  Let \(X_s\) be a vector of tradeable state coordinates,
such as log-prices, forward returns, or centered variance factors.  The multi-date linear option
manifold is
\[
        C_t^{(s_1,\dots,s_m)}(\alpha,K)
        :=
        \EE^\QQ\!\left[
        \frac{B_t}{B_{s_m}}
        \left(\sum_{k=1}^m\alpha_k^\top X_{s_k}-K\right)^+
        \,\middle|\,\cH_t
        \right],
\]
for all admitted \(\alpha=(\alpha_1,\dots,\alpha_m)\in\R^{mN}\) and strikes \(K\in\R\).
\end{definition}

\begin{theorem}[Finite-dimensional path-law recovery]\label{thm:path-law-recovery}
For a fixed grid \(s_1<\cdots<s_m\), assume the public discount factor \(P(t,s_m)\) is known, or
that the zero-coupon payoff at \(s_m\) is included in the admitted public span.  The full
multi-date linear option manifold determines the \(s_m\)-forward conditional pricing law of
\[
        (X_{s_1},\dots,X_{s_m}).
\]
If these manifolds are observed for every finite grid in \([t,T_\ast]\), and the recovered
conditional pricing laws are projectively consistent and tight on the chosen path space, they
determine the conditional path pricing law of \((X_s)_{s\in[t,T_\ast]}\).  If all involved
discount factors are \(\cH_t\)-measurable, the same statements hold for the raw
\(\QQ\)-conditional laws.
\end{theorem}

\begin{proof}
After normalization by the public discount factor \(P(t,s_m)\), each \(\alpha\)-slice gives a
full call surface in \(K\) and therefore the \(s_m\)-forward pricing law of the scalar projection
\(\sum_k\alpha_k^\top X_{s_k}\).  The Cramer--Wold theorem recovers the joint pricing law on
\(\R^{mN}\).  Projective consistency across grids and tightness then give the conditional path
pricing law by Kolmogorov extension, with the usual path-space regularity condition when
continuous or c\`adl\`ag paths are imposed.  Public discounting reduces this to raw
conditional-law recovery.
\end{proof}

\begin{corollary}[Gaussian--Volterra covariance recovery]\label{cor:gv-path-covariance-recovery}
Suppose \(X\) is conditionally Gaussian given \(\cH_t\) and admits the Volterra representation
\[
        X_s=\int_t^s K(s,r)\dd W_r,\qquad t\le r\le s\le T_\ast,
\]
with a square-integrable causal kernel \(K(s,r)\).  The multi-date manifold recovers the
conditional covariance kernel
\[
        \Gamma(s,u)
        =
        \Cov^\QQ(X_s,X_u\mid\cH_t)
        =
        \int_t^{s\wedge u}K(s,r)K(u,r)^\top\dd r .
\]
The kernel \(K\) is identified up to orthogonal driver gauge.  If a causal lower-triangular
Volterra gauge with positive diagonal is imposed, the kernel is uniquely identified by the
operator Cholesky factorization of \(\Gamma\).
\end{corollary}

\begin{proof}
Theorem~\ref{thm:path-law-recovery} recovers the conditional finite-dimensional laws of the process
on every grid.  For a conditionally Gaussian process, these laws are determined by their conditional
means and covariance matrices; in the centered Volterra representation the mean is zero and the
covariances are the displayed values of \(\Gamma\).  Two causal square-integrable kernels with the
same covariance operator differ by an orthogonal transformation of the driving noise on the driver
space.  Imposing the causal lower-triangular gauge with positive diagonal selects the operator
Cholesky factor and removes this driver gauge.
\end{proof}


\begin{theorem}[Numeraire invariance of local recovery]\label{thm:numeraire-invariance}
Let \(N\) be a strictly positive traded numeraire and let \(\QQ^N\) be the associated equivalent
measure.  Suppose the money-market response map \(R_t\) and the \(N\)-numeraire response map
\(R_t^N\) are related by a bounded invertible \(\cH_t\)-measurable linear transformation
\[
        R_t^N=T_tR_t
\]
of the quote space.  Then
\[
        D(R_t^N\circ\psi)(u_0)=T_tD(R_t\circ\psi)(u_0).
\]
Local injectivity, quotient rank, finite-panel information rank, and overidentifying normal-space
restrictions are invariant under the numeraire change.  Inverse constants change only by the
operator norms of \(T_t\) and \(T_t^{-1}\).
\end{theorem}

\begin{proof}
The change-of-numeraire identity gives equality of arbitrage-free prices after the public linear
quote transformation \(T_t\).  Since \(T_t\) is invertible, it neither creates nor removes null
directions of the response derivative on the quotient tangent space.  Left inverses transform by
composition with \(T_t^{-1}\), giving the norm statement.
\end{proof}

\section{The risk-neutral Volterra--Perron state}\label{sec:state}

The state is introduced at the quotient level.  Gaussian--Volterra assumptions enter later as one
source family verifying the required response ranks.

\begin{definition}[Risk-neutral Volterra--Perron state]\label{def:q-state}
The risk-neutral Volterra--Perron state at date \(t\) is
\[
        \Xi_t^\QQ
        =
        (G_t^\QQ,\mathfrak{H}_t^\QQ,\rho_t^\QQ,e_t^\QQ,\ell_t^\QQ,\varpi_t^\QQ),
\]
where:
\begin{enumerate}[leftmargin=2em]
\item \(G_t^\QQ\) is a directed Volterra operator on a Hilbert space \(\cU_t\);
\item \(\mathfrak{H}_t^\QQ\) is a rough maturity-scaling coordinate;
\item \((\rho_t^\QQ,e_t^\QQ,\ell_t^\QQ)\) is a right-left Perron spectral triple for a screened
positive synchronization operator \(A_t^\QQ\), normalized by
\(\ell_t^\QQ(e_t^\QQ)=1\);
\item \(\varpi_t^\QQ\in L^2(\QQ,\cK_t)\) is a centered synchronized latent stress coordinate;
\end{enumerate}
A selected pricing state is
\[
        \Theta_t^\QQ=(\Xi_t^\QQ,\theta_t),
\]
where \(\theta_t\) is a residual-density selector used to price latent residual claims after the
state has been recovered or imposed.
\end{definition}

\begin{definition}[Perron spectral projector]\label{def:perron-projector}
Let \(A\) be a screened positive synchronization operator with simple Perron eigenvalue
\(\rho_A\), right Perron vector \(e_A\), and left Perron functional \(\ell_A\), normalized by
\(\ell_A(e_A)=1\).  The Perron spectral projector is
\[
        \Pi_A:=e_A\otimes\ell_A,
        \qquad
        \Pi_A x=e_A\,\ell_A(x).
\]
The canonical Perron coordinate is \((\rho_A,\Pi_A)\).  The triple
\((\rho_A,e_A,\ell_A)\) is retained for computations, but the projector is the
normalization-free coordinate of the Perron mode.
\end{definition}

\begin{proposition}[Gauge invariance of the Perron projector]\label{prop:perron-projector-gauge}
The projector \(\Pi_A=e_A\otimes\ell_A\) is invariant under reciprocal normalization changes
\[
        e_A\mapsto c e_A,\qquad \ell_A\mapsto c^{-1}\ell_A,\qquad c>0.
\]
If \(A\) is conjugated by an admissible hidden-basis change \(S\), so that
\(\widetilde A=SAS^{-1}\), then
\[
        \Pi_{\widetilde A}=S\Pi_A S^{-1}.
\]
Thus \((\rho_A,\Pi_A)\) is the natural quotient coordinate of the Perron synchronization mode.
\end{proposition}

\begin{proof}
Reciprocal scaling gives \((ce_A)\otimes(c^{-1}\ell_A)=e_A\otimes\ell_A\).  Under conjugacy,
\(Se_A\) and \(\ell_A S^{-1}\) satisfy
\[
        SAS^{-1}(Se_A)=\rho_A Se_A,
        \qquad
        (\ell_A S^{-1})SAS^{-1}=\rho_A(\ell_A S^{-1}),
\]
and \((\ell_A S^{-1})(Se_A)=1\).  The rank-one projector is therefore transported by
\(S\Pi_A S^{-1}\).
\end{proof}

\begin{definition}[Option pricing map]\label{def:pricing-map}
Let \(\mathfrak X_t\) be an admissible state space for \(\Theta_t^\QQ\).  The option pricing map is
an adapted map
\[
        \mathfrak P_t:\mathfrak X_t\to\cY_t,
        \qquad
        \Theta_t^\QQ\mapsto \cM_t^{\mathrm{opt}}.
\]
The map is public because \(\mathfrak P_t(\Theta_t^\QQ)\) is \(\cH_t\)-measurable even when
\(\Theta_t^\QQ\) contains \(\cK_t\)-measurable latent components.
\end{definition}

\begin{definition}[Conditional public response]\label{def:conditional-public-response}
When a source state contains latent \(\cK_t\)-measurable coordinates, the public response map
is a conditional response.  If
\(\widetilde{\mathfrak P}_t(\Theta)\in L^1(\QQ,\cK_t;\cY_t)\) denotes the latent quote response
generated before public compression, the public option manifold is
\[
        \mathfrak P_t(\Theta)
        :=
        \EE^\QQ\!\left[
        \widetilde{\mathfrak P}_t(\Theta)
        \,\middle|\,\cH_t
        \right].
\]
If the recovered state coordinates are already represented by \(\cH_t\)-measurable canonical
parameters, this convention reduces to the ordinary deterministic response chart.  All
public-indistinguishability and recovery statements below are statements about this
\(\cH_t\)-measurable response, not about direct observation of hidden sample-path coordinates.
\end{definition}

\begin{remark}
The recovery theorems identify the canonical class of \(\Xi_t^\QQ\).  When the selector
\(\theta_t\) enters only through public-annihilating residual directions, admitted public option
quotes leave it unidentified.  A post-recovery rule must fix it, or an augmented manifold must add
residual-price instruments.
\end{remark}

\begin{definition}[Calibration ledger]\label{def:calibration-ledger}
At date \(t\), a finite option-manifold calibration separates four maps:
\[
        \text{public fit:}\quad
        y_t\approx \mathfrak P_t(\Theta),
\]
\[
        \text{source recovery:}\quad
        y_t\mapsto [\Xi_t^\QQ]\in\overline{\mathfrak X}_t,
\]
\[
        \text{residual selection:}\quad
        [\Xi_t^\QQ]\mapsto \theta_t,
\]
and
\[
        \text{tradeable completion:}\quad
        X\mapsto \Pi_{\cS_t}X .
\]
The first is quote calibration in the public option panel.  The second is source recovery on the
canonical quotient.  The third chooses a public-equivalent residual price inside the remaining
envelope.  The fourth is projection onto the finite public instrument span.
\end{definition}

\begin{remark}[Calibration, recovery, selection, and hedging]\label{rem:calibration-not-single-inverse}
In a complete Black--Scholes-style finite model these four maps may collapse into one fitted
parameter vector and one hedge.  In the Volterra--Perron option manifold they are logically
separate.  A surface can be well calibrated without source recovery; a source can be recovered
without identifying the residual selector; and a selected residual price can remain outside the
tradeable span.  The recovery theorems below concern the second map.
\end{remark}

\begin{definition}[Gauge equivalence]\label{def:gauge}
Two states \(\Theta,\widetilde\Theta\in\mathfrak X_t\) are gauge equivalent, written
\(\Theta\sim_g\widetilde\Theta\), if they differ only by transformations that leave all latent
loadings and all option-manifold prices unchanged, namely sign conventions for one-dimensional
latent coordinates, bounded invertible basis changes in hidden factor coordinates, common rescalings
absorbed by the Perron normalization \(\ell(e)=1\), and relabelings of unobserved driver bases.
The canonical state space is the quotient
\[
        \overline{\mathfrak X}_t:=\mathfrak X_t/\!\sim_g.
\]
\end{definition}

\begin{remark}
The quotient is essential.  Perron eigenvectors, Volterra driver bases, and one-dimensional
stress coordinates have unavoidable normalization choices.  Recovery of the latent state can only
mean recovery of the equivalence class \([\Xi_t^\QQ]\in\overline{\mathfrak X}_t\).
\end{remark}

\begin{remark}
The coordinate \(\varpi_t^\QQ\) is latent in parametrization but recoverable only through its loading,
conditional moment, or canonical horizon lift encoded by the option manifold.
The recovery statements do not assert observation of an arbitrary future sample-path realization.
When a terminal stress payoff is needed, \(\varpi_{t,T}^\QQ\) denotes the horizon-lifted synchronized
stress direction generated by the recovered source chart.
\end{remark}

\begin{assumption}[Local gauge slice]\label{ass:local-gauge-slice}
Near a reference state \(\Xi_0\), the gauge transformations form a \(C^1\) local group
\(\mathcal G_t\) acting on the source chart, and the public response map is constant on
\(\mathcal G_t\)-orbits.  The orbit tangent
\[
        E_g:=T_{\Xi_0}(\mathcal G_t\cdot\Xi_0)
\]
is closed and admits a closed complement \(E_c\), so \(E=E_g\oplus E_c\).  A gauge-normalized
chart is a local slice \(U_c\subset E_c\) meeting each nearby orbit once.  The quotient tangent
space is identified with \(E_c\).
\end{assumption}

\begin{remark}[Local gauge fixing]\label{rem:local-not-global-gauge}
Assumption~\ref{ass:local-gauge-slice} is intentionally local.  It gives a cross-section near
the recovered state used in the inverse-function and stability arguments.  It does not assert the
existence of a global, singularity-free section of the full orbit space
\(\mathfrak X_t/\mathcal G_t\).  Global stabilizers, orientation changes, eigenvector
normalization seams, or nontrivial orbit topology may force an atlas of slices rather than a
single canonical representative.  All recovery statements below are therefore local or patched
through the recovery atlas of Section~\ref{sec:dynamic-recovery}; the quotient state is global
only up to the declared gauge equivalence.
\end{remark}

\begin{definition}[Gauge-normalized source chart]\label{def:gauge-normalized-chart}
A source chart is gauge-normalized near \(\Xi_0\) if each state is represented by canonical
coordinates
\[
        u
        =
        (u_{\mathfrak H},u_G,u_\rho,u_e,u_P)
        \in
        E_{\mathfrak H}\oplus E_G\oplus E_\rho\oplus E_e\oplus E_P,
\]
corresponding to
\[
        (\mathfrak{H}^\QQ,G^\QQ,\rho^\QQ,e^\QQ,\ell^\QQ,P^\QQ),
\]
with the following conventions:
\begin{enumerate}[leftmargin=2em]
\item \(\ell^\QQ(e^\QQ)=1\), \(e^\QQ\) lies in a fixed positive cone, and the remaining Perron
normalization is absorbed into \(P^\QQ\);
\item the latent stress coordinate \(P^\QQ\) is centered and has a fixed sign convention against
a nonzero reference functional;
\item hidden Volterra driver bases are represented in a fixed orthonormal or triangular gauge.
\end{enumerate}
The block \(u_e\) should be read as the normalized Perron-projector block: the left functional
\(\ell^\QQ\) is induced by the screened operator and the normalization
\(\ell^\QQ(e^\QQ)=1\) inside the chosen slice.  A redundant chart may include an explicit
\(u_\ell\) coordinate, but the quotient tangent then removes the normalization direction and
returns the same canonical block.
The corresponding quotient tangent space is the direct sum of these gauge-normalized coordinate
blocks.
\end{definition}

\section{A Volterra--Perron source chart}\label{sec:functional-forms}

The recovery theorems are stated for general source charts.  This section fixes a finite
Volterra--Perron chart in which the response ranks can be read block by block.  Rough maturity
scaling carries Volterra source geometry; basket and carrier responses carry covariance routes;
the screened route matrix carries the Perron synchronization coordinate.  Residual selectors are
added only afterward as public-equivalent pricing extensions, unless public completion or residual
instruments make them observable.

\begin{definition}[Volterra route kernels]\label{def:volterra-route-kernels}
Fix \(N\) underliers and a maximum maturity \(\bar\tau>0\).  Conditional on \(\cH_t\), let
\[
        B^\QQ_{t+s}=(B^\QQ_{1,t+s},\dots,B^\QQ_{N,t+s})
\]
be a risk-neutral Gaussian driver vector with conditionally independent standard components.  A
nontrivial instantaneous driver covariance can be included as an additional gauge-normalized
matrix coordinate; the displayed route formula fixes the independent-driver gauge.  The centered
latent forward-variance stress of name \(i\) over future time \(s\in[0,\bar\tau]\) is represented by
\[
        Y_{i,t}^\QQ(s)
        =
        \sum_{j=1}^N
        \int_0^s K_{ij,t}^\QQ(s-r)\dd B^\QQ_{j,t+r}.
\]
The route from source \(j\) to target \(i\) has kernel
\[
        K_{ij,t}^\QQ(x)
        =
        a_{ij,t}\,
        x^{H_{ij,t}^\QQ-\frac12}
        \exp(-\beta_{ij,t}x)
        \left(1+\sum_{k=1}^{K_0} c_{ij,k,t}x^k\right),
        \qquad 0<x\le\bar\tau,
\]
where \(H_{ij,t}^\QQ\in(0,1)\), \(\beta_{ij,t}\ge0\), and the polynomial modulation is bounded
away from zero on the recovery neighborhood.  The directed Volterra operator \(G_t^\QQ\) is the
matrix of off-diagonal route coordinates
\[
        G_{ij,t}^\QQ=(a_{ij,t},H_{ij,t}^\QQ,\beta_{ij,t},c_{ij,1,t},\dots,c_{ij,K_0,t}),
        \qquad i\ne j.
\]
\end{definition}

The observable covariance routes generated by these kernels are explicit.  For
\(\tau\le\bar\tau\), define
\[
        V_{ij,t}^\QQ(\tau)
        :=
        \Cov^\QQ\!\left(
        \int_0^\tau Y_{i,t}^\QQ(s)\dd s,
        \int_0^\tau Y_{j,t}^\QQ(s)\dd s
        \,\middle|\,\cH_t
        \right).
\]
Equivalently,
\[
        V_{ij,t}^\QQ(\tau)
        =
        \sum_{q=1}^N
        \int_0^\tau
        \left(\int_r^\tau K_{iq,t}^\QQ(s-r)\dd s\right)
        \left(\int_r^\tau K_{jq,t}^\QQ(s-r)\dd s\right)\dd r .
\]
\begin{lemma}[Analytic leading power of a route covariance]\label{lem:route-leading-power}
Suppose the \(q\)-source routes entering targets \(i\) and \(j\) satisfy, as \(x\downarrow0\),
\[
        K_{iq,t}^\QQ(x)
        =
        a_{iq,t}x^{H_{iq,t}^\QQ-\frac12}\bigl(1+O(x)\bigr),
        \qquad
        K_{jq,t}^\QQ(x)
        =
        a_{jq,t}x^{H_{jq,t}^\QQ-\frac12}\bigl(1+O(x)\bigr),
\]
with nonzero amplitudes.  Then the corresponding route contribution to
\(V_{ij,t}^\QQ(\tau)\) has the short-maturity expansion
\[
        V_{ij,t}^{(q),\QQ}(\tau)
        =
        A_{ijq,t}\,
        \tau^{H_{iq,t}^\QQ+H_{jq,t}^\QQ+2}
        \bigl(1+O(\tau)\bigr),
\]
where
\[
        A_{ijq,t}
        =
        \frac{a_{iq,t}a_{jq,t}}
        {(H_{iq,t}^\QQ+\frac12)(H_{jq,t}^\QQ+\frac12)
        (H_{iq,t}^\QQ+H_{jq,t}^\QQ+2)}
\]
up to the displayed normalization of the leading kernel amplitudes.
\end{lemma}

\begin{proof}
For \(0\le r\le\tau\), integrate the leading kernel expansion in the variable \(s-r\):
\[
        \int_r^\tau K_{iq,t}^\QQ(s-r)\dd s
        =
        \frac{a_{iq,t}}{H_{iq,t}^\QQ+\frac12}
        (\tau-r)^{H_{iq,t}^\QQ+\frac12}
        \bigl(1+O(\tau-r)\bigr).
\]
The same expression holds for the \(j\)-route.  Their product has leading power
\((\tau-r)^{H_{iq,t}^\QQ+H_{jq,t}^\QQ+1}\).  Hence
\[
        \int_0^\tau
        (\tau-r)^{H_{iq,t}^\QQ+H_{jq,t}^\QQ+1}\dd r
        =
        \frac{\tau^{H_{iq,t}^\QQ+H_{jq,t}^\QQ+2}}
        {H_{iq,t}^\QQ+H_{jq,t}^\QQ+2},
\]
and the displayed coefficient follows after multiplying by the two one-step integration constants.
\end{proof}

Roughness and directed-route coefficients therefore enter the public manifold through analytic
maturity powers and logarithmic derivatives.  The finite dictionary gives a sufficient response
family for exact finite-dimensional recovery.  The lemma supplies the local invariant used when the
full covariance route is not assumed to lie exactly in a finite power-exponential span.

\begin{definition}[Rough power-exponential dictionary]\label{def:rough-power-exponential-dictionary}
On a maturity interval \((0,\bar\tau)\), a finite Volterra response dictionary consists of
functions of the form
\[
        \tau^\alpha e^{-\beta\tau}p(\tau),
        \qquad \alpha>-1,\quad \beta\ge0,
\]
where \(p\) is a polynomial of degree at most \(K\).  Dictionary atoms are grouped by their
pairs \((\alpha,\beta)\).
A dictionary used for identification is taken in reduced form.  After expanding each polynomial,
the scalar atoms
\[
        \{\tau^{\alpha_\ell+k}e^{-\beta_\ell\tau}:0\le k\le\deg p_\ell\}
\]
are pairwise distinct and linearly independent on every open subinterval.  A convenient
sufficient check is that, within each fixed \(\beta\), no two exponents created by the admitted
polynomial shifts coincide.  This removes notational duplicates such as
\(\tau^\alpha\tau\) and \(\tau^{\alpha+1}\).
\end{definition}

\begin{theorem}[Power-exponential route identifiability]\label{thm:power-exponential-identifiability}
Let
\[
        f(\tau)
        =
        \sum_{\ell=1}^m
        \tau^{\alpha_\ell}e^{-\beta_\ell\tau}p_\ell(\tau),
        \qquad 0<\tau<\bar\tau,
\]
where the pairs \((\alpha_\ell,\beta_\ell)\) are distinct, each \(p_\ell\) is a polynomial of
known maximum degree, and the expanded scalar atoms form a reduced dictionary in the sense of
Definition~\ref{def:rough-power-exponential-dictionary}.  If \(f\) vanishes on a nonempty open
interval, then every expanded scalar coefficient is zero; equivalently every reduced dictionary
coefficient vanishes.  Hence the coefficient vector of a finite Volterra response dictionary is
identified by the continuum maturity response.  There is also a nonzero
analytic finite-panel determinant whose nonvanishing identifies the same coefficient vector, and
the deficient panels form a zero set with empty interior.
\end{theorem}

\begin{proof}
Expand each polynomial and collect the response in scalar atoms
\(\tau^{\gamma_r}e^{-\beta_r\tau}\).  These atoms are real analytic on \((0,\infty)\).  If the
response vanishes on a nonempty open subinterval, the identity theorem gives vanishing on the whole
positive half-line.  Among the remaining nonzero coefficients, choose the smallest exponential
rate \(\beta_\ast\).  Multiplying by \(e^{\beta_\ast\tau}\) and letting \(\tau\to\infty\) removes all
larger-rate groups exponentially.  Inside the \(\beta_\ast\)-group, distinct powers are separated by
dividing successively by the largest remaining power of \(\tau\).  Its coefficient must be zero;
iteration eliminates the group.  Repeating over the finite list of exponential rates eliminates all
expanded scalar coefficients.  Thus the continuum maturity response is injective on the reduced
coefficient space.

It remains to pass from the continuum response to a finite panel.  Let \(E\) be the finite
coefficient space and let \(\varepsilon_\tau:E\to\R\) be evaluation at maturity \(\tau\).  If no
finite family of evaluations separated \(E\), the span of \(\{\varepsilon_\tau:\tau\in I\}\) in
\(E^\ast\) would be a proper subspace.  A nonzero vector in its annihilator would generate a
nonzero dictionary combination vanishing on all of \(I\), contradicting the continuum injectivity.
Thus some finite panel is injective.  A full-rank square minor of the associated collocation matrix
is nonzero at that panel.  Since the entries are analytic in the maturity coordinates, this minor is
a nonzero analytic determinant, and deficient panels lie in its zero set, which has empty interior.
\end{proof}

Option-manifold recovery of a route coordinate is not high-frequency market admissibility.  The
power-exponential chart may recover smooth or semi-smooth option-route coordinates on the public
surface.  Rough market admissibility is the smaller source-screened class with local
high-frequency normalization and \(H_{ij}<1/2\).  Recovered option routes enter the rough market
Perron layer only after that additional admissibility condition is imposed.

\begin{corollary}[Rough exponent as a short-maturity invariant]\label{cor:rough-log-derivative}
If a route response has leading form
\[
        f(\tau)=A\tau^\alpha e^{-\beta\tau}(1+c_1\tau+O(\tau^2)),
        \qquad A\ne0,
\]
then
\[
        \alpha
        =
        \lim_{\tau\downarrow0}\tau\frac{f'(\tau)}{f(\tau)}
\]
and
\[
        \beta-c_1
        =
        -\lim_{\tau\downarrow0}
        \left(\frac{f'(\tau)}{f(\tau)}-\frac{\alpha}{\tau}\right).
\]
Thus the rough exponent is an observable short-maturity invariant whenever the leading route is
separated.
\end{corollary}

\begin{proof}
Since \(A\ne0\), the response has fixed sign for all sufficiently small positive maturities.  On
that interval,
\[
        \log |f(\tau)|
        =
        \log |A|+\alpha\log\tau-\beta\tau+c_1\tau+O(\tau^2).
\]
Differentiation gives
\[
        \frac{f'(\tau)}{f(\tau)}
        =
        \frac{\alpha}{\tau}-\beta+c_1+O(\tau),
\]
and the two limits follow.
\end{proof}


\begin{definition}[Surface coordinates generated by covariance routes]\label{def:surface-functional-forms}
Fix a pricing chart \(\chi\in\{\mathrm B,\mathrm N\}\).  Let \(F_{i,t,T}\) be the public forward
of name \(i\) to \(T=t+\tau\), and let an index or basket \(u\) have holdings
\(w^u=(w_1^u,\dots,w_N^u)\) and forward \(F_{u,t,T}=\sum_iw_i^uF_{i,t,T}\).  The route covariance
matrix \(V_t^\QQ(\tau)\) is written in normalized forward-return coordinates.  The chart-specific
exposure vector is
\[
        \alpha_i^{u,\mathrm B}(t,T)
        :=
        \frac{w_i^uF_{i,t,T}}{F_{u,t,T}},
        \qquad
        \alpha_i^{u,\mathrm N}(t,T)
        :=
        w_i^uF_{i,t,T}.
\]
For a single name these formulas are read with the corresponding unit holding.  The local
chart-specific total-variance coordinate is
\[
        \mathsf{TV}_{u,t}^{\chi}(\tau)
        =
        \mathsf{TV}^{\chi,0}_{u,t}(\tau)
        +
        \sum_{i,j=1}^N
        \alpha_i^{u,\chi}(t,T)\alpha_j^{u,\chi}(t,T)
        V_{ij,t}^\QQ(\tau).
\]
In the Black chart \(\mathsf{TV}^{\mathrm B}\) is dimensionless lognormal-forward total
variance, while in the Bachelier chart \(\mathsf{TV}^{\mathrm N}\) is normal-forward total
variance in squared cash-forward units.  The corresponding level coordinates are
\[
        \mathsf{IVL}_{u,t}^{\mathrm B}(\tau)
        =
        \sqrt{\mathsf{TV}_{u,t}^{\mathrm B}(\tau)/\tau},
        \qquad
        \mathsf{NVL}_{u,t}^{\mathrm N}(\tau)
        =
        \sqrt{\mathsf{TV}_{u,t}^{\mathrm N}(\tau)/\tau}.
\]
The local moneyness expansion of the surface in chart \(\chi\) is
\[
        v_{u,t}^{\chi}(\tau,m^\chi)
        =
        \mathsf{TV}_{u,t}^{\chi}(\tau)
        +
        \mathsf{SKW}_{u,t}^{\chi}(\tau)m^\chi
        +
        \frac12\mathsf{CURV}_{u,t}^{\chi}(\tau)(m^\chi)^2
        +
        O((m^\chi)^3),
\]
where the downside-skew and curvature coordinates are finite linear functionals of
\(\{V_{ij,t}^\QQ(\tau)\}_{i,j}\), for example
\[
        \mathsf{SKW}_{u,t}^{\chi}(\tau)
        =
        s^{\chi,0}_{u,t}(\tau)
        +
        \sum_{i,j} b^{\chi,\mathrm{skw}}_{u,ij,t}(\tau)V_{ij,t}^\QQ(\tau).
\]
For a call or put coordinate \(a=(u,\tau,m^\chi,\kappa,\chi)\), the local pricing map is obtained
by substituting \(v_{u,t}^{\chi}(\tau,m^\chi)\) in the corresponding no-arbitrage one-period
price chart:
\[
        C_t^\Theta(a)
        =
        \mathsf B_\kappa^\chi
        \!\left(D_{t,T},F_{u,t,T},K^\chi(m^\chi),v_{u,t}^{\chi}(\tau,m^\chi)\right).
\]
The construction is a local response chart around an arbitrage-free public surface.  It does not assert
a global parametrization of every strike and maturity.  The recovery neighborhood is chosen inside the
admitted no-arbitrage chart, so positivity, monotonicity, convexity, and calendar constraints
remain valid on the finite or continuum panel under consideration.
The two total-variance coordinates are chart-specific objects.  A lognormal total variance and a
normal-forward variance are not numerically interchangeable without an additional conversion
model.  The rank statements below use only the local price chart and its nonzero variance vega.
\end{definition}

\begin{definition}[Black and Bachelier price charts]\label{def:black-bachelier-charts}
Let \(D>0\) be a discount factor, \(F>0\) a forward, \(K\) a strike, and \(v>0\) a total-variance
coordinate in the chosen chart.  These are the Bachelier normal-forward and
Black--Scholes--Merton lognormal-forward charts \cite{bachelier1900,blackScholes1973,merton1973}.
Let \(\Phi_0\) and \(\varphi_0\) denote the standard normal distribution function and density.  In
the Black chart,
\[
        d_\pm=\frac{\log(F/K)}{\sqrt v}\pm\frac12\sqrt v,
\]
\[
        \mathsf B_C^{\mathrm B}(D,F,K,v)
        =
        D\bigl(F\Phi_0(d_+)-K\Phi_0(d_-)\bigr),
        \qquad
        \mathsf B_P^{\mathrm B}(D,F,K,v)
        =
        D\bigl(K\Phi_0(-d_-)-F\Phi_0(-d_+)\bigr).
\]
In the Bachelier chart,
\[
        d=\frac{F-K}{\sqrt v},
\]
\begin{align*}
        \mathsf B_C^{\mathrm N}(D,F,K,v)
        =
        D\bigl((F-K)\Phi_0(d)+\sqrt v\,\varphi_0(d)\bigr),\\
        \mathsf B_P^{\mathrm N}(D,F,K,v)
        =
        D\bigl((K-F)\Phi_0(-d)+\sqrt v\,\varphi_0(d)\bigr).
\end{align*}
Both charts satisfy put--call parity
\[
        \mathsf B_C^{\chi}(D,F,K,v)-\mathsf B_P^{\chi}(D,F,K,v)=D(F-K).
\]
\end{definition}

\begin{lemma}[Variance vegas]\label{lem:variance-vegas}
In the Black chart,
\[
        \partial_v\mathsf B_C^{\mathrm B}
        =
        \partial_v\mathsf B_P^{\mathrm B}
        =
        \frac{DF\varphi_0(d_+)}{2\sqrt v}.
\]
In the Bachelier chart,
\[
        \partial_v\mathsf B_C^{\mathrm N}
        =
        \partial_v\mathsf B_P^{\mathrm N}
        =
        \frac{D\varphi_0(d)}{2\sqrt v}.
\]
\end{lemma}

\begin{proof}
In the Black chart,
\[
        \partial_v d_+=-\frac{\log(F/K)}{2v^{3/2}}+\frac{1}{4\sqrt v},
        \qquad
        \partial_v d_- = \partial_v d_+-\frac{1}{2\sqrt v}.
\]
The density identity
\[
        F\varphi_0(d_+)=K\varphi_0(d_-)
\]
cancels the \(d_\pm\)-terms generated by differentiating the two distribution functions.  The
remaining Black contribution is
\[
        \frac{D F\varphi_0(d_+)}{2\sqrt v}.
\]
In the Bachelier chart, \(d=(F-K)/\sqrt v\).  The derivatives of
\((F-K)\Phi_0(d)\) and \(\sqrt v\varphi_0(d)\) cancel except for
\(\varphi_0(d)/(2\sqrt v)\).  Put prices have the same variance derivative by put--call parity,
since \(D(F-K)\) does not depend on \(v\).
\end{proof}

\begin{theorem}[Pricing-chart rank invariance]\label{thm:pricing-chart-rank-invariance}
Let \(u\mapsto v(u)\in\R^M\) be a finite vector of chart-specific total-variance coordinates,
and suppose public discount factors, forwards, strikes, maturities, and option types are fixed.
Let
\[
        C_r(u)
        =
        \mathsf B_{\kappa_r}^{\chi_r}
        \!\left(D_r,F_r,K_r,v_r(u)\right),
        \qquad r=1,\dots,M.
\]
If the variance vegas \(w_r:=\partial_v\mathsf B_{\kappa_r}^{\chi_r}(D_r,F_r,K_r,v_r(u_0))\)
satisfy \(0<\underline w\le |w_r|\le\bar w<\infty\) at \(u_0\), then
\[
        DC(u_0)=W\,Dv(u_0),
        \qquad
        W=\operatorname{diag}(w_1,\dots,w_M).
\]
Consequently \(DC(u_0)\) and \(Dv(u_0)\) have the same rank, the same kernel, and comparable
nonzero singular values with constants depending only on \(\underline w\) and \(\bar w\).
\end{theorem}

\begin{proof}
The chain rule and Lemma~\ref{lem:variance-vegas} give the diagonal factorization.  Multiplying
by an invertible diagonal matrix preserves rank and kernel.  The singular-value comparison
follows from
\[
        \underline w\,\|Dv(u_0)h\|
        \le
        \|DC(u_0)h\|
        \le
        \bar w\,\|Dv(u_0)h\|.
\]
\end{proof}

\begin{lemma}[Parity projection and no-arbitrage residuals]\label{lem:parity-projection}
For fixed public \(D,F,K\), the call--put parity row
\[
        C-P-D(F-K)=0
\]
is a static public no-arbitrage restriction, not an independent Volterra--Perron source
response.  Recovery ranks may therefore be computed after projecting quoted call--put pairs onto
the parity-consistent subspace.  The orthogonal distance to that subspace, in the quote norm of
Definition~\ref{def:quote-space-norm}, is a quote-quality or no-arbitrage-cleaning statistic
rather than evidence for a latent residual kernel.
\end{lemma}

\begin{proof}
The identity is deterministic once \(D,F,K\) are fixed.  It does not depend on the latent source
chart, hence its differential annihilates all source directions.  Projecting away the parity row
removes a redundant public coordinate and leaves the source-response rank unchanged.  A parity
violation is therefore a violation of the admitted public price cone, not a new latent pricing
direction.
\end{proof}

\begin{definition}[Screened synchronization and carrier response]\label{def:screened-sync-functional}
Fix a synchronization maturity \(\tau_*\).  Let \(B_t^\QQ\) denote the signed directed Volterra
route matrix in the chosen source chart.  Its entries are the active directed route coefficients
whose public covariance-route images are \(V_{ij,t}^\QQ(\tau)\).  The screened risk-neutral
synchronization matrix is the positive option-manifold image
\[
        A_t^\QQ=\Phi_t(B_t^\QQ).
\]
On the active screened stratum, \(\Phi_t\) is the differentiable map with off-diagonal entries
\[
        A_{ij,t}^\QQ
        =
        q_{ij,t}
        \frac{\bigl(V_{ij,t}^\QQ(\tau_*)\bigr)_+}
        {\sqrt{(V_{ii,t}^\QQ(\tau_*)+\varepsilon)
        (V_{jj,t}^\QQ(\tau_*)+\varepsilon)}},
        \qquad i\ne j,
\]
and \(A_{ii,t}^\QQ=0\).  Here \(q_{ij,t}\ge0\) is a source-screening weight and
\(\varepsilon>0\) is fixed.  The normalization by the diagonal variance terms makes
\(A_t^\QQ\) a dimensionless screened one-step synchronization operator; Perron critical surfaces
therefore refer to this normalized object rather than to an arbitrary scale of option variance.
Passing from \(B_t^\QQ\) to \(A_t^\QQ\) loses sign information and retains the positive
synchronization geometry observed on the option manifold.  Here the screened synchronization
operator is the finite
matrix \(A_t^\QQ\in\R_+^{N\times N}\) on the active underlier set.  Infinite-dimensional
Krein--Rutman variants would require compact strong positivity assumptions and are not used here.
If \(A_t^\QQ\) is irreducible with simple Perron root, then
\[
        A_t^\QQ e_t^\QQ=\rho_t^\QQ e_t^\QQ,\qquad
        \ell_t^\QQ A_t^\QQ=\rho_t^\QQ\ell_t^\QQ,\qquad
        \ell_t^\QQ(e_t^\QQ)=1.
\]
Given carrier portfolios \(b^1,\dots,b^r\in\R^N\), their synchronization response vectors are
\[
        \mathcal C_k(\tau)
        =
        \sum_{i,j}b_i^k b_j^k V_{ij,t}^\QQ(\tau),
        \qquad k=1,\dots,r,
\]
and the carrier Gram matrix on a maturity grid \(\mathcal T\) is
\[
        \Gamma_{\mathrm{car},kl}
        =
        \sum_{\tau\in\mathcal T}
        \mathcal C_k(\tau)\mathcal C_l(\tau).
\]
Carrier absorption is nondegenerate when
\(\Gamma_{\mathrm{car}}\succeq\kappa_{\mathrm{car}}I\).
\end{definition}

\begin{lemma}[Screened-edge derivative]\label{lem:screened-edge-derivative}
On an active screened stratum where \(V_{ij,t}^\QQ(\tau_*)>0\) and \(q_{ij,t}\) is differentiable,
write
\[
        A_{ij}
        =
        q_{ij}\frac{V_{ij}}{D_{ij}},
        \qquad
        D_{ij}:=\sqrt{(V_{ii}+\eps)(V_{jj}+\eps)} .
\]
Then
\[
        \delta A_{ij}
        =
        \frac{q_{ij}}{D_{ij}}\delta V_{ij}
        -
        \frac{A_{ij}}2
        \left(
        \frac{\delta V_{ii}}{V_{ii}+\eps}
        +
        \frac{\delta V_{jj}}{V_{jj}+\eps}
        \right)
        +
        \frac{V_{ij}}{D_{ij}}\delta q_{ij}.
\]
If \(q_{ij}\) is fixed, the last term is absent.  The fixed-sign margin in
Proposition~\ref{prop:functional-rank} is exactly the condition that keeps this derivative away
from the kink of the positive-part screening map.
\end{lemma}

\begin{proof}
On the active stratum the positive-part map has derivative one, so
\(A_{ij}=q_{ij}V_{ij}D_{ij}^{-1}\).  Differentiating this quotient gives
\[
        \delta A_{ij}
        =
        \frac{q_{ij}}{D_{ij}}\delta V_{ij}
        +
        \frac{V_{ij}}{D_{ij}}\delta q_{ij}
        -
        A_{ij}\frac{\delta D_{ij}}{D_{ij}}.
\]
Since
\[
        \frac{\delta D_{ij}}{D_{ij}}
        =
        \frac12
        \left(
        \frac{\delta V_{ii}}{V_{ii}+\eps}
        +
        \frac{\delta V_{jj}}{V_{jj}+\eps}
        \right),
\]
the displayed formula follows.
\end{proof}

\begin{definition}[Interval screened matrix]\label{def:interval-screened-matrix}
Suppose finite-panel recovery or bid--ask uncertainty gives elementwise bounds on the
nonnegative screened synchronization matrix,
\[
        L_{ij}\le A_{ij}\le U_{ij},\qquad L,U\ge0 .
\]
The compatible matrix set is
\[
        \cA[L,U]:=\{A\ge0:L\le A\le U\text{ elementwise}\}.
\]
\end{definition}

Perron--Frobenius terminology and matrix perturbation conventions are used in their
finite-dimensional sense \cite{bermanPlemmons1994,hornJohnson2012}.

\begin{theorem}[Certified Perron phase bounds]\label{thm:certified-perron-phase}
For every \(A\in\cA[L,U]\),
\[
        \rho(L)\le \rho(A)\le \rho(U),
\]
where \(\rho(\cdot)\) denotes spectral radius.  For any strictly positive vector \(x\),
\[
        \min_i\frac{(Lx)_i}{x_i}
        \le
        \rho(A)
        \le
        \max_i\frac{(Ux)_i}{x_i}.
\]
Consequently, if \(\rho(U)<\rho_{\mathrm{crit}}\), the recovered state is certified below
the synchronization threshold for every admissible screened matrix; if
\(\rho(L)>\rho_{\mathrm{crit}}\), it is certified above the threshold.  Otherwise the phase
label is not identified at the available quote precision.
\end{theorem}

\begin{proof}
Elementwise order gives \(L\le A\le U\).  Monotonicity of the Perron root for nonnegative matrices
therefore yields \(\rho(L)\le\rho(A)\le\rho(U)\).  For any strictly positive \(x\), the
Collatz--Wielandt inequalities give
\[
        \min_i\frac{(Lx)_i}{x_i}
        \le \rho(L),
        \qquad
        \rho(U)\le
        \max_i\frac{(Ux)_i}{x_i},
\]
and hence the displayed bound for every compatible \(A\).  The phase certification is the direct
comparison of these bounds with \(\rho_{\mathrm{crit}}\).
\end{proof}

\begin{definition}[Bounded residual selector]\label{def:functional-residual-selector}
Let \(R^1,\dots,R^m\) be bounded residual directions satisfying
\[
        \EE^\QQ[R^r\mid\cH_t]=0,
        \qquad
        \EE^\QQ[R^r X_a\mid\cH_t]=0
        \quad\text{for every admitted public payoff }X_a .
\]
For recovered state coordinates
\[
        z_t
        =
        \left(
        \varpi_t^\QQ,\log\rho_t^\QQ,\bar{\mathfrak{H}}_t^\QQ,
        \sum_{i,j}\ell_{i,t}^\QQ G_{ij,t}^\QQ e_{j,t}^\QQ
        \right),
\]
define
\[
        \theta_{r,t}
        =
        \bar\theta_{r,t}\,
        \tanh\!\left(\eta_{r,0,t}+\eta_{r,t}^{\top}z_t\right),
        \qquad
        \sum_{r=1}^m |\bar\theta_{r,t}|\norm{R^r}_\infty<1.
\]
The associated public-equivalent density family is
\[
        Z_{t,S}^{\theta}
        =
        1+\sum_{r=1}^m \theta_{r,t}R^r .
\]
This form keeps positivity and public equivalence explicit while allowing residual prices to
depend on recovered Perron stress, spectral radius, roughness, and directed-route pressure.  The rule
operates after state recovery; it becomes a public-response coordinate only when augmented residual
instruments are admitted.
\end{definition}

\begin{proposition}[Functional-form rank criterion]\label{prop:functional-rank}
In the source chart of Definitions~\ref{def:volterra-route-kernels}--\ref{def:functional-residual-selector},
suppose:
\begin{enumerate}[leftmargin=2em]
\item the maturity-response atoms used for recovery form a reduced, linearly independent
power--exponential dictionary after the chosen gauge normalizations;
\item the corresponding route amplitudes \(A_{ijq,t}\) are nonzero;
\item the active screened edges satisfy a fixed-sign margin
\[
        |V_{ij,t}^\QQ(\tau_*)|\ge v_*>0
\]
on the recovery neighborhood, while inactive edges have fixed screening weight \(q_{ij,t}=0\);
\item the index and carrier weight systems separate the cross-covariance routes entering
\(\{V_{ij,t}^\QQ\}\);
\item the screened synchronization matrix has a simple Perron root and the carrier Gram matrix
has a positive lower bound.
\end{enumerate}
Then the public observable response derivative for the canonical state is block triangular with
full-rank diagonal blocks.
Consequently the block-rank observability criterion and the conditional recovery theorem apply.
If an augmented manifold also contains residual-price instruments, the same criterion applies to
the additional selector block; for a purely public option manifold the selector block is zero.
\end{proposition}

\begin{proof}
The route kernels generate analytic maturity functions in a reduced power--exponential dictionary.
These atoms are linearly independent on any open maturity interval, and on a sufficiently resolving
finite panel by Theorem~\ref{thm:finite-panel-recovery}.  The nonzero route amplitudes make the
directed Volterra block visible.  The fixed-sign margin keeps the positive-part screening map
differentiable on the active stratum and fixes inactive edges outside the quotient tangent space.
The index and carrier weight conditions make the linear map from cross-covariance routes to
index/carrier responses injective on the admitted source coordinates.  The simple,
well-conditioned Perron eigenvalue gives differentiable recovery of the spectral triple, while the
carrier Gram lower bound gives a bounded inverse for carrier absorption coordinates.

After gauge directions have been removed, any public kernel must lie in one of the diagonal blocks
or in a selector coordinate not captured by the pure public response.  The stated rank conditions
exclude the diagonal-block kernels.  Public-annihilating residual selectors remain outside the pure
public response derivative unless residual instruments are appended.  The block operator is
therefore triangular with full-rank diagonal blocks, and Theorem~\ref{thm:block-rank} applies.
\end{proof}

\begin{theorem}[Standard-model recovery mechanism]\label{thm:standard-model-recovery}
In the source chart of this section, suppose the rank conditions of
Proposition~\ref{prop:functional-rank} hold on a recovery neighborhood.  Then the admitted option
manifold locally recovers the canonical source state
\[
        (\mathfrak{H}^\QQ,G^\QQ,\rho^\QQ,e^\QQ,\ell^\QQ,P^\QQ)
\]
up to gauge.  The inverse is assembled in four conceptual blocks:
\[
\begin{aligned}
        \text{maturity scaling}
        &\Longrightarrow
        \mathfrak{H}^\QQ\text{ and the Volterra source dictionary},\\
        \text{basket/carrier polarization}
        &\Longrightarrow
        G^\QQ\text{ and the active covariance routes},\\
        \text{screened synchronization}
        &\Longrightarrow
        (\rho^\QQ,e^\QQ,\ell^\QQ,P^\QQ),\\
        \text{residual completion or selector choice}
        &\Longrightarrow
        \theta\text{ only outside the pure public response.}
\end{aligned}
\]
Consequently the pure public option manifold can recover the Volterra--Perron source state on this
stratum.  Residual-selector recovery requires public completion, augmented residual instruments, or
a separate canonical selector rule.
\end{theorem}

\begin{proof}
The maturity and route blocks are injective by the analytic dictionary and carrier-separation
conditions in Proposition~\ref{prop:functional-rank}.  On the simple-Perron stratum, the screened
synchronization block has a differentiable spectral chart whose derivative is controlled by the
reduced resolvent; the carrier Gram lower bound gives a bounded inverse for the carrier absorption
coordinates.  Hence the quotient derivative has the block-triangular full-rank form of
Definition~\ref{def:block-source-chart}.  Theorem~\ref{thm:block-rank} gives a bounded left
inverse, and the Volterra--Perron option-manifold equivalence theorem gives local recovery of the
canonical source state.  Because the residual selector lies outside the pure public response block, it is recovered only under
the additional residual-completion alternatives stated in the theorem.
\end{proof}

\section{Risk-neutral market phases and transition surfaces}\label{sec:phases}

The market-phase theory of \cite{vidal2026market} is imposed here on the recovered risk-neutral
quotient.  Phases are not exogenous regimes.  They are functions of the recovered option-manifold
state, recording whether public prices represent local risk, synchronized carrier-absorbed risk,
residual stress, or synchronized risk without an adequate carrier span.

\begin{definition}[Risk-neutral phase coordinates]\label{def:phase-coordinates}
In the source chart of Section~\ref{sec:functional-forms}, use the Perron spectral radius
\(\rho_t^\QQ\) and define the remaining phase coordinates
\[
        K_t^{\mathrm{car}}:=\lambda_{\min}(\Gamma_{\mathrm{car},t}),
\]
\[
        R_t^\theta
        :=
        \left(
        \EE^\QQ\!\left[
        \left(\sum_{r=1}^m\theta_{r,t}R^r\right)^2
        \middle|\cH_t\right]
        \right)^{1/2}
        =
        \left(\theta_t^\top G_t^R\theta_t\right)^{1/2},
        \qquad
        (G_t^R)_{rs}:=\EE^\QQ[R^rR^s\mid\cH_t],
\]
and
\[
        D_t^\QQ
        :=
        \sum_{i,j}
        \ell_{i,t}^\QQ\,q_{ij,t}\,
        a_{ij,t}\,
        e_{j,t}^\QQ .
\]
Here \(\rho_t^\QQ\) measures synchronized spectral amplification, \(K_t^{\mathrm{car}}\)
measures carrier absorption, \(R_t^\theta\) measures residual-kernel pressure with the full
conditional residual Gram matrix, and
\(D_t^\QQ\) measures directed Perron-weighted route pressure.
\end{definition}

\begin{remark}
For smooth phase weights, the recovery neighborhood is chosen so that the smallest carrier Gram
eigenvalue is simple and \(R_t^\theta>0\), or else the smooth scores below are read with the
regularized substitutes \(\log\det(\Gamma_{\mathrm{car},t}+\eps I)\) and
\(\theta_t^\top G_t^R\theta_t\).  The hard phase labels only require the continuous
nonnegative coordinates.
\end{remark}

\begin{definition}[Hard phase regions]\label{def:hard-phases}
Fix public thresholds
\[
        \rho_*>0,\qquad K_*>0,\qquad R_*>0.
\]
The recovered state is in the local phase if
\[
        \rho_t^\QQ<\rho_*,
        \qquad
        R_t^\theta<R_*.
\]
It is in the synchronized-carrier phase if
\[
        \rho_t^\QQ\ge\rho_*,
        \qquad
        K_t^{\mathrm{car}}\ge K_*,
        \qquad
        R_t^\theta<R_*.
\]
It is in the residual-stress phase if
\[
        \rho_t^\QQ\ge\rho_*,
        \qquad
        R_t^\theta\ge R_*.
\]
It is in the carrier-thin phase if
\[
        \rho_t^\QQ\ge\rho_*,
        \qquad
        K_t^{\mathrm{car}}<K_*.
\]
When regions overlap, residual-stress and carrier-thin labels take precedence because they mark
failure of public carrier absorption or enlargement of the residual price envelope.
\end{definition}

The thresholds in Definition~\ref{def:hard-phases} are not fitted trading cutoffs.  They are
coordinate surfaces on the recovered \(\QQ\)-state.  The spectral surface separates local and
synchronized option geometry; the carrier surface separates synchronized risk absorbed by the
admitted carrier span from synchronized risk that remains thinly spanned; the residual surface
marks expansion of public-equivalent price freedom.

The intrinsic phase atlas precedes any positive finite-resolution band.  Its quotient seams are the
surfaces where an observability or pricing mechanism changes type: normalized Perron amplification
crosses its critical level, the Perron spectral gap collapses, the carrier Gram lower bound reaches
zero, a maturity or route-rank determinant vanishes, or the residual price envelope collapses.
Positive constants \((\rho_*,K_*,R_*)\) thicken these seams into certificate regions for finite
panels.  Changing the bands changes the resolution of the certificate, not the underlying quotient
stratification.

\begin{theorem}[Risk-neutral phase stratification]\label{thm:phase-stratification}
On a recovery stratum, the phase labels in Definition~\ref{def:hard-phases} are gauge-invariant
functions of the recovered source quotient.  Their transition surfaces are quotient-invariant
rank, gap, Gram, and residual-width surfaces:
\[
        \rho^\QQ=\rho_*,
        \qquad
        \lambda_{\min}(\Gamma_{\mathrm{car}})=K_*,
        \qquad
        R^\theta=R_* .
\]
If the recovered state process is adapted and continuous, phase changes can occur only by hitting
one of these surfaces.  Thus local, synchronized-carrier, residual-stress, and carrier-thin phases
are not additional latent regimes; they are strata of the recovered risk-neutral option quotient.
\end{theorem}

\begin{proof}
The coordinates \(\rho^\QQ\), \(\Gamma_{\mathrm{car}}\), and \(R^\theta\) are defined from the
canonical recovered state and the selected residual rule, so they are invariant under the gauge
normalizations used in the source quotient.  The phase regions are inverse images of Borel subsets
of these quotient coordinates.  If the recovered state is adapted and continuous, the phase
coordinates are adapted and continuous.  A path cannot move between open phase regions without
crossing the boundary surface that separates them.
\end{proof}

\begin{remark}[Spectral-gap seams]\label{rem:spectral-gap-seams}
The surface \(\Delta_P=0\), where the Perron spectral gap closes, should be read as a spectral
discriminant seam rather than as an ordinary smooth phase boundary.  On finite self-adjoint or
symmetry-reduced subfamilies such seams may appear as familiar eigenvalue-intersection loci; in
non-self-adjoint directed Perron operators their codimension depends on the admitted source
family, positivity constraints, and whether the crossing is an eigenvalue collision or a
dominance exchange in modulus.  The theory uses only the invariant fact that the reduced resolvent
and the differential of the Perron projector lose uniform bounds as \(\Delta_P\downarrow0\), and
may also lose bounds when non-normal conditioning degenerates; no universal codimension assignment
is required.  Thus a path hitting this seam exits the smooth recovered phase chart
and must be continued, if possible, in another chart of the recovery atlas.
\end{remark}

\begin{definition}[Smooth phase weights]\label{def:smooth-phase-weights}
For differentiable pricing, replace the hard phase labels by smooth weights.  Define scores
\[
        s_\rho=\rho_t^\QQ-\rho_*,
        \qquad
        s_K=K_t^{\mathrm{car}}-K_*,
        \qquad
        s_R=R_t^\theta-R_*,
        \qquad
        s_D=D_t^\QQ-D_*,
\]
and
\[
        m_{\mathrm{loc}}=-s_\rho-s_R,\qquad
        m_{\mathrm{syn}}=s_\rho+s_K-s_R+\eta_D s_D,
\]
\[
        m_{\mathrm{res}}=s_\rho+s_R+\eta_D s_D,\qquad
        m_{\mathrm{thin}}=s_\rho-s_K .
\]
Here \(D_*\) is a directed-pressure reference level and \(\eta_D\ge0\) controls how strongly
directed route pressure tilts synchronized and residual phases.  For a sharpness parameter
\(\rho>0\), set
\[
        \pi_{\nu,t}
        =
        \frac{\exp(\rho m_{\nu,t})}
        {\sum_{\mu\in\{\mathrm{loc},\mathrm{syn},\mathrm{res},\mathrm{thin}\}}
        \exp(\rho m_{\mu,t})},
        \qquad
        \nu\in\{\mathrm{loc},\mathrm{syn},\mathrm{res},\mathrm{thin}\}.
\]
The vector \(\pi_t\) is the smooth risk-neutral phase allocation induced by the recovered
Volterra--Perron state.
\end{definition}

\begin{definition}[Phase-weighted residual selector]\label{def:phase-weighted-selector}
For each \(\nu\in\{\mathrm{loc},\mathrm{syn},\mathrm{res},\mathrm{thin}\}\), let
\[
        Z_{t,S}^{\nu}
        =
        1+\sum_r\theta_{r,t}^{\nu}R^r
\]
be a positive public-equivalent density family, indexed by maturity \(S>t\).  The
phase-weighted family is
\[
        Z_{t,S}^{\pi}
        :=
        \sum_{\nu}
        \pi_{\nu,t}Z_{t,S}^{\nu}.
\]
The corresponding selected price is
\[
        \Pi_t^\pi(X):=\EE^\QQ[Z_{t,T}^{\pi}X\mid\cH_t].
\]
\end{definition}

\begin{proposition}[Phase transitions are adapted]\label{prop:phase-transitions-adapted}
Suppose the recovered state process \(\Xi_t^\QQ\), together with the selected
residual-density rule \(\theta_t\), is \(\cH_t\)-adapted and the phase
coordinates in Definition~\ref{def:phase-coordinates} are continuous functions of the canonical
state.  Then the hard phase label is \(\cH_t\)-adapted, and the first hitting times of the phase
transition surfaces
\[
        \{\rho^\QQ=\rho_*\},\qquad
        \{K^{\mathrm{car}}=K_*\},\qquad
        \{R^\theta=R_*\}
\]
are \((\cH_t)\)-stopping times.  If the recovered state is continuous, phase changes occur only
by hitting one of these transition surfaces.
\end{proposition}

\begin{proof}
The phase coordinates are adapted as continuous functions of the adapted recovered state.  The
hard phase regions are Borel sets in the canonical state chart, so the phase label is adapted.
The transition surfaces are closed level sets of adapted continuous processes.  Their first
hitting times are therefore stopping times.  If the recovered state path is continuous, the phase
label cannot move from one open region to another without crossing the boundary separating them.
\end{proof}

\begin{proposition}[No-arbitrage preservation under phase weighting]\label{prop:phase-weighting-noarbitrage}
If each \(Z_{t,\cdot}^{\nu}\) in Definition~\ref{def:phase-weighted-selector} is positive and
public-equivalent,
\[
        \EE^\QQ[Z_{t,S}^{\nu}\mid\cH_t]=1,
        \qquad S>t,
\]
and
\[
        \EE^\QQ[Z_{t,T(a)}^{\nu}X_a\mid\cH_t]
        =
        \EE^\QQ[X_a\mid\cH_t]
        \quad
        \text{for every admitted public payoff }X_a,
\]
and the weights \(\pi_{\nu,t}\) are nonnegative, \(\cH_t\)-measurable, and sum to one, then
\[
        Z_{t,S}^{\pi}
        =
        \sum_\nu\pi_{\nu,t}Z_{t,S}^{\nu}
\]
is positive and public-equivalent.  Hence phase-weighted residual pricing preserves all public
option prices while allowing the latent residual charge to vary with the recovered market phase.
\end{proposition}

\begin{proof}
Positivity follows from positivity of the kernels and nonnegativity of the weights.  Since the
weights are \(\cH_t\)-measurable,
\[
        \EE^\QQ[Z_{t,S}^{\pi}\mid\cH_t]
        =
        \sum_\nu
        \pi_{\nu,t}\EE^\QQ[Z_{t,S}^{\nu}\mid\cH_t]
        =
        \sum_\nu\pi_{\nu,t}
        =
        1 .
\]
For any admitted public payoff \(X_a\),
\[
        \EE^\QQ[Z_{t,T(a)}^{\pi}X_a\mid\cH_t]
        =
        \sum_\nu
        \pi_{\nu,t}\EE^\QQ[Z_{t,T(a)}^{\nu}X_a\mid\cH_t]
        =
        \EE^\QQ[X_a\mid\cH_t].
\]
Thus the kernel is public-equivalent.  Public payoff prices are unchanged, while residual payoffs
can receive phase-dependent charges through their covariance with the selected residual
directions.
\end{proof}

\begin{remark}[Datewise preservation versus process consistency]\label{rem:phase-datewise-process}
Proposition~\ref{prop:phase-weighting-noarbitrage} is a datewise preservation result.  It says
that, at a fixed date and horizon, convex phase weighting preserves the admitted public payoff
prices.  It does not by itself produce a single density process whose horizon ratios generate all
the \(Z_{t,S}^{\pi}\) consistently across \(t\) and \(S\).  Dynamic no-arbitrage requires the
stronger process-consistent kernel of Definition~\ref{def:process-consistent-kernel} and the
martingale restrictions of Section~\ref{sec:martingale}.  Thus phase weighting is a public
residual-selection rule until it is embedded in a martingale density process.
\end{remark}

\begin{proposition}[Semimartingale phase dynamics]\label{prop:semimartingale-phase-dynamics}
Suppose the finite-dimensional recovered chart \(U_t\) satisfies the semimartingale dynamics in
Assumption~\ref{ass:semimartingale-chart}, and suppose the score maps
\[
        U\mapsto m_{\nu}(U),
        \qquad
        \nu\in\{\mathrm{loc},\mathrm{syn},\mathrm{res},\mathrm{thin}\},
\]
are \(C^2\) on the recovery neighborhood.  Then the smooth phase weights
\(\pi_{\nu,t}\) are \(\cH_t\)-semimartingales.  Their drift and diffusion coefficients are given
by It\^o's formula applied to the softmax map in Definition~\ref{def:smooth-phase-weights}.
\end{proposition}

\begin{proof}
The recovered chart is an \(\cH_t\)-semimartingale.  The score vector is a \(C^2\) function of
that chart, and the softmax map is \(C^\infty\).  The composition is therefore an
\(\cH_t\)-semimartingale, with coefficients given by It\^o's formula.
\end{proof}

\begin{definition}[Phase value functions]\label{def:phase-value-functions}
Let \(U_t\) be a finite-dimensional recovered state chart that is Markov for the selected claim
family under the phase kernels below.  For a claim \(X\) and phase
\(\nu\in\{\mathrm{loc},\mathrm{syn},\mathrm{res},\mathrm{thin}\}\), define a regular conditional
version
\[
        \Phi_\nu(t,u;X)
        :=
        \EE^\QQ[Z_{t,T}^{\nu}X\mid U_t=u],
        \qquad
        \Phi_\pi(t,u;X)
        :=
        \sum_\nu \pi_\nu(u)\Phi_\nu(t,u;X).
\]
The phase transition charge relative to a reference phase \(\nu_0\) is
\[
        \mathfrak C_{\nu_0}^{\mathrm{ph}}(t,u;X)
        :=
        \Phi_\pi(t,u;X)-\Phi_{\nu_0}(t,u;X)
        =
        \sum_\nu \pi_\nu(u)\bigl(\Phi_\nu(t,u;X)-\Phi_{\nu_0}(t,u;X)\bigr).
\]
\end{definition}

\begin{definition}[Transition Greeks]\label{def:transition-greeks}
Let \(\xi\) be one of the phase scores \(s_\rho,s_K,s_R\).  The pure allocation transition
Greek of \(X\) in the \(\xi\)-direction is
\[
        \mathfrak T_{\xi}^{\mathrm{alloc}}X
        :=
        \sum_\nu
        \frac{\partial \pi_\nu}{\partial \xi}
        \Phi_\nu(t,U_t;X).
\]
Using the softmax weights of Definition~\ref{def:smooth-phase-weights},
\[
        \mathfrak T_{\xi}^{\mathrm{alloc}}X
        =
        \rho
        \sum_\nu
        \pi_{\nu,t}
        \left(
        \partial_\xi m_{\nu,t}
        -
        \sum_\mu\pi_{\mu,t}\partial_\xi m_{\mu,t}
        \right)
        \Phi_\nu(t,U_t;X).
\]
The full transition Greek adds the direct state sensitivity
\[
        \mathfrak T_\xi X
        :=
        \mathfrak T_{\xi}^{\mathrm{alloc}}X
        +
        \sum_\nu \pi_{\nu,t}\,\partial_\xi\Phi_\nu(t,U_t;X),
\]
whenever the displayed derivatives exist.
\end{definition}

\begin{theorem}[Phase-weighted risk-neutral pricing equation]\label{thm:phase-pricing-equation}
Suppose the recovered chart satisfies
\[
        \dd U_t=b(t,U_t)\dd t+\Sigma(t,U_t)\dd W_t^\QQ
\]
on a recovery neighborhood, and define
\[
        \cA f
        :=
        Df[b]
        +
        \frac12\operatorname{Tr}\!\left(\Sigma^\ast D^2f\,\Sigma\right).
\]
Let the phase weights \(\pi_\nu\) and phase value functions \(\Phi_\nu\) be \(C^{1,2}\).  If the
discounted phase-weighted selected price is a local \(\QQ\)-martingale, then
\[
        (\partial_t+\cA-r_t)\Phi_\pi(t,U_t;X)=0 .
\]
Equivalently,
\[
        0
        =
        \sum_\nu
        \pi_\nu(\partial_t+\cA-r_t)\Phi_\nu
        +
        \sum_\nu
        \Phi_\nu(\partial_t+\cA)\pi_\nu
        +
        \sum_\nu
        \ip{\Sigma^\ast D\pi_\nu}{\Sigma^\ast D\Phi_\nu},
\]
evaluated at \((t,U_t)\).  The last two sums are the transition drift and transition covariance
charges generated by phase weights that move with the recovered \(\QQ\)-state.
\end{theorem}

\begin{proof}
The first display is It\^o's formula applied to \(e^{-\int_0^t r_s\dd s}\Phi_\pi(t,U_t;X)\)
and the local martingale condition.  For the expanded form, use
\[
        \Phi_\pi=\sum_\nu\pi_\nu\Phi_\nu
\]
and apply the product rule to \((\partial_t+\cA)(\pi_\nu\Phi_\nu)\).  The second-order cross term
is
\[
        \ip{\Sigma^\ast D\pi_\nu}{\Sigma^\ast D\Phi_\nu},
\]
which gives the displayed transition covariance charge.
\end{proof}

\begin{corollary}[Phase charges vanish only under frozen or balanced phases]\label{cor:phase-charge-balance}
If every phase value function separately satisfies
\[
        (\partial_t+\cA-r_t)\Phi_\nu=0,
\]
then the phase-weighted selected price is arbitrage-free only if
\[
        \sum_\nu
        \Phi_\nu(\partial_t+\cA)\pi_\nu
        +
        \sum_\nu
        \ip{\Sigma^\ast D\pi_\nu}{\Sigma^\ast D\Phi_\nu}
        =
        0 .
\]
In particular, changing risk-neutral phase weights are priced objects unless their drift and
diffusion effects are exactly balanced by phase-value spreads.
\end{corollary}

\begin{proof}
Substitute the separate phase pricing equations into
Theorem~\ref{thm:phase-pricing-equation}.  The remaining terms must vanish for the
phase-weighted discounted price to be a local martingale.
\end{proof}

\section{Public non-identification}\label{sec:nonidentification}

Public compression gives the first obstruction.  Full latent recovery requires a source-family
restriction before the public quotient can be inverted.

\begin{definition}[Public indistinguishability]\label{def:public-indistinguishability}
Fix date \(t\) and an admitted option-pricing map
\[
        \mathfrak P_t:\mathfrak X_t\to\cY_t,
        \qquad
        \Theta\mapsto\cM_t^{\mathrm{opt}}(\Theta).
\]
Two latent states are public-indistinguishable, written
\[
        \Theta\sim_{\mathrm{pub},t}\Theta',
\]
if \(\mathfrak P_t(\Theta)=\mathfrak P_t(\Theta')\) on every admitted public option coordinate at date \(t\).  The
public-observable state space is the quotient
\[
        \mathfrak X_t^{\mathrm{pub}}
        :=
        \mathfrak X_t/\!\sim_{\mathrm{pub},t},
\]
with canonical projection \(\pi_t^{\mathrm{pub}}:\mathfrak X_t\to\mathfrak X_t^{\mathrm{pub}}\).
\end{definition}

\begin{theorem}[Maximality of the public-observable quotient]\label{thm:public-quotient-maximality}
Let \(Z:\mathfrak X_t\to\cZ\) be any state functional.  The following are equivalent.
\begin{enumerate}[leftmargin=2em]
\item \(Z\) is identified by the public option manifold, in the sense that there is a map
\(\zeta:\mathfrak P_t(\mathfrak X_t)\to\cZ\) such that \(Z=\zeta\circ \mathfrak P_t\).
\item \(Z\) is constant on public-indistinguishability classes.
\item \(Z\) factors through the quotient, in the sense that there is a map
\(\bar Z:\mathfrak X_t^{\mathrm{pub}}\to\cZ\) such that
\[
        Z=\bar Z\circ\pi_t^{\mathrm{pub}}.
\]
\end{enumerate}
Consequently, \(\mathfrak X_t^{\mathrm{pub}}\) is the largest latent object identifiable from
the admitted public option manifold without additional source-family restrictions.
\end{theorem}

\begin{proof}
If \(Z=\zeta\circ \mathfrak P_t\), equal public manifolds imply equal \(Z\)-values, so \(Z\) is constant on
public-indistinguishability classes.  Conversely, if \(Z\) is constant on those classes, define
\[
        \bar Z([\Theta]_{\mathrm{pub}}):=Z(\Theta).
\]
The definition is independent of the representative exactly because the classes are the fibers of
\(\mathfrak P_t\).  It also defines a map on the public image by
\[
        \zeta(\mathfrak P_t(\Theta)):=Z(\Theta),
\]
again well defined on each fiber.  Thus
\(Z=\bar Z\circ\pi_t^{\mathrm{pub}}=\zeta\circ \mathfrak P_t\).  Every identified functional factors
through \(\mathfrak X_t^{\mathrm{pub}}\), which is the maximal public-observable quotient.
\end{proof}

\begin{corollary}[Source-family recovery as quotient embedding]\label{cor:source-quotient-embedding}
Let \(\mathfrak S_t\subset\mathfrak X_t\) be a gauge-normalized Volterra--Perron source family,
and let \(\mathfrak S_t/\!\sim_g\) be its gauge quotient.  Local recovery from the public
manifold is equivalent to local injectivity of
\[
        \mathfrak S_t/\!\sim_g
        \longrightarrow
        \mathfrak X_t^{\mathrm{pub}},
        \qquad
        [\Theta]_g\mapsto[\Theta]_{\mathrm{pub}}.
\]
Thus the source-family rank conditions are exactly the conditions under which the canonical
source quotient embeds into the public-observable quotient.
\end{corollary}

\begin{proof}
By Theorem~\ref{thm:public-quotient-maximality}, the public manifold identifies exactly the public
quotient \(\mathfrak X_t^{\mathrm{pub}}\).  A recovery map on the source family must therefore
invert the induced map from the gauge quotient \(\mathfrak S_t/\!\sim_g\) into this public quotient.
Local recovery is local injectivity of that induced map, and stable local recovery is the same
statement with a positive inverse modulus.  The source-family rank conditions are the differential
certificate for this quotient embedding.
\end{proof}


\begin{definition}[Dynamic public-equivalent density family]\label{def:dynamic-public-equivalent}
At date \(t\), a positive density family
\[
        Z_{t,\cdot}=(Z_{t,S})_{S>t},
        \qquad Z_{t,S}\in L^1(\QQ,\cK_S),
\]
is public-equivalent if, for every \(S>t\),
\[
        \EE^\QQ[Z_{t,S}\mid\cH_t]=1
\]
and, for every admitted public payoff \(X_a\) with maturity \(T(a)\),
\[
        \EE^\QQ[Z_{t,T(a)}X_a\mid\cH_t]
        =
        \EE^\QQ[X_a\mid\cH_t].
\]
It prices a claim \(X\in L^1(\QQ,\cF_T)\) by
\[
        \Pi_t^Z(X):=\EE^\QQ[Z_{t,T}X\mid\cH_t].
\]
\end{definition}

\begin{lemma}[Public preservation]\label{lem:public-preservation}
If \(Z_{t,\cdot}\) is public-equivalent and \(Y\in L^\infty(\QQ,\cH_t)\), then, for every
\(S>t\),
\[
        \EE^\QQ[Z_{t,S}Y\mid\cH_t]=Y.
\]
Consequently the density family preserves every date-\(t\) public payoff test.
\end{lemma}

\begin{proof}
The public-equivalence normalization is \(\EE^\QQ[Z_{t,S}\mid\cH_t]=1\).  Since
\(Y\) is \(\cH_t\)-measurable,
\[
        \EE^\QQ[Z_{t,S}Y\mid\cH_t]
        =
        Y\,\EE^\QQ[Z_{t,S}\mid\cH_t]
        =
        Y.
\]
Multiplication by a date-\(t\) public test therefore does not leave the public equivalence class.
\end{proof}

\begin{lemma}[Residual annihilators generate public-equivalent density families]\label{lem:residual-annihilator}
Let \(R\in L^\infty(\QQ,\cK_t)\) satisfy
\[
        \EE^\QQ[R\mid\cH_t]=0,
        \qquad
        \EE^\QQ[RX_a\mid\cH_t]=0
        \quad\text{for every admitted public payoff }X_a .
\]
If \(|\eta|<\norm{R}_\infty^{-1}\), then the density family
\[
        Z_{t,S}^{\eta}=1+\eta R,\qquad S>t,
\]
is positive and public-equivalent.
\end{lemma}

\begin{proof}
The bound \(|\eta|<\norm{R}_\infty^{-1}\) gives
\(1+\eta R\ge 1-|\eta|\norm{R}_\infty>0\) almost surely.  The normalization condition gives
\[
        \EE^\QQ[Z_{t,S}^{\eta}\mid\cH_t]
        =
        1+\eta\EE^\QQ[R\mid\cH_t]
        =
        1.
\]
For every admitted payoff \(X_a\), the annihilator condition gives
\[
        \EE^\QQ[Z_{t,T(a)}^{\eta}X_a\mid\cH_t]
        =
        \EE^\QQ[X_a\mid\cH_t]
        +
        \eta\EE^\QQ[RX_a\mid\cH_t]
        =
        \EE^\QQ[X_a\mid\cH_t].
\]
The density family is positive and public-equivalent.
\end{proof}

\begin{theorem}[Public non-identification of latent stress]\label{thm:public-nonidentification}
Fix \(t<T\).  Suppose there is a bounded \(R\in L^\infty(\QQ,\cK_t)\) such that
\[
        \EE^\QQ[R\mid\cH_t]=0
        \quad\text{and}\quad
        \EE^\QQ[RX_a\mid\cH_t]=0
        \quad\text{for every admitted public payoff }X_a,
        \tag{A}
        \label{eq:public-annihilator}
 \]
but
\[
        \EE^\QQ[L_{t,T}^\QQ(X)R\mid\cH_t]\ne0
\]
on a set of positive probability for some latent-stress claim \(X\).  Then, for sufficiently small
nonzero \(\eta\), the density family
\[
        Z_{t,S}^\eta=1+\eta R,\qquad S>t,
\]
preserves the public option manifold at date \(t\) but changes the selected price of \(X\):
\[
        \Pi_t^{Z^\eta}(X)
        =
        \EE^\QQ[X\mid\cH_t]
        +
        \eta\,\EE^\QQ[L_{t,T}^\QQ(X)R\mid\cH_t].
\]
Thus public calibration alone cannot recover the latent stress price.
\end{theorem}

\begin{proof}
Choose \(|\eta|<\norm{R}_\infty^{-1}\).  Lemma~\ref{lem:residual-annihilator} shows that
\(Z_{t,\cdot}^{\eta}\) is public-equivalent and therefore preserves the public option manifold.
For the latent claim,
\[
        \EE^\QQ[Z_{t,T}^\eta X\mid\cH_t]
        =
        \EE^\QQ[X\mid\cH_t]
        +
        \eta\EE^\QQ[XR\mid\cH_t].
\]
Since \(R\in L^\infty(\QQ,\cK_t)\), the tower property gives
\[
        \EE^\QQ[XR\mid\cH_t]
        =
        \EE^\QQ[\EE^\QQ[X\mid\cK_t]R\mid\cH_t]
        =
        \EE^\QQ[L_{t,T}^\QQ(X)R\mid\cH_t].
\]
The last term is nonzero on a set of positive probability by assumption, so the selected price is
moved while all public prices are fixed.
\end{proof}

\begin{corollary}[No universal recovery from public compression]\label{cor:no-universal-recovery}
If the residual space \(L^2(\QQ,\cK_t)\ominus L^2(\QQ,\cH_t)\) contains a bounded direction
correlated with a latent-stress loading, then no procedure using only the public option manifold
can identify all latent-stress prices without additional source-family restrictions.
\end{corollary}

\begin{proof}
Theorem~\ref{thm:public-nonidentification} constructs two positive public-equivalent density
families with the same value of every admitted public payoff and different conditional prices for
the latent claim.  A reconstruction rule whose input is only \(\cM_t^{\mathrm{opt}}\) has the
same input under both families.  It therefore cannot identify both selected latent-stress prices.
\end{proof}

\begin{corollary}[Exact residual obstruction]\label{cor:exact-residual-obstruction}
Fix the admitted public payoff family at date \(t\), and write
\[
        L_a:=\EE^\QQ[X_a\mid\cK_t],
        \qquad
        L_X:=\EE^\QQ[X\mid\cK_t].
\]
Let
\[
        \cP_t^{\mathrm{pub}}
        :=
        \overline{\left\{
        \eta_0+\sum_{j=1}^n\eta_j L_{a_j}:
        n<\infty,\ a_j\in\cI_t,\ \eta_0,\eta_j\in L^\infty(\QQ,\cH_t)
        \right\}}^{\,L^2(\QQ,\cK_t)}
        \subset L^2(\QQ,\cK_t)
\]
be the closed conditional public-loading module generated by the date-\(t\) public payoffs.
Assume bounded elements are dense in \((\cP_t^{\mathrm{pub}})^\perp\).  A latent claim \(X\)
has a public-equivalent residual price degree of freedom at date \(t\) if and only if
\[
        (I-C_{\cP_t^{\mathrm{pub}}}^\QQ)L_X\ne0.
\]
In that case there is a bounded residual direction \(R\) preserving all admitted public prices
and changing the selected price of \(X\).
\end{corollary}

\begin{proof}
Orthogonality to the conditional module is equivalent to the conditional annihilator equations:
\[
        R\in(\cP_t^{\mathrm{pub}})^\perp
        \quad\Longleftrightarrow\quad
        \EE^\QQ[R\mid\cH_t]=0,
        \qquad
        \EE^\QQ[RL_a\mid\cH_t]=0\quad\text{for every }a\in\cI_t .
\]
The equivalence follows by testing \(R\) against
\(Y(\eta_0+\sum_j\eta_jL_{a_j})\) with bounded \(\cH_t\)-measurable coefficients and using the
defining property of conditional expectation.  Since \(R\) is \(\cK_t\)-measurable,
\(\EE^\QQ[RX_a\mid\cH_t]=\EE^\QQ[RL_a\mid\cH_t]\), so these are exactly the public-equivalence
annihilator equations.

If \((I-C_{\cP_t^{\mathrm{pub}}}^\QQ)L_X=0\), then \(L_X\) belongs to the conditional
public-loading module, and every public-annihilating residual direction has
\(\EE^\QQ[L_XR\mid\cH_t]=0\).  Public-equivalent tilts cannot change the price of \(X\).

If the projection is nonzero, set
\[
        R_0=(I-C_{\cP_t^{\mathrm{pub}}}^\QQ)L_X .
\]
Then \(R_0\in(\cP_t^{\mathrm{pub}})^\perp\) and
\(\EE^\QQ[L_XR_0\mid\cH_t]=\EE^\QQ[R_0^2\mid\cH_t]\) is nonzero on the event where the residual
projection is nonzero.  By the density assumption, choose a bounded residual direction
\(R\in(\cP_t^{\mathrm{pub}})^\perp\) sufficiently close to \(R_0\) in conditional \(L^2\) to
preserve a nonzero conditional pairing with \(L_X\).  Lemma~\ref{lem:residual-annihilator} then
generates positive public-equivalent density families that leave the public manifold fixed and
move the price of \(X\).
\end{proof}

\begin{definition}[Public payoff-module completion]\label{def:public-completion}
For a fixed horizon \(T\), define the closed conditional public payoff module
\[
        \cA_{t,T}^{\mathrm{pub}}
        :=
        \overline{\Span\{Y\,L_{t,T(a)}^\QQ(X_a):
        Y\in L^\infty(\QQ,\cH_t),\ a\in\cI_t,\ T(a)=T\}}^{\,L^2(\QQ,\cK_t)}
        \subset L^2(\QQ,\cK_t).
\]
The admitted public manifold is complete for the date-\(t\) latent loading class
\(\cL_{t,T}\subset L^2(\QQ,\cK_t)\) if
\[
        \cL_{t,T}\subseteq \cA_{t,T}^{\mathrm{pub}}.
\]
\end{definition}

\begin{theorem}[Public completion and residual collapse]\label{thm:public-completion}
Fix \(t<T\), and let \(\cL_{t,T}\subset L^2(\QQ,\cK_t)\) be the latent loading class whose
claims are to be priced.  The following are equivalent.
\begin{enumerate}[leftmargin=2em]
\item The public manifold is complete for \(\cL_{t,T}\).
\item Every \(L\in\cL_{t,T}\) has zero orthogonal residual relative to
\(\cA_{t,T}^{\mathrm{pub}}\).
\item No public-equivalent density perturbation generated by a bounded public-annihilating
residual direction can change the price of any claim with loading in \(\cL_{t,T}\).
\end{enumerate}
When these conditions hold, residual price intervals collapse to singletons on
\(\cL_{t,T}\).  If they fail and bounded annihilating residual directions are dense in the
orthogonal complement, every nonzero loading residual generates a genuine residual pricing fiber.
\end{theorem}

\begin{proof}
The first two statements are the Hilbert projection decomposition of each loading
\(L\in\cL_{t,T}\) into its component in \(\cA_{t,T}^{\mathrm{pub}}\) and its orthogonal
residual.  If every residual component is zero, then every public-annihilating direction
\(R\in(\cA_{t,T}^{\mathrm{pub}})^\perp\) has zero conditional pairing with every such loading.
Lemma~\ref{lem:residual-annihilator} therefore cannot move the selected price of any claim in the
class.

Conversely, if some loading has a nonzero orthogonal residual, the density assumption supplies a
bounded annihilating direction with nonzero conditional pairing against that residual.  The
positive densities generated by Lemma~\ref{lem:residual-annihilator} preserve all admitted public
prices and move the claim price.  The image of this one-dimensional tilt is a nontrivial residual
pricing fiber, and it collapses exactly when the projection residual is zero.
\end{proof}

\begin{corollary}[Selector redundancy at public completion]\label{cor:selector-redundancy}
On a publicly complete latent loading class, the residual selector \(\theta_t\) is redundant.
Outside public completion, the selector chooses a point in the unspanned pricing fiber left by the
public payoff module.
\end{corollary}

\begin{proof}
Theorem~\ref{thm:public-completion} identifies public completion with vanishing residual
projection for every loading in the claim class.  In that case all public-annihilating density
tilts have zero conditional pairing with those loadings, so the selected residual term is fixed and
the selector carries no pricing degree of freedom.  If completion fails, the orthogonal residual
projection is nonzero for some loading; Lemma~\ref{lem:residual-annihilator} then gives
public-equivalent tilts that move the corresponding price.  The selector chooses one element of
that residual fiber.
\end{proof}


\begin{theorem}[Option-manifold impossibility theorem]\label{thm:option-manifold-impossibility}
Fix date \(t\) and an arbitrage-free admitted public option family.  The following obstructions
are sharp.
\begin{enumerate}[leftmargin=2em]
\item Without source-family restrictions, the public manifold identifies only
\(\mathfrak X_t^{\mathrm{pub}}\).  No calibration or estimator using only the admitted public
coordinates can distinguish two latent states in the same public equivalence class.
\item Single-name vanilla surfaces identify marginal risk-neutral laws, not joint covariance
routes, directed Volterra contagion, or Perron synchronization.  Distinct couplings of the same
marginals may have the same single-name surfaces and different index or latent-stress prices.
\item If the finite basket or carrier design has a nontrivial null direction on the active
covariance-route subspace, then option prices on that finite panel cannot separate perturbations in
that null direction while they remain inside the covariance cone.
\item If the Perron reduced-resolvent norm diverges, in particular along gap-closing seams, the
Perron projector is not uniformly Lipschitz as a function of the public response; any recovery map
for \(e^\QQ\), \(\ell^\QQ\), or \(P^\QQ\) has an inverse modulus that diverges along such strata.
\item If the maturity dictionary has a nontrivial null combination on the observed maturities,
distinct Volterra source geometries produce the same maturity panel.
\item If the public payoff module is incomplete for a claim class, public-equivalent residual
density families preserve every admitted public option price while assigning different prices to
some latent-stress claims in that class.
\end{enumerate}
Thus every positive recovery statement requires the corresponding rank, gap,
source-family, or public-completion hypothesis.  The failures are mathematical obstructions in the
option manifold, not deficiencies of a particular numerical procedure.
\end{theorem}

\begin{proof}
The first statement is Theorem~\ref{thm:public-quotient-maximality}.  For the second, fixed
single-name marginal laws do not determine a joint coupling; distinct couplings with the same
marginals can change cross-covariances and hence index or latent-stress prices.  The third is the
nullspace argument used in the covariance-recovery and elliptope results of
Theorems~\ref{thm:quadratic-basket-covariance} and
\ref{thm:elliptope-perron-uncertainty}.  The fourth follows from the reduced-resolvent formula for
a simple isolated eigenvalue, because the derivative of the Perron projector is controlled by that
resolvent.  The fifth is rank failure of the finite maturity response matrix.  The sixth is
Theorem~\ref{thm:public-completion} together with Lemma~\ref{lem:residual-annihilator}.
Each obstruction therefore leaves either a kernel in the public response, an unbounded inverse
modulus, or a nontrivial residual pricing fiber.
\end{proof}

\section{Risk-neutral observability of the option manifold}\label{sec:observability}

Recovery requires more than public calibration.  The option manifold must separate canonical
coordinates after gauge directions have been removed.  This section gives the infinitesimal and
finite-panel verification for Theorem~\ref{thm:spine}.  Observability is the derivative-level form
of quotient embedding; the smallest quotient singular value is the local conditioning scale for
source-chart recovery from public option data.

\begin{definition}[Risk-neutral observability operator]\label{def:observability-operator}
Let \(\psi:U\subset E\to\mathfrak X_t^{\mathrm{str}}\) be a gauge-normalized source
chart and let
\[
        F=\mathcal R_t\circ\psi:U\to\cY_t
\]
be the observable option-response map.  The risk-neutral observability operator at \(u_0\) is
\[
        \mathcal O_t^\QQ:=DF(u_0):E\to\cY_t.
\]
Its canonical version is the induced quotient operator
\[
        \overline{\mathcal O}_t^\QQ:\bar E=E/E_g\to\Ran DF(u_0),
        \qquad
        \overline{\mathcal O}_t^\QQ [h]=DF(u_0)h .
\]
The state is locally observable at \(u_0\) if
\(\overline{\mathcal O}_t^\QQ\) has a bounded left inverse on its range.
\end{definition}

\begin{definition}[Single-name and joint option manifolds]\label{def:single-joint-manifolds}
The single-name vanilla manifold is the restriction of \(\cM_t^{\mathrm{opt}}\) to payoffs
depending on one underlier at a time.  The joint index/carrier manifold adds index, basket,
variance, or carrier-spread option coordinates whose payoffs depend on at least two underliers
through the same maturity window.
\end{definition}

\begin{theorem}[Vanilla insufficiency for joint stress]\label{thm:vanilla-insufficiency}
Suppose two risk-neutral Volterra--Perron states have the same conditional marginal law for each
single underlier over all maturities in the admitted vanilla family, but have different conditional
cross-variation response
\[
        \Cov^\QQ\!\left(
        \int_t^{t+\tau}V_i(s)\dd s,
        \int_t^{t+\tau}V_j(s)\dd s
        \,\middle|\,\cH_t
        \right)
\]
for some \(i\ne j\) and some maturity \(\tau\).  Then their single-name vanilla manifolds agree,
but a nondegenerate index or basket variance response differs.  Consequently, single-name
vanilla surfaces alone do not recover the directed Volterra contagion operator, the Perron
synchronization state, or synchronized latent stress.
\end{theorem}

\begin{proof}
A single-name vanilla payoff is a Borel function of one terminal underlier.  Its date-\(t\) price
is determined by the corresponding conditional marginal law.  Equal conditional marginal laws
therefore imply identical single-name vanilla surfaces.  An index or basket variance claim
contains cross terms of the form
\[
        2\omega_i\omega_j
        \Cov^\QQ\!\left(
        \int_t^{t+\tau}V_i(s)\dd s,
        \int_t^{t+\tau}V_j(s)\dd s
        \,\middle|\,\cH_t
        \right),
\]
for nonzero weights \(\omega_i\omega_j\).  If this cross response differs and the basket
denominator is nondegenerate, the joint option response differs.  The directed Volterra and
Perron synchronization coordinates enter exactly through these joint responses, so they are not
identified by marginal vanilla prices alone.
\end{proof}

\begin{proposition}[Basket design observability]\label{prop:basket-design-observability}
Let \(\cE\) be the active set of unordered cross-pairs \((i,j)\) used in recovery.  For each
admitted basket or index \(u\), define the cross-variance response
\[
        B_u(\tau)
        :=
        \mathsf{TV}_{u,t}(\tau)-\mathsf{TV}_{u,t}^0(\tau)
        -
        \sum_i (w_i^u)^2V_{ii,t}^\QQ(\tau)
        =
        2\sum_{(i,j)\in\cE}w_i^u w_j^uV_{ij,t}^\QQ(\tau).
\]
Let the basket design matrix be
\[
        W_{u,(i,j)}:=2w_i^u w_j^u,\qquad (i,j)\in\cE .
\]
Then, for each maturity \(\tau\), the active cross-route vector
\[
        V_{\cE}(\tau)=\bigl(V_{ij,t}^\QQ(\tau)\bigr)_{(i,j)\in\cE}
\]
is identified from the joint basket responses \(B(\tau)\) if and only if
\[
        \operatorname{rank} W_{\cE}=|\cE|.
\]
When this rank condition holds,
\[
        \norm{\delta V_{\cE}(\tau)}
        \le
        \sigma_{\min}(W_{\cE})^{-1}\norm{\delta B(\tau)} .
\]
\end{proposition}

\begin{proof}
The displayed variance decomposition gives the linear system
\[
        B(\tau)=W_{\cE}V_{\cE}(\tau).
\]
The vector \(V_{\cE}(\tau)\) is identifiable exactly when this linear map is injective on the
active cross-route coordinates, which is equivalent to full column rank.  If \(W_{\cE}\) has full column rank, its Moore--Penrose left inverse has norm
\(\sigma_{\min}(W_{\cE})^{-1}\), which gives the displayed stability estimate.
\end{proof}

\begin{definition}[Dual recovery certificate]\label{def:dual-recovery-certificate}
Let \(F:U\subset\R^d\to\R^m\) be a finite quote response map with full-column-rank Jacobian
\[
        J:=DF(u_0).
\]
A dual recovery certificate for latent coordinate \(k\) is a vector \(a_k\in\R^m\) satisfying
\[
        a_k^\top J=e_k^\top,
\]
where \(e_k\) is the \(k\)th coordinate vector in \(\R^d\).  A matrix \(A\in\R^{d\times m}\)
with \(AJ=I_d\) is a full dual certificate.
\end{definition}

\begin{theorem}[Quote portfolios recover infinitesimal coordinates]\label{thm:dual-certificates}
Let \(F:U\subset\R^d\to\R^m\) be \(C^2\), and suppose \(J=DF(u_0)\) has full column rank.  If
\(A\) is any left inverse of \(J\), then, for \(u\) near \(u_0\),
\[
        A\bigl(F(u)-F(u_0)\bigr)
        =
        u-u_0+O(\norm{u-u_0}^2).
\]
Equivalently,
\[
        u_k-u_{0,k}
        =
        a_k^\top\bigl(F(u)-F(u_0)\bigr)
        +
        O(\norm{u-u_0}^2).
\]
If quote noise has covariance matrix \(\Sigma_y\), the linear certificate variance for coordinate
\(k\) is \(a_k^\top\Sigma_y a_k\).  When \(\Sigma_y\) is positive definite, the minimum-variance
linear unbiased certificate is the \(k\)th row of
\[
        A_{\Sigma}
        :=
        (J^\top\Sigma_y^{-1}J)^{-1}J^\top\Sigma_y^{-1}.
\]
\end{theorem}

\begin{proof}
Taylor expansion gives
\(F(u)-F(u_0)=J(u-u_0)+O(\norm{u-u_0}^2)\).  Multiplication by \(A\) and the identity
\(AJ=I_d\) give the coordinate recovery expansion.  A certificate \(a_k\) applies the linear quote
functional \(a_k^\top\delta y\), so its noise variance is \(a_k^\top\Sigma_y a_k\).  Minimizing
this variance subject to the unbiasedness constraint \(AJ=I_d\) gives the normal equations
\(A\Sigma_y^{-1}J=I_d\), and therefore the generalized least-squares left inverse
\((J^\top\Sigma_y^{-1}J)^{-1}J^\top\Sigma_y^{-1}\).
\end{proof}

\begin{corollary}[Functional certificate in infinite quote spaces]\label{cor:functional-certificate}
Let \(DF(u_0):E\to Y\) have a bounded left inverse \(L:Y\to E\).  For every continuous coordinate
functional \(\ell\in E^\ast\), the composition
\[
        \alpha_\ell:=\ell\circ L\in Y^\ast
\]
satisfies
\[
        \alpha_\ell(DF(u_0)h)=\ell(h).
\]
Thus every recovered infinitesimal latent coordinate has a dual quote functional that extracts it
from the first-order option-manifold response.
\end{corollary}

\begin{proof}
The identity is the definition of a left inverse.  If \(LDF(u_0)=I_E\), then for every
\(h\in E\),
\[
        \alpha_\ell(DF(u_0)h)
        =
        \ell(LDF(u_0)h)
        =
        \ell(h).
\]
Continuity follows from \(\ell\in E^\ast\) and boundedness of \(L\).
\end{proof}


\begin{corollary}[One index identifies one synchronization aggregate]\label{cor:one-index-aggregate}
If the active unordered edge set satisfies \(|\cE|>1\) and the only joint coordinate is one
index or basket \(u\), then \(W_{\cE}\) has one row.  It cannot identify the full edge vector
\(V_{\cE}(\tau)\).  It identifies only the aggregate
\[
        2\sum_{(i,j)\in\cE}w_i^u w_j^uV_{ij,t}^\QQ(\tau).
\]
Full cross-route recovery therefore requires additional carrier baskets, sector or index
panels, a smaller active edge set, or source-family restrictions that reduce the effective
dimension.
\end{corollary}

\begin{proof}
With one basket, \(W_{\cE}\) has a single row, so its rank is at most one.  When
\(|\cE|>1\), the map from the active edge vector to public basket variance has a nontrivial kernel.
The only identified cross term is the displayed row action of \(W_{\cE}\) on
\(V_{\cE}(\tau)\).  Separating the remaining edge coordinates requires more independent rows or a
source restriction that lowers the quotient dimension.
\end{proof}


\begin{corollary}[Generic basket designs]\label{cor:generic-basket-designs}
Suppose \(r\ge|\cE|\) basket weight vectors are drawn from a distribution with density on an
open subset of \(\R^N\), and suppose the design class contains one choice of weights for which
\(W_{\cE}\) has rank \(|\cE|\).  Then
\[
        \operatorname{rank}W_{\cE}=|\cE|
\]
with probability one.
\end{corollary}

\begin{proof}
Some \(|\cE|\times|\cE|\) minor determinant is then a nonzero polynomial in the basket weights.
The zero set of a nonzero polynomial has Lebesgue measure zero, and a distribution with density
assigns that set probability zero.
\end{proof}

\begin{definition}[Variance-strip covariance matrix]\label{def:variance-strip-covariance}
Fix a maturity window \(\tau\), and let \(X_i(\tau)\) be the centered integrated return,
variance, or volatility factor whose basket variance is observed.  The variance-strip covariance
matrix is
\[
        C(\tau)=(C_{ij}(\tau))_{i,j=1}^N,
        \qquad
        C_{ij}(\tau)=\Cov^\QQ(X_i(\tau),X_j(\tau)\mid\cH_t).
\]
For a basket weight \(w\in\R^N\), the model-free quadratic basket response is
\[
        V_w(\tau)=w^\top C(\tau)w.
\]
\end{definition}

\begin{theorem}[Quadratic basket manifold and covariance recovery]\label{thm:quadratic-basket-covariance}
For a fixed \(\tau\), knowledge of \(V_w(\tau)\) for all \(w\in\R^N\) is equivalent to knowledge
of the full symmetric covariance matrix \(C(\tau)\).  In particular,
\[
        C_{ij}(\tau)
        =
        \frac12\left(
        V_{e_i+e_j}(\tau)-V_{e_i}(\tau)-V_{e_j}(\tau)
        \right).
\]
A finite family of basket weights \(w_1,\dots,w_m\) identifies \(C(\tau)\) on a target subspace
\(\mathcal V\subseteq\mathbb S^N\) if and only if the linear functionals
\[
        C\mapsto \operatorname{Tr}(Cw_aw_a^\top),\qquad a=1,\dots,m,
\]
separate points of \(\mathcal V\).
\end{theorem}

\begin{proof}
The map \(C\mapsto w^\top Cw=\operatorname{Tr}(Cww^\top)\) is linear in the symmetric matrix
\(C\).  Polarization gives the displayed identity.  The finite-weight statement is exactly
injectivity of this linear observation map on \(\mathcal V\).
\end{proof}

\begin{theorem}[Elliptope feasibility and Perron uncertainty]\label{thm:elliptope-perron-uncertainty}
Suppose the single-name variance levels \(C_{ii}(\tau)>0\) are known, and set
\[
        R_{ij}(\tau)
        =
        \frac{C_{ij}(\tau)}{\sqrt{C_{ii}(\tau)C_{jj}(\tau)}}.
\]
A candidate cross-covariance matrix is statically feasible only if
\[
        R(\tau)\succeq0,\qquad R_{ii}(\tau)=1.
\]
If only a finite basket family is observed, the model-free feasible set is
\[
        \mathcal C_\tau
        :=
        \{C\succeq0:\operatorname{diag}(C)=v,\ 
        \operatorname{Tr}(Cw_aw_a^\top)=V_{w_a},\ a=1,\dots,m\}.
\]
For any screened Perron matrix \(A(C)\) constructed from \(C\), the model-free spectral-radius
uncertainty interval is
\[
        [\rho^-_\tau,\rho^+_\tau]
        :=
        \left[
        \inf_{C\in\mathcal C_\tau}\rho(A(C)),
        \sup_{C\in\mathcal C_\tau}\rho(A(C))
        \right].
\]
A Perron phase classification from the finite basket panel is model-free certified only when
this interval lies entirely on one side of the phase threshold.
\end{theorem}

\begin{proof}
Positive semidefiniteness and unit diagonal are necessary for any correlation matrix and
sufficient for a centered Gaussian vector with that correlation matrix.  The finite basket
constraints are affine in \(C\), so the feasible set is an affine slice of the positive
semidefinite cone.  If the panel is feasible, the constraints \(\operatorname{diag}(C)=v\) with
positive variance levels make this slice closed and bounded, hence compact.  The map
\(C\mapsto\rho(A(C))\) is continuous for the screened Perron matrices considered here, so the
infimum and supremum are attained on \(\mathcal C_\tau\).  If the displayed spectral interval
crosses a phase threshold, two statically feasible covariance completions imply different phase
labels.
\end{proof}

\begin{definition}[Direction-resolving response block]\label{def:direction-resolving-block}
Let \(G_{i\leftarrow j}\) denote the ordered Volterra route from source \(j\) to target \(i\),
and let \(\vec{\cE}\) be an active ordered edge set.  A public observation block is
direction-resolving on \(\vec{\cE}\) if the derivative of its response map with respect to the
ordered route coordinates,
\[
        J_{\mathrm{dir}}:
        (\delta G_{i\leftarrow j})_{(i,j)\in\vec{\cE}}
        \longmapsto
        \delta\cM_t^{\mathrm{opt}},
\]
is injective after quotienting by driver-basis gauge.
\end{definition}

\begin{theorem}[Symmetric covariance does not identify arbitrary direction]
\label{thm:symmetric-covariance-direction}
Suppose the admitted public response depends on ordered routes only through symmetric
combinations
\[
        S_{ij}
        =
        \Psi_{ij}(G_{i\leftarrow j},G_{j\leftarrow i})
        =
        \Psi_{ij}(G_{j\leftarrow i},G_{i\leftarrow j})
\]
after the remaining source coordinates are fixed.  Then any tangent direction satisfying
\[
        D\Psi_{ij}[\delta G_{i\leftarrow j},\delta G_{j\leftarrow i}]=0
\]
for every active unordered pair lies in the kernel of the public observability operator.  In
particular, if locally \(S_{ij}=G_{i\leftarrow j}+G_{j\leftarrow i}\), the anti-symmetric
perturbation
\[
        \delta G_{i\leftarrow j}=H,\qquad
        \delta G_{j\leftarrow i}=-H
\]
is invisible to a symmetric covariance manifold.
\end{theorem}

\begin{proof}
If the public response factors through the symmetric map \(S=(S_{ij})\), then its derivative
factors as \(D\bar F\circ DS\).  Any tangent vector in \(\Ker DS\) is therefore in the kernel of
the public response derivative.  In the additive example, \(DS[H,-H]=H-H=0\), so the anti-symmetric perturbation lies in that
kernel.
\end{proof}

\begin{theorem}[Sufficient orientation separation]\label{thm:orientation-separation}
Suppose that, after gauge normalization and after imposing all symmetric covariance
constraints, each ordered route response on the active set admits the analytic representation
\[
        R_{i\leftarrow j}(\tau)
        =
        \sum_{q=1}^{d_{ij}}\gamma_{ijq}f_{ijq}(\tau),
        \qquad \tau\in I,
\]
where the family \(\{f_{ijq}:(i,j,q)\in\vec{\cE}\}\) is linearly independent on the admitted
maturity interval \(I\).  Then an open maturity interval identifies the ordered route
coefficients \(\gamma_{ijq}\).  A generic finite analytic maturity panel with sufficiently many
points identifies the same coefficients.
\end{theorem}

\begin{proof}
The derivative of the response with respect to \(\gamma_{ijq}\) is the analytic function
\(f_{ijq}\).  If a quotient tangent vector is annihilated on the open interval \(I\), the
corresponding linear combination of the functions \(f_{ijq}\) vanishes on \(I\).  Linear
independence forces all tangent coefficients to be zero, so the continuum derivative is injective.
The finite-panel conclusion follows from Corollary~\ref{cor:generic-finite-panel}.
\end{proof}

\begin{definition}[Conditional basket cumulant tensor]\label{def:basket-cumulant-tensor}
For a return vector \(\Delta_\tau X:=X_{t+\tau}-X_t\), suppose the conditional moment generating
function is finite near the origin.  Define
\[
        \kappa_t(z,\tau)
        :=
        \log\EE^\QQ[e^{z^\top\Delta_\tau X}\mid\cH_t].
\]
The order-\(n\) conditional cumulant tensor is
\[
        K_t^{(n)}(\tau)_{i_1\cdots i_n}
        :=
        \partial_{z_{i_1}}\cdots\partial_{z_{i_n}}\kappa_t(0,\tau).
\]
A full signed basket smile for \(z^\top\Delta_\tau X\) determines \(\kappa_t(z,\tau)\) near
zero and hence all finite cumulant tensors whose moments exist.
\end{definition}

\begin{proposition}[Skew-tensor orientation test]\label{prop:skew-tensor-orientation}
In a Volterra leverage source family, suppose an ordered edge \(j\to i\) contributes to the
short-maturity third cumulant through
\[
        K_t^{(3)}(\tau)_{iij}
        =
        A_{i\leftarrow j}\tau^{\alpha_{i\leftarrow j}}
        +
        o(\tau^{\alpha_{i\leftarrow j}}),
        \qquad A_{i\leftarrow j}\ne0,
\]
while the reverse edge contributes to
\[
        K_t^{(3)}(\tau)_{jji}
        =
        A_{j\leftarrow i}\tau^{\alpha_{j\leftarrow i}}
        +
        o(\tau^{\alpha_{j\leftarrow i}}).
\]
If the pairs \((\alpha_{i\leftarrow j},A_{i\leftarrow j})\) and
\((\alpha_{j\leftarrow i},A_{j\leftarrow i})\) are separated by the source-family gauge, then
the observed short-maturity skew tensor identifies the orientation of the edge.  If the third
cumulant vanishes by symmetry, the same statement holds with the first nonzero mixed cumulant
tensor.
\end{proposition}

\begin{proof}
The short-maturity power and leading coefficient are recovered from the cumulant response by the
same leading-jet argument as Corollary~\ref{cor:rough-log-derivative}.  Separation of the two
ordered leading pairs assigns the observed mixed cumulant to one orientation rather than the
other.  If the third cumulant is zero, the first nonzero mixed cumulant supplies the leading jet.
\end{proof}

\begin{theorem}[Minimal option-manifold hierarchy]\label{thm:minimal-manifold-hierarchy}
In a gauge-normalized source family with block tangent decomposition
\[
        \bar E^{\mathrm{str}}=E_{\mathfrak H}\oplus E_G\oplus E_{\mathrm{syn}},
\]
the following hierarchy holds locally.
\begin{enumerate}[leftmargin=2em]
\item Single-name vanilla term structures can identify marginal maturity-scaling directions in
\(E_{\mathfrak H}\), but leave \(E_G\oplus E_{\mathrm{syn}}\) unobserved unless cross-kernel coordinates are
fixed by the source family.
\item Adding nondegenerate index or basket option responses creates a synchronization observation
block for \(E_{\mathrm{syn}}\).
\item Adding a nondegenerate carrier span with Gram lower bound
\(\kappa_{\mathrm{car}}>0\) separates carrier-replicated single-name variation from the
synchronized index component.
\item An open maturity interval, or a finite analytic maturity panel with full rank, upgrades
synchronization curves into recovery of rough maturity and Volterra source coordinates.
\item Residual-selector coordinates lie outside the purely public source tangent space.  They
remain unobserved unless residual price instruments are included or a selector restriction is
imposed.
\end{enumerate}
Thus full canonical recovery requires a joint index/single-name/carrier/maturity manifold beyond
calibrated marginal surfaces.
\end{theorem}

\begin{proof}
Each statement follows by restricting the observability operator
\(\overline{\mathcal O}_t^\QQ\) to the indicated public subsystem.  Marginal vanilla payoffs
differentiate only marginal law and maturity-scaling directions, so the joint blocks remain in
the kernel.  Index and basket responses contain cross terms and therefore generate a
synchronization block.  Carrier Gram nondegeneracy makes the carrier projection invertible on the
admitted carrier coordinates, separating replicated single-name variation from common
synchronization.  The finite-panel assertion is the analytic rank argument later formalized in
Theorem~\ref{thm:finite-panel-recovery}.  Finally, residual selectors change prices only along
public-equivalent residual directions; absent residual instruments or a selector restriction,
those directions remain outside the public observation map.
\end{proof}

\begin{definition}[Finite-quote observation error]\label{def:finite-quote-error}
A noisy finite quote panel is
\[
        \widehat y_t=y_t+\eps_t,
        \qquad y_t=\cM_t^{\mathrm{opt}},
\]
where \(y_t\) lies in a recovery neighborhood and \(\eps_t\in\cY_t\) is a quote, smoothing, or
finite-panel perturbation measured in the quote-space norm.
\end{definition}

\begin{theorem}[Finite-quote stability and observability breakdown]\label{thm:finite-quote-stability}
Suppose \(y_t\) and \(\widehat y_t\) remain in the same local recovery neighborhood and
\(\widehat y_t\in W_t\cap F_t(V_t)\), so that the local inverse is defined.  Let
\[
        \widehat\Xi_t^\QQ=\mathcal I_t(\widehat y_t),
        \qquad
        \Xi_t^\QQ=\mathcal I_t(y_t).
\]
Then
\[
        \dist_{\overline{\mathfrak X}_t}
        \!\left(\widehat\Xi_t^\QQ,\Xi_t^\QQ\right)
        \le
        2\norm{L_t}\,\norm{\eps_t}_{\cY_t}.
\]
In finite-dimensional panels, if
\[
        \beta_t
        :=
        \sigma_{\min}\!\left(\overline{\mathcal O}_t^\QQ\right)>0,
\]
then \(\norm{L_t}\le \beta_t^{-1}\).  Hence recovery becomes ill-conditioned when
\(\beta_t\downarrow0\).  In the Volterra--Perron chart this can occur through collapse of the
Perron gap or reduced-resolvent conditioning, carrier Gram lower bound, maturity-response rank,
directed-route rank, synchronization rank, or, in an augmented residual-instrument manifold,
residual-selector rank.
\end{theorem}

\begin{proof}
The first estimate is the local inverse estimate from
Theorem~\ref{thm:conditional-recovery} applied to \(y_t\) and \(y_t+\eps_t\).  In a
finite-dimensional panel, the Moore--Penrose left inverse of the quotient observability matrix has
operator norm \(\sigma_{\min}(\overline{\mathcal O}_t^\QQ)^{-1}\) on the observed range.  The block
decomposition then identifies the ways in which this singular value can vanish.  Rank loss in a
diagonal block, divergence of the Perron reduced-resolvent constant, collapse of the carrier Gram
bound, or loss of a residual-selector rank block removes a bounded left inverse and forces the
local recovery constant to diverge.
\end{proof}

\begin{remark}
If a noisy quote vector does not lie on the model image, the inverse \(\mathcal I_t\) is not
defined.  In that case the regularized finite-panel estimator of
Definition~\ref{def:regularized-panel-reconstruction} is the appropriate reconstruction object.
\end{remark}

\begin{definition}[Bid--ask quote set]\label{def:bid-ask-quote-set}
For a finite panel, let \(\underline y_t,\overline y_t\in\R^M\) be lower and upper executable
quote vectors.  Let \(\cA\) denote the finite set of static-arbitrage linear inequalities used
to filter the panel.  The arbitrage-filtered quote set is
\[
        \cQ_t
        :=
        \{y\in\R^M:\underline y_t\le y\le\overline y_t,\ \cA y\ge0\}.
\]
For a finite-panel response map \(F_\tau:U\to\R^M\), the set of model states compatible with the
quote set is
\[
        \cI_t(\cQ_t):=\{u\in U:F_\tau(u)\in\cQ_t\}.
\]
\end{definition}

\begin{theorem}[Set-valued recovery under bid--ask uncertainty]\label{thm:set-valued-recovery}
Suppose \(F_\tau\) has a Lipschitz inverse on the chosen recovery neighborhood with constant
\(C_\tau\).  Then
\[
        \operatorname{diam}\bigl(\cI_t(\cQ_t)\bigr)
        \le
        C_\tau\,\operatorname{diam}\bigl(\cQ_t\cap F_\tau(U)\bigr).
\]
If the whole bid--ask box has quote-space diameter at most \(\delta_t\), then
\[
        \operatorname{diam}\bigl(\cI_t(\cQ_t)\bigr)
        \le C_\tau\delta_t .
\]
For any selected claim price \(\Pi(X;u)\) that is Lipschitz in the recovered state with constant
\(L_X\),
\[
        \operatorname{diam}\{\Pi(X;u):u\in\cI_t(\cQ_t)\}
        \le L_XC_\tau\delta_t,
\]
plus the residual-envelope width if residual selectors are allowed to vary.
\end{theorem}

\begin{proof}
For any \(u,v\in\cI_t(\cQ_t)\), both \(F_\tau(u)\) and \(F_\tau(v)\) belong to
\(\cQ_t\cap F_\tau(U)\).  The inverse estimate gives
\[
        \norm{u-v}\le C_\tau\norm{F_\tau(u)-F_\tau(v)}.
\]
Taking the supremum over compatible states gives the diameter bound.  If the whole filtered
bid--ask box has diameter at most \(\delta_t\), the same estimate gives the second bound.  Applying
the Lipschitz property of \(\Pi(X;u)\) to the recovered-state set gives the price-diameter bound;
any free selector coordinate adds its residual-envelope width separately.
\end{proof}

\section{Latent-state recovery from rich option manifolds}\label{sec:recovery}

Restricted source families can restore recovery despite global non-identification.  Public
compression alone is insufficient; observability supplies the additional condition.  A resolving
option manifold contains enough cross-sectional and maturity variation to recover the source quotient
inside the specified family.

\begin{assumption}[Gauge-normalized source family]\label{ass:source-chart}
There is a Banach manifold \(\mathfrak X_t^{\mathrm{str}}\) of admissible Volterra--Perron states
and a gauge-normalized local chart
\[
        \psi:U\subset E\to\mathfrak X_t^{\mathrm{str}},
        \qquad
        E=E_{\mathfrak H}\oplus E_G\oplus E_\rho\oplus E_e\oplus E_P.
\]
The local gauge-slice condition in Assumption~\ref{ass:local-gauge-slice} holds.  Thus
gauge-equivalent infinitesimal directions form a closed subspace \(E_g\subset E\), and the
canonical tangent space is identified with the slice \(E_c\simeq E/E_g=:\bar E\).  The
chart represents the coordinates
\[
        (\mathfrak{H}^\QQ,G^\QQ,\rho^\QQ,e^\QQ,\ell^\QQ,P^\QQ)
\]
using the normalizations of Definition~\ref{def:gauge-normalized-chart}.
\end{assumption}

\begin{assumption}[Observable response rank]\label{ass:full-response}
At the reference state \(\Xi_0=\psi(u_0)\), the observable response map
\[
        F:=\mathcal R_t\circ\psi:U\to\cY_t
\]
is continuously Fr\'echet differentiable.  Its derivative annihilates gauge directions and
induces
\[
        \overline{DF}(u_0):\bar E\to \Ran DF(u_0)\subseteq \cY_t .
\]
This induced derivative has a bounded left inverse \(L_0:\Ran DF(u_0)\to\bar E\):
\[
        L_0\,\overline{DF}(u_0)=\Id_{\bar E}.
\]
Equivalently, with \(B:=\norm{L_0}\),
\[
        \norm{h}_{\bar E}
        \le
        B\norm{\overline{DF}(u_0)h}_{\cY_t},
        \qquad h\in\bar E.
\]
\end{assumption}

\begin{assumption}[Rank and conditioning]\label{ass:structural-nondegeneracy}
The reference state satisfies the following lower-bound conditions.
\begin{enumerate}[leftmargin=2em]
\item The screened synchronization operator has a simple Perron eigenvalue with spectral gap
\(\gamma_P>0\) and finite reduced-resolvent constant
\(\kappa_P:=\norm{(A_t^\QQ-\rho_t^\QQ I)_{\mathrm{red}}^{-1}}\) in the chosen Perron chart.
\item The admitted carrier coordinates have Gram matrix \(\Gamma_{\mathrm{car}}\) satisfying
\(\ip{x}{\Gamma_{\mathrm{car}}x}\ge\kappa_{\mathrm{car}}\norm{x}^2\) for some
\(\kappa_{\mathrm{car}}>0\).
\item The rough maturity-response block has full rank either on an open maturity interval or on a
finite analytic maturity panel.
\item The directed Volterra source-family block and the index/single-name synchronization block
have full rank after quotienting by gauge.
\item Any residual selector used for pricing is either fixed by a post-recovery rule or observed
only through an augmented residual-instrument block, not through the pure public option response.
\end{enumerate}
\end{assumption}

\begin{theorem}[Local identification of the canonical Volterra--Perron source quotient]\label{thm:conditional-recovery}
Under Assumptions~\ref{ass:source-chart}--\ref{ass:structural-nondegeneracy}, there are
neighborhoods \(V\subset U\) of \(u_0\) and \(W\subset \cY_t\) of \(F(u_0)\) such that the public
option manifold identifies the canonical source state locally:
\[
        F(u)=F(v),\quad u,v\in V
        \qquad\Longrightarrow\qquad
        \psi(u)\sim_g\psi(v).
\]
Equivalently, the quotient map
\[
        \bar F:V/E_g\to W\cap F(V)
\]
has a continuous local inverse on its image.  After shrinking \(V\), the inverse obeys
a local Lipschitz estimate
\[
        \norm{[u]-[v]}_{\bar E}
        \le
        2\norm{L_0}\,\norm{F(u)-F(v)}_{\cY_t},
        \qquad u,v\in V.
\]
\end{theorem}

\begin{proof}
Because \(DF(u_0)\) annihilates gauge directions and the response is constant on local gauge
orbits, \(F\) descends locally to a differentiable map \(\bar F\) on the quotient chart.  Quotient
representatives are again denoted by \(u,v\).  Set
\[
        J:=D\bar F(u_0).
\]
By the bounded-below form of Assumption~\ref{ass:full-response},
\[
        \norm{h}_{\bar E}\le B\norm{Jh}_{\cY_t},\qquad h\in\bar E.
\]
Continuity of \(D\bar F\) gives a convex neighborhood \(V\) on which
\[
        \norm{D\bar F(z)-J}_{\cL(\bar E,\cY_t)}
        \le \frac{1}{2B}.
\]
For \(u,v\in V\), the fundamental theorem of calculus in the Banach chart gives
\[
        \bar F(u)-\bar F(v)
        =
        J(u-v)
        +
        R(u,v),
        \qquad
        R(u,v):=\int_0^1
        \bigl(D\bar F(v+s(u-v))-D\bar F(u_0)\bigr)(u-v)\dd s .
\]
The neighborhood choice implies
\[
        \norm{R(u,v)}_{\cY_t}\le \frac{1}{2B}\norm{u-v}_{\bar E}.
\]
Therefore
\begin{align*}
        \norm{u-v}_{\bar E}
        &\le
        B\norm{J(u-v)}_{\cY_t}  \\
        &\le
        B\norm{\bar F(u)-\bar F(v)}_{\cY_t}
        +B\norm{R(u,v)}_{\cY_t}  \\
        &\le
        B\norm{\bar F(u)-\bar F(v)}_{\cY_t}
        +\frac12\norm{u-v}_{\bar E}.
\end{align*}
This gives the displayed Lipschitz estimate and hence local injectivity on the quotient.  If
\(F(u)=F(v)\), then \([u]=[v]\), so \(\psi(u)\sim_g\psi(v)\).  The same estimate gives continuity
of the inverse on \(F(V)\).
\end{proof}

\begin{remark}
The theorem is conditional and compatible with the non-identification theorem.  Recovery becomes
possible only after restricting the latent state to a source family whose option responses are
locally complete.  Public-annihilating residual selectors remain unidentified from the same public
quotes.
\end{remark}

\begin{theorem}[Stability of recovered states]\label{thm:recovery-stability}
Suppose Assumptions~\ref{ass:source-chart}--\ref{ass:structural-nondegeneracy} hold and let
\[
        B:=\norm{L_0}<\infty.
\]
Then, after shrinking the recovery neighborhood if necessary, there are constants \(C>0\) and
\(\delta>0\) such that whenever two option manifolds \(y_1,y_2\in W\cap F(V)\) satisfy
\(\norm{y_1-y_2}_{\cY_t}<\delta\), their recovered canonical states satisfy
\[
        \dist_{\bar E}\!\left(\bar F^{-1}(y_1),\bar F^{-1}(y_2)\right)
        \le C\norm{y_1-y_2}_{\cY_t}.
\]
One may take \(C=2B\) on a sufficiently small neighborhood.  In the Volterra--Perron chart,
\(B\) is finite only while the Perron reduced-resolvent constant \(\kappa_P\) remains bounded,
the carrier Gram lower bound \(\kappa_{\mathrm{car}}\) remains bounded away from zero, and the
diagonal response-block singular values remain bounded away from zero.
\end{theorem}

\begin{proof}
The first estimate is the Lipschitz inverse bound proved in
Theorem~\ref{thm:conditional-recovery}.  The Perron contribution is controlled by the reduced
resolvent of the simple eigenvalue; it can blow up as the spectral gap closes or as non-normal
conditioning degenerates.  The dependence on \(\kappa_{\mathrm{car}}\) follows because the carrier
coordinates are recovered by inverting the carrier Gram matrix.  The remaining diagonal response
blocks enter through the norm of the block left inverse.
\end{proof}

\begin{definition}[Observation modulus]\label{def:observation-modulus}
Let \(F:U\subset E\to Y\) be a quotient response map on a recovery chart.  For \(r>0\), define
\[
        \omega_F(r)
        :=
        \inf\{\norm{F(u)-F(v)}_Y:
        u,v\in U,\ \norm{u-v}_E\ge r\}.
\]
For a finite linear panel \(F(u)=F(u_0)+J(u-u_0)\) with full column-rank derivative,
\(\omega_F(r)=r\,\sigma_{\min}(J)\) locally.
\end{definition}

\begin{theorem}[No reconstruction beats the inverse modulus]\label{thm:inverse-modulus-lower-bound}
Suppose only a perturbed quote vector \(\widehat y\) is observed, with deterministic error bound
\[
        \norm{\widehat y-F(u)}_Y\le\varepsilon .
\]
Let \(\widehat u(\widehat y)\) be any reconstruction rule.  If there exist \(u,v\in U\) such that
\[
        \norm{u-v}_E\ge r,
        \qquad
        \norm{F(u)-F(v)}_Y\le2\varepsilon,
\]
then, for at least one admissible observation \(\widehat y\),
\[
        \max\{\norm{\widehat u(\widehat y)-u}_E,\,
        \norm{\widehat u(\widehat y)-v}_E\}
        \ge \frac r2.
\]
Consequently, if \(\omega_F(r)\le2\varepsilon\), the worst-case recovery error at quote-error
level \(\varepsilon\) is at least \(r/2\).
\end{theorem}

\begin{proof}
In a normed linear quote space take \(\widehat y=(F(u)+F(v))/2\).  The bound
\(\norm{F(u)-F(v)}_Y\le2\varepsilon\) makes \(\widehat y\) admissible for both latent states.  The
same observed quote vector is therefore compatible with \(u\) and with \(v\).  If one output
\(\widehat u(\widehat y)\) were within \(r/2\) of both states, the triangle inequality would give
\(\norm{u-v}_E<r\), contradicting the separation condition.  Hence at least one of the two errors
is at least \(r/2\).
\end{proof}

\begin{corollary}[Sharp local recovery scale]\label{cor:sharp-local-recovery-scale}
For a finite panel with least quotient singular value \(\beta>0\), deterministic quote
uncertainty \(\varepsilon\) implies an unavoidable local recovery scale of order
\[
        \frac{\varepsilon}{\beta}.
\]
Thus, as \(\beta\downarrow0\), recovery instability is an information-theoretic limit for every
procedure using the same quote data.
\end{corollary}

\begin{proof}
On the quotient tangent space, the least singular value \(\beta\) gives the smallest first-order
quote displacement per unit latent displacement.  Taking a unit singular direction associated with
\(\beta\), two nearby latent states separated by order \(\varepsilon/\beta\) produce quote vectors
separated by order \(\varepsilon\), up to higher-order terms.  The local minimax lower bound in
Theorem~\ref{thm:inverse-modulus-lower-bound} then gives the same scale, while the inverse estimate
from the bounded left inverse gives the matching upper order.
\end{proof}


\begin{corollary}[Recovery and residual freedom are complementary]\label{cor:recovery-complementarity}
In a locally complete source family, the public option manifold recovers the canonical
Volterra--Perron state.  Outside such a family, public-equivalent residual density families remain
available and the residual price interval of \cite{vidal2026pricing} is the correct pricing
object.
\end{corollary}

\begin{proof}
The first statement is Theorem~\ref{thm:conditional-recovery}.  The second is
Theorem~\ref{thm:public-nonidentification}.  They apply to different information sets, namely a
restricted source family with full response rank and an unrestricted residual kernel class.
\end{proof}

\subsection{Block recovery and finite maturity panels}

The quotient inverse theorem separates recovery into response blocks.  In the option manifold,
these blocks have financial content.  Marginal single-name term structures carry rough maturity
geometry; index-versus-carrier spreads carry synchronization.  A residual charge becomes a
recoverable coordinate only in an augmented manifold containing residual instruments.

\begin{definition}[Observable block response]\label{def:observable-block-response}
The observable response map is block-adapted if, after a local linear recombination of public
quote coordinates, its response space decomposes into
\[
        \cY_t
        =
        \cY_H\oplus\cY_G\oplus\cY_{\mathrm{syn}} .
\]
The blocks are interpreted as maturity-scaling responses, directed cross-kernel responses, and
index-versus-carrier synchronization responses.  An augmented residual-instrument manifold may
append an additional block \(\cY_\theta\).
\end{definition}

\begin{definition}[Block source chart]\label{def:block-source-chart}
A source chart is block triangular if its quotient tangent space decomposes as
\[
        \bar E
        =
        E_{\mathfrak H}\oplus E_G\oplus E_{\mathrm{syn}},
        \qquad
        E_{\mathrm{syn}}:=E_\rho\oplus E_e\oplus E_P,
\]
the observable response is block-adapted,
so that the derivative of the pricing map has the block form
\[
        \overline{DF}(u_0)
        =
        \begin{pmatrix}
        J_H & 0 & 0\\
        * & J_G & 0\\
        * & * & J_{\mathrm{syn}}
        \end{pmatrix}.
\]
The blocks correspond respectively to rough maturity scaling, directed Volterra/cross-kernel
geometry, and Perron synchronization.
\end{definition}

\begin{theorem}[Block-rank observability criterion]\label{thm:block-rank}
Suppose the source chart is block triangular and each diagonal block
\(J_H,J_G,J_{\mathrm{syn}}\) has a bounded left inverse on its range.  Then
\(\overline{DF}(u_0)\) has a bounded left inverse on its range.  Consequently the conditional
recovery theorem applies to the canonical state.  If an augmented residual-instrument block is
present and its diagonal block \(J_\theta\) also has a bounded left inverse, the same triangular
argument recovers the selected residual coordinate as part of the augmented state.
\end{theorem}

\begin{proof}
Write a tangent vector as \(h=(h_H,h_G,h_{\mathrm{syn}})\).  Given the response
\(y=\overline{DF}(u_0)h\), recover \(h_H\) from the first block by applying the left inverse of
\(J_H\).  Subtract the known contribution of \(h_H\) from the second block and recover \(h_G\) by
the left inverse of \(J_G\).  Continue recursively through the synchronization block.  The finite
number of bounded operations gives a bounded left inverse
for the full triangular derivative.  Theorem~\ref{thm:conditional-recovery} then gives local
recovery.
\end{proof}

\begin{theorem}[Triangular block inverse bound]\label{thm:triangular-block-inverse}
Let
\[
        J=
        \begin{pmatrix}
        J_1 & 0 & \cdots & 0\\
        C_{21} & J_2 & \cdots & 0\\
        \vdots & \vdots & \ddots & \vdots\\
        C_{m1} & C_{m2} & \cdots & J_m
        \end{pmatrix}
\]
act between finite products of normed block spaces.  Suppose each diagonal block has a bounded
left inverse \(L_i\) with \(\norm{L_i}\le b_i\), and set \(c_{ij}:=\norm{C_{ij}}\) for
\(j<i\).  Define
\[
        B_1=b_1,\qquad
        B_i=b_i\left(1+\sum_{j<i}c_{ij}B_j\right),\quad i=2,\dots,m.
\]
Then \(J\) has a bounded left inverse on its range, and one may choose it so that
\[
        \norm{J_{\mathrm{left}}^{-}}
        \le
        \left(\sum_{i=1}^mB_i^2\right)^{1/2}
\]
under product norms.  In the Volterra--Perron chart, the \(b_i\)'s are controlled by the inverse
rough-maturity rank, directed-route rank, Perron reduced-resolvent conditioning, carrier Gram
lower bound, synchronization rank, and any augmented residual-selector rank.
\end{theorem}

\begin{proof}
Given \(y=Jh\), solve recursively.  First \(h_1=L_1y_1\), so
\(\norm{h_1}\le b_1\norm y\).  For \(i\ge2\),
\[
        h_i=L_i\left(y_i-\sum_{j<i}C_{ij}h_j\right).
\]
Thus
\[
        \norm{h_i}
        \le
        b_i\left(\norm y+\sum_{j<i}c_{ij}\norm{h_j}\right).
\]
The displayed recursion gives \(\norm{h_i}\le B_i\norm y\).  Summing over blocks gives the
claimed operator norm bound.
\end{proof}

\begin{definition}[Finite maturity observation operator]\label{def:finite-maturity-operator}
Let \(I\subset(0,\infty)\) be an open maturity interval and let
\[
        R_\Theta:I\to\R^m
\]
be a vector of option-response functions generated by a state \(\Theta\).  For a finite maturity
set \(\tau=(\tau_1,\dots,\tau_k)\subset I\), define
\[
        \mathcal O_\tau(\Theta)
        :=
        \bigl(R_\Theta(\tau_1),\dots,R_\Theta(\tau_k)\bigr)\in\R^{mk}.
\]
\end{definition}

\begin{theorem}[Finite panels recover analytic source charts]\label{thm:finite-panel-recovery}
Assume the source chart is finite dimensional after quotienting by gauge, and for each state the
response function \(\tau\mapsto R_\Theta(\tau)\) is real analytic on \(I\).  If the continuum
response map
\[
        \Theta\mapsto \{R_\Theta(\tau):\tau\in I\}
\]
has full local rank at \(\Theta_0\), then there exists a finite maturity set \(\tau\subset I\)
such that the finite observation operator \(\mathcal O_\tau\) has full local rank at
\(\Theta_0\).  Hence a finite index/single-name maturity panel locally recovers the canonical source state.
\end{theorem}

\begin{proof}
Let \(d\) be the dimension of the quotient source chart.  Full rank of the continuum response
means that the \(d\) derivative functions
\[
        \tau\mapsto D R_{\Theta_0}(\tau)[v_j],
        \qquad j=1,\dots,d,
\]
are linearly independent for some basis \(v_1,\dots,v_d\) of the quotient tangent space.  Let
\((q_n)_{n\ge1}\) be a countable dense subset of \(I\).  If every finite maturity evaluation were
rank deficient, then for each \(n\) there would be a vector \(a^{(n)}\in\R^d\), \(\norm{a^{(n)}}=1\),
such that
\[
        \sum_{j=1}^d
        a_j^{(n)}D R_{\Theta_0}(q_k)[v_j]
        =
        0,
        \qquad k=1,\dots,n.
\]
By compactness of the unit sphere, pass to a subsequence \(a^{(n_\ell)}\to a\) with
\(\norm a=1\).  For each fixed \(k\), the preceding display holds for all sufficiently large
\(\ell\), hence
\[
        \sum_{j=1}^d a_jD R_{\Theta_0}(q_k)[v_j]=0.
\]
The analytic function
\[
        \tau\mapsto
        \sum_{j=1}^d a_jD R_{\Theta_0}(\tau)[v_j]
\]
therefore vanishes on a dense subset of \(I\), and hence identically.  This contradicts linear
independence.  Therefore some finite maturity set has full rank.  The finite-dimensional inverse
function theorem gives local recovery from that finite panel.
\end{proof}

\begin{corollary}[Generic finite analytic panels]\label{cor:generic-finite-panel}
Under the assumptions of Theorem~\ref{thm:finite-panel-recovery}, there exist \(k\) and a
\(d\times d\) minor of the finite-panel Jacobian whose determinant is a nonzero real-analytic
function of \((\tau_1,\dots,\tau_k)\in I^k\).  Hence the set of \(k\)-point maturity panels with
deficient local rank is contained in the zero set of a nonzero real-analytic function and has
Lebesgue measure zero in \(I^k\).
\end{corollary}

\begin{proof}
The proof of Theorem~\ref{thm:finite-panel-recovery} gives at least one finite panel whose
Jacobian has rank \(d\).  Therefore some \(d\times d\) minor is nonzero at that panel.  The
entries of the finite-panel Jacobian are analytic in the maturity coordinates, so this minor is a
nonzero analytic function.  The zero set of a nonzero real-analytic function has Lebesgue measure
zero.
\end{proof}

\begin{corollary}[Finite-panel observability determinant]\label{cor:observability-determinant}
For a fixed finite panel \(\tau\), let \(J_\tau\) be the quotient finite-panel Jacobian at the
reference state.  If the quotient chart has dimension \(d\) and \(J_\tau\) has at least \(d\)
quote rows, define
\[
        \mathfrak D_t(\tau):=\det(J_\tau^\ast J_\tau).
\]
Then \(\mathfrak D_t(\tau)>0\) is equivalent to local finite-panel observability.  When this
holds,
\[
        \norm{J_\tau^{-}}
        \le
        \sigma_{\min}(J_\tau)^{-1}
        =
        \mathfrak D_t(\tau)^{-1/2}
        \prod_{j=1}^{d-1}\sigma_j(J_\tau),
\]
where \(\sigma_1(J_\tau)\ge\cdots\ge\sigma_d(J_\tau)>0\) are the singular values.  Loss of
recovery is therefore visible as the finite-panel observability determinant approaching zero.
\end{corollary}

\begin{proof}
The quotient finite-panel map is locally observable exactly when its Jacobian has rank \(d\).  This
is equivalent to \(J_\tau^\ast J_\tau\) being positive definite, hence to
\(\mathfrak D_t(\tau)>0\).  If \(\sigma_1\ge\cdots\ge\sigma_d>0\) are the singular values, then
\[
        \mathfrak D_t(\tau)
        =
        \prod_{j=1}^d\sigma_j(J_\tau)^2.
\]
The Moore--Penrose left inverse has norm \(\sigma_d(J_\tau)^{-1}\), which gives the displayed
identity after solving the determinant formula for \(\sigma_d^{-1}\).  Thus determinant collapse is
exactly singular-value collapse in the quotient chart.
\end{proof}


\begin{definition}[Regularized finite-panel reconstruction]\label{def:regularized-panel-reconstruction}
Let \(F_\tau:\bar U\subset\bar E\to\R^M\) be a finite-panel observation map on the
gauge quotient, and let
\[
        \widehat y_\tau=F_\tau(u_0)+\eps_\tau
\]
be a finite quote vector with smoothing or observation error \(\eps_\tau\).  For a compact
recovery neighborhood \(K\subset\bar U\), a reference point \(u_{\mathrm{pr}}\in K\), and a
regularization level \(\alpha\ge0\), define
\[
        \widehat u_{\alpha,\tau}
        \in
        \argmin_{u\in K}
        \left\{
        \norm{F_\tau(u)-\widehat y_\tau}_{\R^M}^2
        +
        \alpha\norm{u-u_{\mathrm{pr}}}_{\bar E}^2
        \right\}.
\]
The recovered risk-neutral state is \(\widehat\Xi_{\alpha,\tau}^\QQ=\psi(\widehat u_{\alpha,\tau})\),
with the gauge normalizations fixed by the chart.
\end{definition}

\begin{theorem}[Regularized finite-panel recovery]\label{thm:regularized-panel-recovery}
Suppose \(F_\tau\) is continuous on the compact recovery neighborhood \(K\) and has a local
inverse on \(K\) with Lipschitz constant \(C_\tau\):
\[
        \norm{u-v}_{\bar E}
        \le
        C_\tau\norm{F_\tau(u)-F_\tau(v)}_{\R^M},
        \qquad u,v\in K.
\]
Let \(u_0\in K\), let \(\widehat u_{\alpha,\tau}\) be any minimizer from
Definition~\ref{def:regularized-panel-reconstruction}, and set
\[
        r_{\mathrm{pr}}:=\norm{u_0-u_{\mathrm{pr}}}_{\bar E}.
\]
Then, whenever the minimizer remains in the same recovery neighborhood,
\[
        \norm{\widehat u_{\alpha,\tau}-u_0}_{\bar E}
        \le
        C_\tau
        \left(
        \norm{\eps_\tau}_{\R^M}
        +
        \sqrt{\norm{\eps_\tau}_{\R^M}^2+\alpha r_{\mathrm{pr}}^2}
        \right)
        \le
        C_\tau
        \left(
        2\norm{\eps_\tau}_{\R^M}
        +
        \sqrt{\alpha}\,r_{\mathrm{pr}}
        \right).
\]
If the finite-panel observability matrix has least quotient singular value
\(\beta_\tau>0\), then \(C_\tau\) may be chosen proportional to \(\beta_\tau^{-1}\) on a
sufficiently small neighborhood.  Thus finite-panel recovery is stable only to the extent that
quote error, smoothing error, regularization bias, and observability conditioning are controlled.
\end{theorem}

\begin{proof}
A minimizer exists because \(K\) is compact and the objective in
Definition~\ref{def:regularized-panel-reconstruction} is continuous on \(K\).  By optimality of
\(\widehat u_{\alpha,\tau}\),
\[
        \norm{F_\tau(\widehat u_{\alpha,\tau})-\widehat y_\tau}^2
        +
        \alpha\norm{\widehat u_{\alpha,\tau}-u_{\mathrm{pr}}}^2
        \le
        \norm{F_\tau(u_0)-\widehat y_\tau}^2
        +
        \alpha\norm{u_0-u_{\mathrm{pr}}}^2
        =
        \norm{\eps_\tau}^2+\alpha r_{\mathrm{pr}}^2 .
\]
Therefore
\[
        \norm{F_\tau(\widehat u_{\alpha,\tau})-\widehat y_\tau}
        \le
        \sqrt{\norm{\eps_\tau}^2+\alpha r_{\mathrm{pr}}^2}.
\]
The triangle inequality gives
\[
        \norm{F_\tau(\widehat u_{\alpha,\tau})-F_\tau(u_0)}
        \le
        \sqrt{\norm{\eps_\tau}^2+\alpha r_{\mathrm{pr}}^2}
        +
        \norm{\eps_\tau}.
\]
Applying the local inverse estimate yields the first displayed bound.  The second follows from
\(\sqrt{a^2+b^2}\le a+b\).  The singular-value statement is the finite-dimensional inverse bound
for the quotient observability matrix.
\end{proof}

\begin{corollary}[Sieve consistency of recovered source states]\label{cor:sieve-consistency}
Consider a sequence of finite panels \(\tau_n\), recovery neighborhoods \(K_n\), and source
charts whose approximation error to the target canonical state is \(a_n\).  If
\[
        C_{\tau_n}
        \left(
        \norm{\eps_{\tau_n}}
        +
        \sqrt{\alpha_n}\,r_{\mathrm{pr},n}
        \right)
        +a_n
        \longrightarrow0,
\]
then the reconstructed states converge to the target canonical risk-neutral state in quotient
distance.
In particular, when \(C_{\tau_n}r_{\mathrm{pr},n}\) is bounded, a convenient sufficient
regularization schedule is
\[
        \alpha_n\downarrow0,
        \qquad
        \frac{\norm{\eps_{\tau_n}}^2}{\alpha_n}\longrightarrow0,
        \qquad
        C_{\tau_n}\norm{\eps_{\tau_n}}\longrightarrow0,
        \qquad
        a_n\longrightarrow0 .
\]
The middle condition is the quotient-space analogue of the classical discrepancy scaling.  The
penalty vanishes asymptotically, but not so quickly that the reconstruction can spend the noise
envelope as source motion.
\end{corollary}

\begin{proof}
Apply Theorem~\ref{thm:regularized-panel-recovery} inside each finite source chart and add the
chart approximation error \(a_n\) by the triangle inequality.
\end{proof}

\begin{theorem}[Volterra--Perron option-manifold equivalence theorem]\label{thm:identification-frontier}\label{thm:spine}
Fix date \(t\), an arbitrage-free public option manifold, an ambient latent state space
\(\mathfrak X_t\), and a gauge-normalized Volterra--Perron source family
\(\mathfrak S_t\subset\mathfrak X_t\).  Let
\[
        \bar F_t:\mathfrak S_t/\!\sim_g\longrightarrow \mathfrak X_t^{\mathrm{pub}}
\]
be the public response on the source gauge quotient.  Around a reference source state, the
following three statements are equivalent.
\begin{enumerate}[leftmargin=2em]
\item The canonical \(\QQ\)-state is locally and stably recoverable from the public option manifold,
up to the stated scale, sign, latent-basis, and Perron-normalization gauge conventions:
the public response has a locally Lipschitz inverse on the model image of the recovery stratum.
\item The source quotient embeds locally and regularly into the public-observable quotient:
\(\bar F_t\) is a \(C^1\) local embedding on the recovery stratum, its image is a local embedded
submanifold of \(\mathfrak X_t^{\mathrm{pub}}\), and its inverse on the image is locally
Lipschitz in the quote norm.
\item The quotient observability operator has no source kernel and has positive local modulus:
\[
        \Ker D\bar F_t=\{0\},
        \qquad
        \mathfrak m_t
        :=
        \inf_{\norm{h}_{\bar E}=1}\norm{D\bar F_t h}_{\cY_t}
        >0
\]
on the recovery stratum.
\end{enumerate}
In the block-triangular source chart, these equivalent conditions are verified by a
full-rank diagonal-block certificate.  The maturity/source-response block, directed
covariance-route block, and index/carrier synchronization block each have a bounded left inverse
on their quotient domains.  Thus maturity scaling separates rough source geometry, carrier and
basket responses separate active Volterra routes, and synchronization responses separate the
simple Perron coordinate.
When these equivalent conditions hold, the local inverse satisfies
\[
        \dist_{\bar E}\bigl(\bar F_t^{-1}(y_1),\bar F_t^{-1}(y_2)\bigr)
        \le
        C\,\mathfrak m_t^{-1}\norm{y_1-y_2}_{\cY_t}
\]
for nearby public manifolds in the model image.  For a finite panel with least quotient singular
value \(\beta_t>0\), deterministic recovery error is at most order
\(\varepsilon/\beta_t\) under quote error \(\varepsilon\), and no reconstruction rule can improve
this order uniformly.

The residual statement is separate from state recovery.  For a latent loading class
\(\cL_{t,T}\), selected residual prices collapse to singletons if and only if the admitted public
payoff module is complete for that class:
\[
        \cL_{t,T}\subseteq \cA_{t,T}^{\mathrm{pub}} .
\]
Outside public completion, the remaining unidentified directions are exactly the
public-annihilator residual pricing fibers generated by public-equivalent density families.  Thus
recovery of the Volterra--Perron state does not by itself imply residual-free pricing,
tradeability, or hedgeability of every latent-stress claim.  Dynamic admissibility further requires
the static calendar convex-order restrictions and, in a recovered semimartingale chart, the
martingale drift restrictions for admitted option coordinates.
\end{theorem}

\begin{proof}
The public option manifold identifies exactly \(\mathfrak X_t^{\mathrm{pub}}\) by
Theorem~\ref{thm:public-quotient-maximality}.  A recovery rule on the source family must therefore
invert the induced map from \(\mathfrak S_t/\!\sim_g\) into this public quotient.  Local stable
recovery is exactly local Lipschitz invertibility of that induced map on the model image, which
is the embedded-stratum statement.

In a recovery chart, a bounded left inverse of the quotient derivative gives the local inverse
estimate of Theorem~\ref{thm:conditional-recovery}.  In finite dimension this is equivalent to a
positive least quotient singular value; in the Banach chart it is the bounded-below hypothesis used
there.  Conversely, a \(C^1\) local embedding with locally Lipschitz inverse has no quotient-source
kernel and has positive local modulus on the recovery stratum.  Singular analytic identifiability
with zero derivative belongs to a seam and is treated separately by
Theorem~\ref{thm:holder-recovery-seams}.

The block-rank condition is the triangular verification mechanism proved in
Theorem~\ref{thm:block-rank}.  It constructs a bounded left inverse by solving successively for the
maturity, directed-route, and synchronization coordinates.  If a diagonal response block ceases to
be bounded below and the lower blocks do not supply an independent response for that lost quotient
direction, the public map has either a kernel or a vanishing inverse modulus along that direction.
The finite-panel singular values record this obstruction in the block chart.

The stability estimate is Theorem~\ref{thm:recovery-stability}; the finite-panel upper bound is
Theorem~\ref{thm:regularized-panel-recovery}; and the matching lower bound is
Theorem~\ref{thm:inverse-modulus-lower-bound}.  The residual clause is
Theorem~\ref{thm:public-completion}.  Its proof identifies residual-free pricing with membership
of the claim loading in the public payoff module and identifies non-uniqueness with the
public-annihilator directions of Lemma~\ref{lem:residual-annihilator}.  The dynamic admissibility
clause combines Theorem~\ref{thm:calendar-convex-order} with the recovered-state martingale
restriction proved later in Theorem~\ref{thm:martingale-restriction}.
\end{proof}

\begin{remark}[What the equivalence does not say]
The theorem separates recoverability, residual-free pricing, and tradeable hedgeability:
\[
        \text{recoverability}\ne\text{residual-free pricing}\ne\text{tradeable hedgeability}.
\]
The option manifold may recover the state while non-public claims still lie in residual pricing
fibers.  Conversely, a public claim can be priced or hedged without recovering the full
Volterra--Perron state.  This distinction is part of the mathematical-finance content of the
theory rather than a limitation of the recovery theorem.
\end{remark}

\subsection{Failure modes}

The equivalence theorem also classifies non-identification.  Failures can occur at the public
no-arbitrage layer, the source quotient, the carrier span, the Perron mode, the maturity
dictionary, the residual payoff module, or the dynamic martingale layer.

\begin{definition}[Failure certificates]\label{def:failure-atlas-certificates}
For a finite option panel and a fixed source chart, the failure certificate vector is the
collection
\[
        \mathfrak C_t
        :=
        \left(
        \delta_{\mathrm{arb}},\delta_{\mathrm{cx}},
        \operatorname{rank}W_{\cE},
        \lambda_{\min}\Gamma_{\mathrm{car}},
        \Delta_P,
        \kappa_P,
        \sigma_{\min}J_{\mathrm{mat}},
        \mathfrak m_t,
        \dim \cN_{t,T}^{\mathrm{pub}},
        \operatorname{width}\mathcal I_t(X)
        \right).
\]
Here \(\delta_{\mathrm{arb}}\) is distance to the static no-arbitrage cone,
\(\delta_{\mathrm{cx}}\) is a convex-order violation across maturities, \(W_{\cE}\) is the basket
design matrix, \(\Gamma_{\mathrm{car}}\) is the carrier Gram matrix, \(\Delta_P\) is the Perron
spectral gap, \(\kappa_P\) is the reduced-resolvent norm for the simple Perron eigenvalue, and
\(J_{\mathrm{mat}}\) is the maturity-response Jacobian,
\[
        \mathfrak m_t
        :=
        \inf_{\norm{h}_{\bar E}=1}\norm{\overline{\mathcal O}_t^\QQ h}_{\cY_t}
\]
is the quotient observability modulus, \(\cN_{t,T}^{\mathrm{pub}}\) is the public annihilator, and
\(\mathcal I_t(X)\) is the residual price interval for \(X\).
\end{definition}

\begin{theorem}[Failure-mode theorem]\label{thm:failure-atlas}
Consider a finite-dimensional Volterra--Perron source chart satisfying the gauge-slice,
smoothness, and no-arbitrage assumptions above.  Stable dynamic selected pricing from the public
option panel can fail only through one or more of the following local mechanisms:
\begin{enumerate}[leftmargin=2em]
\item the quote vector violates static no-arbitrage or calendar convex-order admissibility;
\item the public response is not locally injective on the source gauge quotient, equivalently
the observability modulus vanishes;
\item basket or carrier design leaves a nontrivial covariance or synchronization fiber, visible
through rank loss in \(W_{\cE}\) or degeneracy of \(\Gamma_{\mathrm{car}}\);
\item the Perron reduced-resolvent norm diverges, in particular through spectral-gap collapse or
non-normal conditioning, so the dominant synchronization mode is no longer a smooth coordinate;
\item the maturity-response dictionary loses rank, so rough or Volterra source coordinates
cannot be separated by the admitted maturity panel;
\item the public payoff module is incomplete for the target claim, so the residual price fiber
has nonzero width;
\item the recovered finite-dimensional state violates the martingale drift restrictions for
admitted option coordinates.
\end{enumerate}
Conversely, away from these failure strata, and under correct source-family specification, the
equivalence theorem gives local gauge-class recovery, stability controlled by
\(\mathfrak m_t^{-1}\), residual-free pricing for publicly completed claims, and dynamically
admissible recovered prices subject to the martingale equations.
\end{theorem}

\begin{proof}
Static and dynamic no-arbitrage are necessary before any risk-neutral interpretation exists.
Local source recovery is equivalent to local injectivity of the response on the gauge quotient;
in finite dimension, stable recovery requires positive quotient singular value.  Basket and
carrier rank failures produce explicit kernels in the joint-observation block.  Perron
reduced-resolvent blow-up invalidates differentiable spectral coordinates; spectral-gap collapse
is one mechanism, and non-normal conditioning is another.  Maturity-rank collapse creates kernel
directions in the Volterra dictionary.  Public payoff incompletion is the existence of a
public-annihilating residual direction that moves the target claim.  The martingale condition is
the dynamic no-arbitrage restriction after recovery.  These mechanisms are the local ways in
which one clause of Theorem~\ref{thm:identification-frontier} can fail.
\end{proof}

\begin{definition}[Singularly regular observation map]\label{def:singular-regular-map}
Let \(F:U\subset\R^d\to\R^M\) be an analytic finite-panel response map on a gauge slice.  The
map is \(m\)-regular at \(u_0\) on a closed cone \(\cC\subset\R^d\) if there are \(c>0\) and
\(r>0\) such that, for all \(u,v\in u_0+\cC\) with
\(\norm{u-u_0},\norm{v-u_0}<r\),
\[
        \norm{F(u)-F(v)}
        \ge
        c\norm{u-v}^m .
\]
The case \(m=1\) is ordinary Lipschitz recovery.  The case \(m>1\) is singular but still
identifiable on the cone.
\end{definition}

\begin{theorem}[H\"older recovery across analytic seams]\label{thm:holder-recovery-seams}
If \(F\) is \(m\)-regular at \(u_0\) on the gauge-normalized cone \(\cC\), then \(F\) is locally
injective on that cone and its inverse satisfies
\[
        \norm{u-v}
        \le
        c^{-1/m}\norm{F(u)-F(v)}^{1/m}.
\]
Consequently, let \(y=F(u)\), let \(\norm{\widehat y-y}\le\delta\), and let a reconstruction
\(\widehat u\in u_0+\cC\) satisfy
\[
        \norm{F(\widehat u)-\widehat y}=o(\delta).
\]
Then the local recovery bound
\[
        \norm{\widehat u-u}
        \le
        c^{-1/m}\delta^{1/m}+o(\delta^{1/m})
\]
holds on the same local cone.
\end{theorem}

\begin{proof}
If \(F(u)=F(v)\) for two points in the cone, the defining inequality forces \(u=v\).  Thus \(F\)
is locally injective on that cone.  The same inequality rearranges to the displayed inverse bound.
For the noisy estimate,
\[
        \norm{F(\widehat u)-F(u)}
        \le
        \norm{F(\widehat u)-\widehat y}
        +
        \norm{\widehat y-y}
        =
        \delta+o(\delta).
\]
Applying the inverse bound gives
\[
        \norm{\widehat u-u}
        \le
        c^{-1/m}(\delta+o(\delta))^{1/m}
        =
        c^{-1/m}\delta^{1/m}+o(\delta^{1/m}).
\]
A nonzero projection residual would enter through the first term in the preceding triangle
inequality.
\end{proof}

\begin{definition}[Pseudo-true recovered state]\label{def:pseudo-true-state}
Let \(F:U\subset\R^d\to\R^M\) be a finite-panel response map and let \(y\in\R^M\) be an
arbitrage-free quote vector, not necessarily in \(F(U)\).  For a positive definite weight \(W\),
a pseudo-true recovered state is
\[
        u^\dagger
        \in
        \argmin_{u\in U}
        (y-F(u))^\top W(y-F(u)).
\]
The model residual is \(r^\dagger:=y-F(u^\dagger)\).
\end{definition}

\begin{theorem}[Misspecified projection and normal equations]\label{thm:misspecified-projection}
Assume \(F\) is \(C^2\), \(u^\dagger\) is an interior minimizer, and
\(J^\dagger:=DF(u^\dagger)\) has full column rank.  Then
\[
        (J^\dagger)^\top W r^\dagger=0.
\]
If
\[
        (J^\dagger)^\top WJ^\dagger
        -
        \sum_{a=1}^M (Wr^\dagger)_aD^2F_a(u^\dagger)
\]
is positive definite, then \(u^\dagger\) is a strict local projection.  Small quote perturbations
move \(u^\dagger\) differentiably with derivative
\[
        \dd u^\dagger
        =
        \left[
        (J^\dagger)^\top WJ^\dagger
        -
        \sum_{a=1}^M (Wr^\dagger)_aD^2F_a(u^\dagger)
        \right]^{-1}
        (J^\dagger)^\top W\,\dd y .
\]
When \(r^\dagger=0\), this reduces to the Gauss--Newton inverse.
\end{theorem}

\begin{proof}
For
\[
        Q(u;y)=(y-F(u))^\top W(y-F(u)),
\]
the first variation is
\[
        D_uQ(u;y)[h]
        =
        -2h^\top DF(u)^\top W(y-F(u)).
\]
At an interior minimizer this gives the normal equation
\((J^\dagger)^\top Wr^\dagger=0\).  The Hessian in the state variable is the displayed
positive-definite matrix up to the common factor two, so the second-order condition gives a strict
local projection.  Applying the implicit-function theorem to the normal equation
\(DF(u)^\top W(y-F(u))=0\) and differentiating with respect to \(y\) gives the stated sensitivity
formula.
\end{proof}

\begin{definition}[Weighted finite-panel estimator]\label{def:weighted-finite-panel-estimator}
Let \(F:\R^d\to\R^M\) be the finite-panel response map on a gauge slice.  Observed quotes satisfy
\[
        Y_n=F(u_0)+n^{-1/2}\eps_n,
        \qquad
        \eps_n\Rightarrow \mathcal N(0,\Omega),
\]
where \(\Omega\) is positive definite on the quote coordinates.  For a positive definite weight
\(W\), define
\[
        \widehat u_n
        \in
        \argmin_{u\in U}(Y_n-F(u))^\top W(Y_n-F(u)).
\]
\end{definition}

\begin{theorem}[Asymptotic covariance of recovered states]\label{thm:asymptotic-recovery-covariance}
Assume \(F\) is \(C^2\), \(J:=DF(u_0)\) has full column rank, \(u_0\) is an interior point of the
gauge slice, and the estimator in Definition~\ref{def:weighted-finite-panel-estimator} is
locally consistent.  Then
\[
        \sqrt n(\widehat u_n-u_0)
        \Rightarrow
        \mathcal N\!\left(
        0,\,
        (J^\top WJ)^{-1}J^\top W\Omega WJ(J^\top WJ)^{-1}
        \right).
\]
The efficient weight is \(W=\Omega^{-1}\), giving covariance
\[
        \mathcal I(u_0)^{-1}
        :=
        (J^\top\Omega^{-1}J)^{-1}.
\]
This is the local Cram\'er--Rao lower bound for linear functionals in the Gaussian quote
experiment generated by the same panel.
\end{theorem}

\begin{proof}
The first-order condition is
\[
        DF(\widehat u_n)^\top W(Y_n-F(\widehat u_n))=0.
\]
Local consistency and the expansion \(Y_n=F(u_0)+n^{-1/2}\eps_n\) give
\[
        J^\top W\eps_n
        -
        J^\top WJ\sqrt n(\widehat u_n-u_0)
        +o_p(1)=0.
\]
Thus
\[
        \sqrt n(\widehat u_n-u_0)
        =
        (J^\top WJ)^{-1}J^\top W\eps_n+o_p(1),
\]
and the stated covariance follows from \(\eps_n\Rightarrow\mathcal N(0,\Omega)\).  Minimizing this
Loewner covariance over positive definite weights gives the generalized least-squares weight
\(W=\Omega^{-1}\), hence the inverse Fisher information of the local Gaussian shift experiment.
\end{proof}

\begin{corollary}[Quote-design criterion]\label{cor:quote-design}
For a candidate finite quote panel with Jacobian \(J_P\) and quote-noise covariance
\(\Omega_P\), the local information matrix is
\[
        \mathcal I_P:=J_P^\top\Omega_P^{-1}J_P.
\]
A panel is locally well designed for recovery if \(\lambda_{\min}(\mathcal I_P)\) is bounded
away from zero.  For a scalar functional \(g(u)\), the local asymptotic variance is
\[
        \nabla g(u_0)^\top\mathcal I_P^{-1}\nabla g(u_0).
\]
Thus the theory identifies which strikes, maturities, carriers, and baskets improve recovery of
Perron stress, roughness, or directed-route coefficients.
\end{corollary}

\begin{proof}
Theorem~\ref{thm:asymptotic-recovery-covariance} gives asymptotic covariance
\((J_P^\top\Omega_P^{-1}J_P)^{-1}\) for the efficiently weighted local recovery estimator.  Hence
\(\lambda_{\min}(\mathcal I_P)>0\) is the finite-panel information condition, and a lower bound on
this eigenvalue prevents loss of quotient directions.  The scalar-functional variance follows from
the delta method applied to \(g\).
\end{proof}


\section{Adapted recovery of the risk-neutral state process}\label{sec:dynamic-recovery}

Dynamic observability follows the recovered-state construction of \cite{vidal2026dynamic}.  The
preceding theorems are pointwise in \(t\).  A dynamic pricing model also requires the recovered
state to be adapted through time.  Pointwise inverses assemble into an observable \(\QQ\)-state
process only under measurable charts and uniform quotient conditioning.

\begin{definition}[Local recovery operator]\label{def:local-recovery-operator}
On a recovery neighborhood at date \(t\), define
\[
        \mathcal I_t:W_t\cap F_t(V_t)\to\overline{\mathfrak X}_t,
        \qquad
        \mathcal I_t(y):=\bar F_t^{-1}(y),
\]
where \(F_t=\mathcal R_t\circ\psi_t\).  The recovered canonical state is
\[
        \widehat\Xi_t^\QQ
        :=
        \mathcal I_t(\cM_t^{\mathrm{opt}}).
\]
\end{definition}

\begin{assumption}[Uniform recovery window]\label{ass:uniform-recovery}
On a time interval \([0,T_0]\), there are \(\cH_t\)-measurable recovery events
\(\Omega_t^{\mathrm{rec}}\) on which the realized option manifold remains in local recovery
neighborhoods and the realized response maps satisfy the rank conditions with deterministic
lower bounds.  On these events the recovery charts can be chosen so that:
\begin{enumerate}[leftmargin=2em]
\item \((t,y)\mapsto\mathcal I_t(y)\) is Borel measurable;
\item the inverse response norms satisfy \(\sup_{t\le T_0}\norm{L_t}<\infty\);
\item the Perron gaps, Perron reduced-resolvent constants, and carrier Gram lower bounds satisfy
\[
        \inf_{t\le T_0}\gamma_P(t)>0,
        \qquad
        \sup_{t\le T_0}\kappa_P(t)<\infty,
        \qquad
        \inf_{t\le T_0}\kappa_{\mathrm{car}}(t)>0;
\]
\item the diagonal maturity, directed-Volterra, and synchronization response blocks have uniformly
bounded left inverses; augmented residual-instrument blocks, when used, satisfy the same
condition.
\end{enumerate}
\end{assumption}

\begin{theorem}[Adapted dynamic recovery]\label{thm:adapted-dynamic-recovery}
Suppose Assumption~\ref{ass:uniform-recovery} holds.  Then the recovered canonical state process
\[
        \widehat\Xi_t^\QQ=\mathcal I_t(\cM_t^{\mathrm{opt}})
\]
is \(\cH_t\)-measurable on \(\Omega_t^{\mathrm{rec}}\) for every \(t\).  If the public option manifold is
\((\cH_t)\)-adapted, then \((\widehat\Xi_t^\QQ)_{t\le T_0}\) is an adapted canonical
Volterra--Perron state process.
\end{theorem}

\begin{proof}
The public option manifold \(\cM_t^{\mathrm{opt}}\) is an \(\cH_t\)-measurable quote object.
Assumption~\ref{ass:uniform-recovery} gives a Borel recovery map on the recovery event at date
\(t\).  The composition \(\mathcal I_t(\cM_t^{\mathrm{opt}})\) is therefore
\(\cH_t\)-measurable on \(\Omega_t^{\mathrm{rec}}\).  The argument is pointwise in \(t\), and the
adaptedness statement follows for the recovered process on the interval.
\end{proof}

\begin{theorem}[Pathwise stability of recovered states]\label{thm:pathwise-stability}
Under Assumption~\ref{ass:uniform-recovery}, there is a constant \(C_*>0\) such that for any two
public option-manifold paths \(M,\widetilde M\) remaining in the same recovery neighborhoods,
\[
        \sup_{t\le T_0}
        \dist_{\overline{\mathfrak X}_t}
        \!\left(\mathcal I_t(M_t),\mathcal I_t(\widetilde M_t)\right)
        \le
        C_*
        \sup_{t\le T_0}\norm{M_t-\widetilde M_t}_{\cY_t}.
\]
If \(t\mapsto\mathcal I_t\) is locally continuous and \(t\mapsto\cM_t^{\mathrm{opt}}\) is
c\`adl\`ag in the quote-space topology, then \(t\mapsto\widehat\Xi_t^\QQ\) is c\`adl\`ag in
the canonical state topology.
\end{theorem}

\begin{proof}
Pointwise stability gives a Lipschitz constant controlled by the inverse response
norm, Perron reduced-resolvent bound, carrier Gram lower bound, and diagonal response-block inverse
norms.  The uniform recovery-window assumption bounds these quantities on \([0,T_0]\); taking the
supremum over \(t\) gives the displayed estimate.  The c\`adl\`ag statement follows by composing a
c\`adl\`ag option-manifold path with locally continuous recovery maps.
\end{proof}

\begin{definition}[Recovery atlas]\label{def:recovery-atlas}
A recovery atlas is a family of gauge-normalized charts
\((U_\alpha,F_\alpha,\mathcal I_\alpha)_{\alpha\in A}\) such that
\[
        F_\alpha:U_\alpha\to\cY_t
\]
is locally injective on the canonical quotient, \(\mathcal I_\alpha\) is its inverse on the
model image, and on overlaps the inverses agree up to the prescribed gauge transition:
\[
        \mathcal I_\beta(y)
        =
        g_{\alpha\beta}\bigl(\mathcal I_\alpha(y)\bigr),
        \qquad
        y\in F_\alpha(U_\alpha)\cap F_\beta(U_\beta).
\]
The atlas is uniformly observable on a region \(\cO\) if all local inverse constants are bounded
there by a common constant.
\end{definition}

\begin{theorem}[Patching local recovery]\label{thm:patching-local-recovery}
Let \(\cM_t^{\mathrm{opt}}\) be an adapted option-manifold path taking values in the image of a
uniformly observable recovery atlas on a region \(\cO\).  Suppose the chart transitions are
Borel.  Then the local recovered states patch into a well-defined adapted canonical state
process on the stopping interval \([0,\tau_{\mathrm{loss}})\), where
\[
        \tau_{\mathrm{loss}}
        :=
        \inf\{t:\cM_t^{\mathrm{opt}}\notin\cO
        \text{ or an observability boundary is hit}\}.
\]
The observability boundary may include
\[
        \gamma_P(t)=0,\qquad
        \lambda_{\min}(\Gamma_{\mathrm{car},t})=0,\qquad
        \mathfrak D_t(\tau)=0,
\]
as well as active-edge transitions \(V_{ij,t}=0\) for screened routes whose stratum is being
held fixed.
\end{theorem}

\begin{proof}
On each chart domain, adaptedness follows from the local inverse theorem.  On overlaps,
gauge-compatible transition maps identify the same canonical state.  The first exit time from
the observable region is a stopping time when the boundary functions are adapted and the
recovery region is measurable.  Uniform inverse bounds give pathwise stability before exit.
\end{proof}

\begin{remark}
The dynamic recovery theorem is an adapted inverse-map result, with no observation-noise filtering
model.  Statistical filtering would require an additional observation-error model for finite quotes,
bid--ask frictions, and microstructure smoothing.  The map constructed here is the
quotient inverse that such a filter would approximate.
\end{remark}

\begin{definition}[Recovered pricing semigroup]\label{def:recovered-pricing-semigroup}
Let \(U_t\) denote a finite-dimensional recovered quotient state on a recovery window.  For a
payoff class \(\cD\) and maturity increment \(\tau\ge0\), define
\[
        P_{t,t+\tau}f(U_t)
        :=
        \EE^\QQ\!\left[
        \frac{B_t}{B_{t+\tau}}f(U_{t+\tau})
        \,\middle|\,\cH_t
        \right].
\]
The option manifold prices a core \(\cD\) if these values are observable or statically spanned
for every \(f\in\cD\) and all sufficiently small \(\tau\).
\end{definition}

\begin{theorem}[Generator identification from short maturities]\label{thm:generator-identification}
Assume \(U\) is Markov under \(\QQ\) on the recovery window, and let \(\mathcal A_t\) be its
risk-neutral generator on a core \(\cD\).  If the recovered option manifold identifies
\(P_{t,t+\tau}f(U_t)\) for every \(f\in\cD\) and all sufficiently small \(\tau>0\), then
\[
        \mathcal A_t f(U_t)-r_tf(U_t)
        =
        \lim_{\tau\downarrow0}
        \frac{P_{t,t+\tau}f(U_t)-f(U_t)}{\tau}
\]
for every \(f\in\cD\) for which the limit exists.  If
\[
        \dd U_t=b(t,U_t)\dd t+\Sigma(t,U_t)\dd W_t^\QQ
\]
is an It\^o diffusion, then prices of coordinate functions and their quadratic products identify
first the discounted generator
\[
        L_t f(U_t)
        :=
        \lim_{\tau\downarrow0}
        \frac{P_{t,t+\tau}f(U_t)-f(U_t)}{\tau}
        =
        \mathcal A_t f(U_t)-r_tf(U_t),
\]
and hence the drift and covariance coefficients by
\[
        b_i(t,U_t)=L_tu_i(U_t)+r_tu_i(U_t),
\]
and
\[
        a_{ij}(t,U_t)
        =
        L_t(u_iu_j)(U_t)
        +
        r_tu_i(U_t)u_j(U_t)
        -
        u_i(U_t)b_j(t,U_t)
        -
        u_j(U_t)b_i(t,U_t),
        \qquad a=\Sigma\Sigma^\top .
\]
\end{theorem}

\begin{proof}
On the core \(\cD\), the recovered short-maturity prices identify the discounted pricing
semigroup to first order:
\[
        P_{t,t+\tau}f(U_t)
        =
        f(U_t)+\tau\bigl(\mathcal A_tf(U_t)-r_tf(U_t)\bigr)+o(\tau).
\]
The displayed limit is therefore the generator of the discounted semigroup, so the option manifold
identifies \(L_t=\mathcal A_t-r_t\) on the core, and adding back \(r_tf\) recovers
\(\mathcal A_tf\).  For the diffusion statement, apply
\[
        \mathcal A_tf=b^\top\nabla f+\frac12\operatorname{Tr}(a\nabla^2f)
\]
to \(f(u)=u_i\) and \(f(u)=u_iu_j\), then replace \(\mathcal A_t\) by \(L_t+r_t\) in the
identified quantities.
\end{proof}

\begin{remark}[Ideal generator recovery]\label{rem:ideal-generator-recovery}
Theorem~\ref{thm:generator-identification} is an ideal short-maturity result.  It requires the
public manifold, or an admitted static span of it, to price the coordinate functions and
quadratic products in the chosen core.  Ordinary listed vanillas do not automatically contain
these state claims.  When the core is only approximated by a finite quote panel, the finite
state-law experiment below should be read as an adapted propagation and innovation diagnostic,
not as exact identification of the infinitesimal generator.
\end{remark}

\begin{definition}[Finite adapted state-law experiment]\label{def:finite-state-law-experiment}
Let \(0=t_0<t_1<\cdots<t_n\le T_0\) be observation dates inside a recovery window, and let
\[
        U_k:=\widehat\Xi_{t_k}^{\QQ}\in\R^d
\]
be a gauge-normalized recovered coordinate vector in a fixed chart.  Write
\(\Delta_k=t_{k+1}-t_k\), \(\Delta U_k=U_{k+1}-U_k\), and
\(\cG_k=\cH_{t_k}\vee\sigma(U_0,\dots,U_k)\).  A finite adapted state-law experiment specifies,
for each \(k\), a coefficient estimate
\[
        \widehat b_{k}^{-}(u)=\widehat a_k+\widehat B_ku
\]
and positive semidefinite innovation matrix \(\widehat A_k^{-}\), both
\(\cG_k\)-measurable and estimated only from transitions
\((U_j,U_{j+1})\) with \(j<k\).  The one-step state-law innovation is
\[
        \varepsilon_{k+1}
        :=
        U_{k+1}-U_k-\Delta_k\widehat b_k^-(U_k),
\]
and the persistence innovation is \(\varepsilon_{k+1}^{0}:=\Delta U_k\).
\end{definition}

\begin{definition}[Phase-local diffusion and Gaussian seam energy]\label{def:phase-local-seam}
Let \(\phi_k=\phi(U_k)\) be any hard or smooth phase label generated by the quotient phase
coordinates of Section~\ref{sec:phases}.  Let \(\widehat A_k^{\mathrm{glob}}\) be a shrunk
covariance estimate of the scaled past innovations
\[
        \frac{\varepsilon_{j+1}}{\sqrt{\Delta_j}},
        \qquad j<k,
\]
and let \(\widehat A_k^{\phi}\) be the analogous estimate restricted to past indices
\(j<k\) with \(\phi_j=\phi_k\), whenever enough such observations exist.  Otherwise set
\(\widehat A_k^\phi=\widehat A_k^{\mathrm{glob}}\) and record that the phase-local covariance
was not separately identified.  With ridge parameter \(\eta_k\ge0\), define the global and
phase-local Gaussian seam energies
\[
        \mathfrak E_{k+1}^{\mathrm{glob}}
        =
        \frac1d
        \left(\frac{\varepsilon_{k+1}}{\sqrt{\Delta_k}}\right)^\top
        \left(\widehat A_k^{\mathrm{glob}}+\eta_k I\right)^\dagger
        \left(\frac{\varepsilon_{k+1}}{\sqrt{\Delta_k}}\right),
\]
and
\[
        \mathfrak E_{k+1}^{\phi}
        =
        \frac1d
        \left(\frac{\varepsilon_{k+1}}{\sqrt{\Delta_k}}\right)^\top
        \left(\widehat A_k^{\phi}+\eta_k I\right)^\dagger
        \left(\frac{\varepsilon_{k+1}}{\sqrt{\Delta_k}}\right).
\]
The same definitions with \(\varepsilon_{k+1}^{0}\) give the corresponding persistence seam
energies.
\end{definition}

\begin{proposition}[Adaptedness of finite state-law diagnostics]\label{prop:finite-state-law-adapted}
In the finite adapted state-law experiment, the fitted drift
\(\widehat b_k^-\), global covariance \(\widehat A_k^{\mathrm{glob}}\), phase-local covariance
\(\widehat A_k^\phi\), ridge \(\eta_k\), and forecast
\[
        U_k+\Delta_k\widehat b_k^-(U_k)
\]
are \(\cG_k\)-measurable.  Therefore the one-step propagation error, innovation scale, and seam
energies are prequential diagnostics: they evaluate the realized next recovered state against an
object formed before \(U_{k+1}\) is observed.
\end{proposition}

\begin{proof}
All coefficient estimates and covariance estimates are functions only of \(\cG_k\)-measurable
past transitions and of the current recovered state \(U_k\).  The phase label \(\phi_k\) is a
measurable function of \(U_k\).  Hence the fitted drift, covariance matrices, ridge, and forecast
are \(\cG_k\)-measurable.  The diagnostics are then obtained by comparing these known-at-\(t_k\)
objects with \(U_{k+1}\).
\end{proof}

\begin{theorem}[Finite state-law consistency]\label{thm:finite-state-law-consistency}
Suppose on a recovery window the recovered state satisfies the It\^o diffusion
\[
        \dd U_t=b(t,U_t)\dd t+\Sigma(t,U_t)\dd W_t^\QQ,
        \qquad A(t,U_t)=\Sigma(t,U_t)\Sigma(t,U_t)^\top ,
\]
with locally Lipschitz coefficients.  Let \(h_k=\max_{j<k}\Delta_j\to0\), and let the regression
and covariance estimators in Definition~\ref{def:finite-state-law-experiment} be consistent local
estimators of \(b(t_k,U_{t_k})\) and \(A(t_k,U_{t_k})\) on the chosen chart or phase stratum.
Then
\[
        \frac{\varepsilon_{k+1}}{\sqrt{\Delta_k}}
        =
        \Sigma(t_k,U_{t_k})Z_{k+1}+o_{\QQ}(1),
        \qquad Z_{k+1}\sim N(0,I)
\]
conditionally to first order, and the seam energy in
Definition~\ref{def:phase-local-seam} is a finite-information innovation statistic for the
recovered \(\QQ\)-state.  If the local affine drift is misspecified, the same statistic remains
an adapted propagation-reserve diagnostic, not a source-transition estimator for the true law.
\end{theorem}

\begin{proof}
The Euler expansion of the It\^o diffusion gives
\[
        U_{t_{k+1}}-U_{t_k}
        =
        b(t_k,U_{t_k})\Delta_k
        +
        \Sigma(t_k,U_{t_k})(W_{t_{k+1}}-W_{t_k})
        +
        o_{\QQ}(\sqrt{\Delta_k})
\]
under the stated regularity.  Consistency of the drift estimator makes the fitted drift error
lower order after multiplication by \(\Delta_k\).  Consistency of the covariance estimator gives
the displayed covariance-normalized innovation statistic.  Without correct local specification,
the same construction is still prequential by Proposition~\ref{prop:finite-state-law-adapted};
it is then a reserve or fragility measure rather than a recovered transition kernel.
\end{proof}

\begin{remark}[Role of the Gaussian seam metric]\label{rem:gaussian-seam-state-law}
The seam energy is the dynamic analogue of the finite Gaussian source-family tests used in the
recovery layer.  Raw Euclidean propagation errors depend on the chosen coordinate scale.  The
covariance-normalized innovation measures the error in the finite information metric of the
recovered chart.  Phase-local versions replace a pooled covariance by the covariance of comparable
past quotient states; they do not create a new phase by thresholding future errors.
\end{remark}

\begin{remark}[Pricing dynamics versus observed state dynamics]\label{rem:q-state-observation-measure}
The consistency statement above is a \(\QQ\)-dynamics statement.  A calendar time series of
option-implied recovered states is usually observed under the statistical data-generating measure,
even though each state is a risk-neutral object recovered from option prices.  Such a time series
therefore estimates observation-measure dynamics of \(\QQ\)-implied coordinates unless an
additional change-of-measure or state-space transfer map is imposed.  The finite prequential
diagnostics remain valid as adapted stability, propagation, and fragility diagnostics for the
observed option manifold; they are not direct estimates of the risk-neutral pricing generator
without that extra transfer map.
\end{remark}

\section{Martingale consistency of recovered risk-neutral dynamics}\label{sec:martingale}

Recovery turns the option manifold into a canonical state process.  Risk-neutral pricing imposes a
separate condition.  Once public prices are written as functionals of the recovered state, their
discounted dynamics must be local martingales.  The resulting drift equations are the dynamic
no-arbitrage restrictions on the Volterra--Perron coordinates.

\begin{definition}[Recovered price functional]\label{def:recovered-price-functional}
In this section \(a\) denotes a fixed calendar contract, such as
\((u,K,T,\kappa)\), rather than a rolling fixed-\(\tau\) surface point.  Let \(a\in\cI_t\) be an
admitted option coordinate.  On a recovery neighborhood, define
\[
        \Phi_a(t,\bar\Theta)
        :=
        C_t^a
        \quad\text{whenever}\quad
        \bar\Theta=\mathcal I_t(\cM_t^{\mathrm{opt}}).
\]
Equivalently, \(\Phi_a\) is the \(a\)-coordinate of the observable response map expressed in
canonical recovered-state coordinates.
\end{definition}

\begin{remark}
If the same restriction is written in rolling surface coordinates
\((u,\tau,m,\kappa,\chi)\), the generator must include the roll of \(\tau=T-t\) and, when \(m\) is
defined relative to a stochastic forward, the It\^o generator of \(m\).  In the simplest fixed
strike parametrization this adds the familiar convective term \(-\partial_\tau\) to the surface
PDE.
\end{remark}

\begin{assumption}[Finite-dimensional semimartingale chart]\label{ass:semimartingale-chart}
On a recovery window, the canonical quotient state admits local coordinates \(U_t\in\R^d\) and
has risk-neutral semimartingale dynamics
\[
        \dd U_t=b_t\dd t+\Sigma_t\dd W_t^\QQ,
\]
where \(b_t\in\R^d\), \(\Sigma_t\in\R^{d\times q}\), and the recovered price functionals
\(\Phi_a(t,u)\) are \(C^{1,2}\) in \((t,u)\) on the coordinate neighborhood.
\end{assumption}

\begin{theorem}[Recovered-state martingale restriction]\label{thm:martingale-restriction}
Suppose Assumption~\ref{ass:semimartingale-chart} holds and, for every admitted public option
coordinate \(a\), the discounted price process
\[
        e^{-\int_0^t r_s\dd s}\Phi_a(t,U_t)
\]
is a local \(\QQ\)-martingale.  Then, along the recovered state path,
\[
        \partial_t\Phi_a(t,U_t)
        +
        D\Phi_a(t,U_t)[b_t]
        +
        \frac12
        \operatorname{Tr}\!\left(
        \Sigma_t^\ast D^2\Phi_a(t,U_t)\Sigma_t
        \right)
        -
        r_t\Phi_a(t,U_t)
        =
        0
\]
for every admitted option coordinate \(a\), up to indistinguishability.
\end{theorem}

\begin{proof}
Apply It\^o's formula to \(e^{-\int_0^t r_s\dd s}\Phi_a(t,U_t)\).  The finite-variation drift is
the discount factor times
\[
        \partial_t\Phi_a
        +
        D\Phi_a[b_t]
        +
        \frac12\operatorname{Tr}(\Sigma_t^\ast D^2\Phi_a\Sigma_t)
        -
        r_t\Phi_a .
\]
A local martingale has zero finite-variation drift.  This gives the displayed restriction.
\end{proof}

\begin{corollary}[Dynamic no-arbitrage closure]\label{cor:dynamic-no-arbitrage-closure}
Conversely, suppose the drift restriction in Theorem~\ref{thm:martingale-restriction} holds for
all admitted coordinates and the stochastic integrals are well defined.  Then the generated
discounted public option price coordinates are local \(\QQ\)-martingales.  Hence the recovered
state dynamics define a dynamically arbitrage-free public option manifold on the admitted payoff
span, subject to the usual true-martingale integrability conditions.
\end{corollary}

\begin{proof}
The same It\^o calculation shows that the finite-variation part of each discounted coordinate
vanishes.  The remaining term is a stochastic integral and hence a local martingale.  A positive
linear conditional-expectation representation on the admitted span gives local no-arbitrage; true
martingale integrability upgrades local prices to genuine conditional-expectation prices.
\end{proof}

\begin{definition}[Recovered Musiela option surface]\label{def:recovered-musiela-surface}
Let \(U_t\) be a finite-dimensional recovered state chart.  For underlier \(u\), define the
discounted rolling surface coordinate in pricing chart \(\chi\in\{\mathrm B,\mathrm N\}\) by
\[
        \mathcal C^{u,\chi}(t,U_t,\tau,m^\chi)
        :=
        B_t^{-1}C_t^u(T=t+\tau,K=K^\chi(F_t^u,m^\chi)),
\]
where \(\tau>0\) is time to maturity,
\[
        K^{\mathrm B}(F,m)=Fe^m,\qquad K^{\mathrm N}(F,m)=F+m .
\]
For a fixed calendar contract \((u,K,T)\), set
\[
        \tau_t=T-t,\qquad
        m_t^{\mathrm B}=\log(K/F_t^u),
        \qquad
        m_t^{\mathrm N}=K-F_t^u .
\]
Thus the Black roll is written in log-moneyness, whereas the normal-forward roll is written in
cash-forward moneyness.  A standardized normal coordinate
\((K-F_t^u)/\sqrt{v_{u,t}^{\mathrm N}(\tau)}\) may be used instead, but then the generator also
contains the roll of the normal variance scale.
Let \(\mathcal L_{U,m^\chi}\) denote the generator of the joint state
\((U_t,m_t^\chi)\).
\end{definition}

\begin{theorem}[Surface roll restriction]\label{thm:surface-roll-restriction}
Suppose the discounted calendar-contract price
\[
        B_t^{-1}C_t^u(K,T)
        =
        \mathcal C^{u,\chi}(t,U_t,\tau_t,m_t^\chi)
\]
is a local \(\QQ\)-martingale and the joint state \((U_t,m_t^\chi)\) has generator
\(\mathcal L_{U,m^\chi}\).  Then, along the recovered state path,
\[
        \partial_t\mathcal C^{u,\chi}
        -
        \partial_\tau\mathcal C^{u,\chi}
        +
        \mathcal L_{U,m^\chi}\mathcal C^{u,\chi}
        =
        0 .
\]
In undiscounted coordinates the right-hand side is \(r_t\mathcal C^{u,\chi}\).  If
\(m^\chi\) is held fixed as a surface coordinate rather than induced by a fixed calendar strike,
the generator must be changed accordingly; an additional chart-specific moneyness-roll term
enters.
\end{theorem}

\begin{proof}
Apply It\^o's formula to
\(\mathcal C^{u,\chi}(t,U_t,T-t,m_t^\chi)\).  The deterministic roll of \(\tau_t\) contributes
\(-\partial_\tau\mathcal C^{u,\chi}\), while the joint diffusion and drift of
\((U_t,m_t^\chi)\) contribute
\(\mathcal L_{U,m^\chi}\mathcal C^{u,\chi}\).  The finite-variation part of a discounted local
martingale must vanish.
\end{proof}

\begin{definition}[Dynamically admissible residual selector]\label{def:dynamic-selector-admissible}
A residual selector \(\theta\) is dynamically admissible on a recovery window if the selected
price functionals
\[
        \Phi_X^\theta(t,U_t):=\Pi_t^\theta(X)
\]
for the selected latent-stress claims satisfy the same discounted local-martingale drift
restriction as public option coordinates, whenever they are included in the priced claim family.
\end{definition}

\begin{proposition}[Residual selectors are constrained but not eliminated]\label{prop:selector-not-eliminated}
Dynamic martingale consistency constrains the drift of selected residual price functionals, while a
separate completion condition is needed to collapse the residual interval to a single selector.  If two positive
public-equivalent selector density families generate selected price functionals satisfying the drift
restriction and agree on public option coordinates, they remain observationally equivalent on the
public manifold while assigning different prices to residual latent-stress claims.
\end{proposition}

\begin{proof}
Agreement on public option coordinates follows from public equivalence.  The martingale
restriction says that each selected discounted price process has zero drift.  This gives dynamic
consistency for an already selected extension.  Completeness of the tradeable span requires a
separate condition.  If the two selector density families differ on a residual direction correlated with the latent
claim, the public non-identification calculation gives different selected prices for that claim
while preserving all public prices.
\end{proof}

\begin{definition}[Process-consistent public-equivalent kernel]\label{def:process-consistent-kernel}
A horizon-wise public-equivalent family \(Z_{t,T}\) is process-consistent if there is a positive
\(\QQ\)-martingale density process \(D^\theta=(D_s^\theta)_{s\ge0}\), \(D_0^\theta=1\), such
that
\[
        Z_{t,T}^\theta=\frac{D_T^\theta}{D_t^\theta}.
\]
It is dynamically public-equivalent on the admitted manifold if, for every date \(t\) and every
admitted public payoff \(X_a\) with maturity \(T(a)>t\),
\[
        \EE^\QQ\!\left[
        \frac{D_{T(a)}^\theta}{D_t^\theta}X_a
        \,\middle|\,\cH_t
        \right]
        =
        \EE^\QQ[X_a\mid\cH_t].
\]
\end{definition}

\begin{theorem}[Martingale residual densities preserve public prices]
\label{thm:martingale-residual-density}
Let \(D^\theta\) be a positive martingale density process satisfying the dynamic
public-equivalence identities in Definition~\ref{def:process-consistent-kernel}.  Then
\(D^\theta\) defines a dynamically public-equivalent extension of the public option manifold.
For any latent claim \(X\) whose discounted payoff has nonzero conditional covariance with
\(D_T^\theta/D_t^\theta\),
\[
        \EE^\QQ\!\left[
        \frac{D_T^\theta}{D_t^\theta}X
        \,\middle|\,\cH_t
        \right]
        \ne
        \EE^\QQ[X\mid\cH_t],
\]
the density process changes the selected residual price without changing admitted public option
prices.
\end{theorem}

\begin{proof}
The first statement is the defining preservation identity applied at every date and maturity.
For a latent claim, the price change is
\[
        \EE^\QQ\!\left[
        \left(\frac{D_T^\theta}{D_t^\theta}-1\right)X
        \middle|\cH_t
        \right].
\]
Because the density ratio has conditional mean one, this term is the conditional covariance of the
payoff with \(D_T^\theta/D_t^\theta\).  Nonzero conditional covariance is therefore exactly the
condition that the selected residual price differs from the baseline conditional price on a set of
positive probability.
\end{proof}

\begin{remark}
The drift restriction is the option-manifold analogue of an HJM restriction.  The state may be
parametrized by Volterra kernels, Perron spectral coordinates, rough maturity coordinates, and
residual selectors, but once these coordinates generate public option prices under \(\QQ\), their
joint drift and diffusion cannot be chosen independently.
\end{remark}

\section{Canonical finite benchmark model}\label{sec:finite-benchmark}

The preceding sections establish when a risk-neutral Volterra--Perron source state is recoverable
and which residual prices remain undetermined.  The finite benchmark fixes a closed chart: the
recovered state is a finite vector, the option formula is explicit, the no-arbitrage equation is a
state PDE, native Greeks are derivatives of the same state, and hedge incompleteness has a
computable residual.

The benchmark uses Bachelier coordinates.  The Bachelier choice is a normal-forward normalization
rather than an economic restriction on the theory.  It avoids basket-lognormal approximations and makes the synchronization
covariance object exact.  Lognormal, displaced, or stochastic-volatility variants can be built on
the same recovered state, but the finite benchmark is the smallest closed model in which the
Volterra--Perron quotient becomes a pricing equation.

\begin{definition}[Finite Bachelier--Volterra--Perron state]\label{def:finite-bvp-state}
Fix \(N\) underliers, a maturity horizon \(\bar\tau\), and an active set of carrier baskets
\(\mathcal B=\{w^1,\dots,w^M\}\subset\R^N\), including single names, indices, and any admitted
carrier portfolios.  A finite Bachelier--Volterra--Perron state at date \(t\) consists of
\[
        U_t
        =
        \bigl(F_t,\nu_t,\mathfrak{H}_t,G_t,\rho_t,e_t,\ell_t,\varpi_t,\theta_t\bigr),
\]
where \(F_t\in\R^N\) is the vector of forward levels, \(\nu_t\in(0,\infty)^N\) is a diagonal
idiosyncratic variance scale, \(\mathfrak{H}_t\in(-1/2,\infty)^Q\) is the finite rough maturity-scaling
vector, \(G_t\) is the finite directed route matrix, \((\rho_t,e_t,\ell_t)\) is the Perron
spectral triple with \(\ell_t(e_t)=1\), \(\varpi_t\) is the synchronized stress coordinate, and
\(\theta_t\) is an admissible residual selector.  The usual rough-volatility benchmark takes
\(\mathfrak{H}_{q,t}\in(0,1/2]\), but the finite normal-forward covariance formula only requires
\(2\mathfrak{H}_{q,t}+1>0\) for variance growth at the origin.

For \(0<\tau\le\bar\tau\), the conditional covariance matrix of forward changes is
\[
        \Gamma_t(\tau;U_t)
        =
        D_t(\tau)
        +
        \sum_{q=1}^Q
        \tau^{2\mathfrak{H}_{q,t}+1}B_{q,t}B_{q,t}^{\top}
        +
        a_\varpi(\tau;U_t)c_tc_t^{\top}
        +
        \Gamma_t^{\mathrm{res}}(\tau;\theta_t),
\]
where
\[
        D_t(\tau)=\tau\,\operatorname{diag}(\nu_t),\qquad
        c_t=\operatorname{diag}(\sqrt{\nu_t})\,e_t,
\]
\[
        a_\varpi(\tau;U_t)=\rho_t\exp(\varpi_t)\tau^{2\bar{\mathfrak{H}}_t+1},
        \qquad
        \bar{\mathfrak{H}}_t=\frac{\sum_q \omega_q \mathfrak{H}_{q,t}}{\sum_q\omega_q},
        \quad \omega_q>0,
\]
and the matrices \(B_{q,t}=B_q(G_t)\) are fixed differentiable functions of the directed route
coordinates.  The residual block is required to be symmetric and small enough that
\[
        \Gamma_t(\tau;U_t)\succeq0,\qquad
        \Gamma_t(\tau_2;U_t)-\Gamma_t(\tau_1;U_t)\succeq0
        \quad(0<\tau_1<\tau_2\le\bar\tau)
\]
on the admitted maturity interval.  Equivalently, in a smooth maturity chart
\(\partial_\tau\Gamma_t(\tau;U_t)\succeq0\).  The first condition gives slice-wise covariance
positivity; the second is the normal-forward calendar condition induced by convex order.  The
residual-free benchmark is obtained by setting \(\Gamma_t^{\mathrm{res}}\equiv0\).
\end{definition}

\begin{remark}[Gauge and positivity]\label{rem:finite-benchmark-gauge}
The gauge is fixed by \(e_t\) in the positive cone and \(\ell_t(e_t)=1\); the scalar size of the
Perron mode is carried by \(\exp(\varpi_t)\rho_t\).  This keeps the covariance contribution
positive while preserving a centered stress coordinate \(\varpi_t\).  Residual selectors do not have
to be positive semidefinite term by term; they must keep the total covariance matrix positive
semidefinite on the admitted panel.  This is the finite counterpart of the residual-envelope
positivity condition.
\end{remark}

\begin{remark}[Covariance-projected residual selectors]\label{rem:covariance-projected-selector}
The residual term \(\Gamma_t^{\mathrm{res}}(\tau;\theta_t)\) is a finite normal-forward
specialization of the general residual-kernel selector.  It represents only the part of the
selected residual pricing rule that can be expressed as an admissible covariance correction in a
Bachelier chart.  General public-equivalent density selectors may also change higher moments,
tail charges, or non-Gaussian residual prices without admitting a covariance-only
representation.  Thus the finite benchmark is a closed pricing model for the covariance-projected
branch of the theory, not a replacement for the full residual-density envelope.
\end{remark}

\begin{definition}[Benchmark option coordinates]\label{def:benchmark-option-coordinates}
For a basket \(w\in\mathcal B\), define
\[
        F_t^w=w^\top F_t,
        \qquad
        V_t^w(\tau;U_t)=w^\top\Gamma_t(\tau;U_t)w .
\]
Let \(P(t,T)>0\) be the public zero-coupon price and let \(\QQ^T\) denote the \(T\)-forward
measure on the admitted payoff span.  The vector \(F_t\) is the date-\(t\) vector of
\(T\)-forward levels for the contract maturity under consideration; the maturity argument is
suppressed only to lighten notation.  Let \(\Phi\) and \(\varphi\) denote the standard normal
distribution and density, and put
\[
        d_t^w(K,T)
        =
        \frac{F_t^w-K}{\sqrt{V_t^w(T-t;U_t)}} .
\]
For \(\kappa=1\) for a call and \(\kappa=-1\) for a put, the finite benchmark price is
\[
        C_t^\kappa(w,K,T;U_t)
        =
        P(t,T)
        \left[
        \kappa(F_t^w-K)\Phi(\kappa d_t^w)
        +
        \sqrt{V_t^w(T-t;U_t)}\,\varphi(d_t^w)
        \right],
\]
with the intrinsic-value convention when \(V_t^w=0\).  If discounting is deterministic or
\(\cH_t\)-measurable, this is equivalently the money-market formula with
\(P(t,T)=\EE^\QQ[B_t/B_T\mid\cH_t]\).
\end{definition}

\begin{theorem}[Canonical finite pricing formula]\label{thm:canonical-finite-pricing}
Suppose that, conditional on \(\cH_t\) under the \(T\)-forward measure \(\QQ^T\), every admitted
basket forward satisfies
\[
        F_{T}^{w}-F_t^w
        \sim
        N\!\left(0,V_t^w(T-t;U_t)\right)
\]
jointly across the basket family, with covariance induced by
\(\Gamma_t(T-t;U_t)\).  Then the price in
Definition~\ref{def:benchmark-option-coordinates} is the arbitrage-free date-\(t\) price of the
Bachelier call or put on \(w^\top F_T\).  The full finite option panel is the map
\[
        \mathfrak P_t^{\mathrm{fin}}:
        U_t
        \longmapsto
        \bigl(C_t^\kappa(w,K,T;U_t)\bigr)_{(w,K,T,\kappa)\in\mathcal I_t}.
\]
If the Jacobian of this map on the gauge-normalized source quotient has full column rank, the
finite panel locally recovers the benchmark state modulo the fixed gauge.
\end{theorem}

\begin{proof}
The conditional Bachelier formula is the expectation of
\((\kappa(w^\top F_T-K))_+\) under a one-dimensional normal law with mean \(F_t^w\) and variance
\(V_t^w\) under \(\QQ^T\), multiplied by \(P(t,T)\).  Joint Gaussianity ensures that all basket
prices come from one positive \(T\)-forward conditional law for the admitted forward vector, so
the finite panel is arbitrage-free on the basket payoff span.  The recovery statement is the
finite-dimensional inverse-function theorem applied to the gauge-normalized response map.
\end{proof}

\begin{theorem}[Recovered-state pricing equation in the finite benchmark]\label{thm:finite-benchmark-pde}
Assume the benchmark state \(U_t\in\R^d\) is Markov under \(\QQ\) on a recovery window,
\[
        \dd U_t=b^\QQ(t,U_t)\dd t+\Sigma^\QQ(t,U_t)\dd W_t^\QQ,
        \qquad
        a^\QQ=\Sigma^\QQ(\Sigma^\QQ)^\top,
\]
and let \(V(t,u)\) be the selected price of a calendar claim written in benchmark coordinates.
Then dynamic no-arbitrage requires
\[
        \partial_tV(t,u)
        +
        b^\QQ(t,u)^\top\nabla_u V(t,u)
        +
        \frac12
        \operatorname{Tr}\!\left(a^\QQ(t,u)D_u^2V(t,u)\right)
        -
        r_tV(t,u)
        =
        0
\]
on the recovery window, with terminal condition given by the claim payoff.  The second-order term
is the native state-volatility charge; it is the finite benchmark form of quotient gamma pricing.
\end{theorem}

\begin{proof}
This is Theorem~\ref{thm:martingale-restriction} in the finite Markov chart.  Applying It\^o's
formula to the discounted selected price \(e^{-\int_0^t r_s\dd s}V(t,U_t)\), and setting its
finite-variation drift equal to zero, gives the displayed equation.
\end{proof}

\begin{proposition}[Closed native Greeks in the Bachelier benchmark]\label{prop:closed-native-greeks}
Let \(z\) be any differentiable coordinate of the finite state \(U_t\).  For an admitted basket
option with \(V=V_t^w(\tau;U_t)>0\), \(F=F_t^w\), and \(d=(F-K)/\sqrt V\),
\[
        \partial_z C_t^\kappa
        =
        P(t,T)
        \left[
        \kappa\Phi(\kappa d)\,\partial_zF
        +
        \frac{\varphi(d)}{2\sqrt V}\,\partial_zV
        \right].
\]
In particular, for the Perron stress coordinate,
\[
        \partial_{\varpi}V_t^w(\tau)
        =
        a_\varpi(\tau;U_t)(w^\top c_t)^2
        +
        w^\top\partial_\varpi\Gamma_t^{\mathrm{res}}(\tau;\theta_t)w ,
\]
for the Perron spectral radius coordinate,
\[
        \partial_{\rho}V_t^w(\tau)
        =
        \exp(\varpi_t)\tau^{2\bar{\mathfrak{H}}_t+1}(w^\top c_t)^2
        +
        w^\top\partial_\rho\Gamma_t^{\mathrm{res}}(\tau;\theta_t)w ,
\]
and, for a roughness coordinate \(\mathfrak{H}_q\) appearing in a route term
\(\tau^{2\mathfrak{H}_q+1}B_qB_q^\top\),
\[
        \partial_{\mathfrak{H}_q}V_t^w(\tau)
        =
        2\log(\tau)\,
        \tau^{2\mathfrak{H}_q+1}(w^\top B_q)(B_q^\top w)
        +
        \partial_{\mathfrak{H}_q}a_\varpi(\tau;U_t)(w^\top c_t)^2
        +
        w^\top\partial_{\mathfrak{H}_q}\Gamma_t^{\mathrm{res}}(\tau;\theta_t)w .
\]
Directed-route, carrier, eigenvector, and residual-selector Greeks are obtained by the same
formula with \(\partial_zV=w^\top(\partial_z\Gamma_t)w\).
\end{proposition}

\begin{proof}
Differentiate the Bachelier formula.  The terms containing \(\partial_z d\) cancel, leaving
\(\partial_F C=P(t,T)\kappa\Phi(\kappa d)\) and
\(\partial_V C=P(t,T)\varphi(d)/(2\sqrt V)\).  The displayed native Greek formulas are the
corresponding derivatives of the finite covariance matrix in
Definition~\ref{def:finite-bvp-state}.  The zero-coupon factor is treated as public and not
differentiated in this native state chart; if a rate factor is included as a state coordinate, its
bond-price derivative adds a separate term.
\end{proof}

\begin{definition}[Finite public hedge residual]\label{def:finite-public-hedge-residual}
Let \(y=\mathfrak P_t^{\mathrm{fin}}(U_t)\in\R^M\) be the finite quote vector and
\(J=D_U\mathfrak P_t^{\mathrm{fin}}(U_t)\) its quotient Jacobian.  Let \(g_X\in(\R^d)^*\) be the
native Greek row of a claim \(X\).  A finite public hedge for \(X\) is a vector
\(h\in\R^M\) of quote-coordinate holdings.  Its first-order native exposure is \(h^\top J\).  The
minimum weighted first-order hedge residual is
\[
        \mathcal R_t^{\mathrm{hedge}}(X)
        =
        \inf_{h\in\R^M}
        \norm{g_X-h^\top J}_{\mathcal W_t^{-1}},
\]
where \(\mathcal W_t\) is a positive definite state-coordinate risk weight.
\end{definition}

\begin{theorem}[Finite hedge formula and incompleteness certificate]\label{thm:finite-hedge-formula}
Let \(J\) and \(g_X\) be as in Definition~\ref{def:finite-public-hedge-residual}.  If \(J\) has
full column rank, then \(\mathcal R_t^{\mathrm{hedge}}(X)=0\), and the minimum quote-norm hedge
solving \(J^\top h=g_X^\top\) is
\[
        h^\star
        =
        J(J^\top J)^{-1}g_X^\top .
\]
If the admitted quote span is restricted to a submatrix \(J_{\mathrm{tr}}\), then
\[
        \mathcal R_{t,\mathrm{tr}}^{\mathrm{hedge}}(X)
        =
        \norm{
        (I-\Pi_{\Ran(\mathcal W_t^{-1/2}J_{\mathrm{tr}}^\top)})
        \mathcal W_t^{-1/2}g_X^\top
        }_2 .
\]
Thus the finite benchmark separates recovered hedgeable claims \((\mathcal R=0)\), recovered
unhedgeable claims \((\mathcal R>0)\), and unrecovered claims
when \(J\) is rank deficient on the quotient.
\end{theorem}

\begin{proof}
The full-rank case is the adjoint Greek equation from
Definition~\ref{def:adjoint-quote-greek} with \(W=I\).  The minimum Euclidean-norm solution is
the displayed Moore--Penrose solution.  With a restricted tradeable quote span, the best hedge is the
orthogonal projection of the weighted native Greek onto the weighted row space generated by the
tradeable Jacobian.  The residual norm is exactly the distance to that row space.
\end{proof}

\begin{remark}[What the finite benchmark does and does not solve]\label{rem:finite-benchmark-scope}
The benchmark is the finite analogue of a canonical pricing equation.  It makes no claim that the
full listed-option manifold is Gaussian or Bachelier.  It gives a closed reference model in which the
Volterra--Perron state, option formula, state PDE, native Greeks, and public hedge residual are
all explicit.  Smile completion, jump tails, and liquidity effects are outside this finite
reference model.  The mathematical role of the benchmark is to show that the quotient
geometry can be compressed into a tractable no-arbitrage model without losing the residual and
carrier distinctions that drive the general theory.
\end{remark}

\section{Gaussian--Volterra source-family recovery}\label{sec:gaussian}

The Gaussian--Volterra branch is a source-family verification of the quotient recovery theorem.
The reduced form keeps the rank mechanism visible: temporal singular profiles, scalar seams, and a
common Perron-aligned option channel supply the response blocks.  Companion~A supplies the primitive
singular-value and scalar-seam estimates \cite{vidal2026gaussianseams}.

\begin{assumption}[Gaussian--Volterra recovery family]\label{ass:gv-family}
The risk-neutral latent volatility system is Gaussian conditional on \(\cH_t\).  Its directed
Volterra kernels have nonvanishing diagonal amplitudes and belong to a finite source chart whose
coordinates determine:
\begin{enumerate}[leftmargin=2em]
\item the leading singular-value profile of each temporal Volterra route;
\item the centered quadratic seam coefficient;
\item the common Perron-aligned option channel;
\item the index/single-name synchronization response matrix over an open maturity interval.
\end{enumerate}
The chart is gauge-normalized by fixing driver basis, Perron normalization, and common-channel
sign as in Definition~\ref{def:gauge-normalized-chart}.
\end{assumption}

\begin{assumption}[Gaussian response rank]\label{ass:gv-rank}
The Gaussian response blocks listed below have full rank after quotienting by gauge:
\begin{enumerate}[leftmargin=2em]
\item the temporal singular-value profile block that identifies rough maturity scaling;
\item the Volterra route block generated by cross-kernel singular coefficients;
\item the common-channel index/carrier response matrix;
\item in an augmented residual-instrument manifold, the residual-selector response on the admitted
residual-direction chart.
\end{enumerate}
The leading singular-value coefficients in the temporal routes are nonzero, and the seam defect is
lower order uniformly on the maturity interval.
\end{assumption}

\begin{theorem}[Reduced Gaussian--Volterra recovery]\label{thm:gv-recovery}
Under Assumptions~\ref{ass:gv-family} and~\ref{ass:gv-rank}, and assuming a simple
well-conditioned Perron eigenvalue and nondegenerate carrier Gram matrix, the Gaussian--Volterra
option manifold locally recovers the canonical risk-neutral Volterra--Perron state.  If the
residual-selector rank block is supplied only by augmented residual instruments, the selector is
recovered only in that augmented manifold; otherwise it remains a selected pricing extension.
\end{theorem}

\begin{proof}
The Gaussian source-family assumptions provide the chart \(U\) of
Assumption~\ref{ass:source-chart}.  The leading singular-value profile and scalar seam identify
the temporal and common-channel coordinates in the chart up to the gauge conventions.  The
Volterra route block identifies the directed cross-kernel coordinates, while the common-channel
matrix identifies the Perron synchronization coordinates.  The response-rank assumption gives
bounded left inverses for the diagonal blocks in Definition~\ref{def:block-source-chart}.
Finite dimensionality of the source chart turns injectivity of these blocks into bounded
inversion on their ranges.  A simple Perron eigenvalue with bounded reduced resolvent makes the
Perron eigenvalue and eigenvectors differentiable functions of the screened operator through the
simple-isolated-eigenvalue perturbation formulas.  The carrier Gram condition prevents collapse of
the carrier coordinates.
Hence the block-rank observability criterion applies, and
Theorem~\ref{thm:conditional-recovery} gives canonical recovery.
\end{proof}

\begin{remark}
The Gaussian companion analysis is used as a verification layer, not as a separate economic claim.
Its role is to justify that primitive Gaussian--Volterra kernels produce the singular-value seam
and common-channel option response required by the recovery theorem.
\end{remark}

\section{Residual kernel selection}\label{sec:residual}

Source recovery leaves residual pricing in place.  The recovered state fixes the canonical
coordinates available to a selector; it does not complete the public payoff module or enlarge the
tradeable span.

\begin{definition}[Conditional entropy residual constraint set]\label{def:entropy-residual-set}
Fix a date \(t\) and horizon \(T\).  Let
\[
        L_0\equiv1,\qquad L_1,\dots,L_m\in L^\infty(\QQ,\cK_t)
\]
be the latent loadings of the admitted public payoff constraints, and let
\(q_1,\dots,q_p\in L^\infty(\QQ,\cK_t)\) be residual scores whose exposure is to be selected.
For a bounded \(\cH_t\)-measurable target vector
\(\eta\in L^\infty(\QQ,\cH_t;\R^p)\), define
\[
        \cZ_t(\eta)
        :=
        \left\{
        Z>0:
        \begin{array}{l}
        \EE^\QQ[ZL_j\mid\cH_t]=\EE^\QQ[L_j\mid\cH_t],\quad j=0,\dots,m,\\
        \EE^\QQ[Zq_k\mid\cH_t]=\eta_k,\quad k=1,\dots,p
        \end{array}
        \right\}.
\]
The conditional entropy of \(Z\) is
\[
        H_t(Z):=\EE^\QQ[Z\log Z\mid\cH_t].
\]
\end{definition}

\begin{theorem}[Entropy-minimal residual density]\label{thm:entropy-minimal-density}
Assume \(\cZ_t(\eta)\) is nonempty, satisfies a conditional Slater condition, and has the
required conditional exponential moments.  Then there is a unique minimizer
\[
        Z_t^\star(\eta)
        =
        \argmin_{Z\in\cZ_t(\eta)}H_t(Z).
\]
It has conditional exponential form
\[
        Z_t^\star(\eta)
        =
        \exp\!\left(
        \alpha_0+\sum_{j=1}^m\alpha_jL_j+\sum_{k=1}^p\gamma_kq_k
        \right),
\]
where the multipliers are \(\cH_t\)-measurable and chosen to satisfy the constraints.  The
selected price
\[
        \Pi_t^{\star,\eta}(X)
        :=
        \EE^\QQ\!\left[
        Z_t^\star(\eta)X
        \,\middle|\,\cH_t
        \right]
\]
is linear, positive, public-equivalent on the admitted public constraints, and differentiable in
\(\eta\) wherever the conditional covariance matrix of the active constraints is nonsingular.
\end{theorem}

\begin{proof}
Conditional relative entropy is strictly convex on positive densities, and the moment constraints
form a convex affine set.  The conditional Slater condition gives the supporting Lagrange
multipliers and excludes boundary minimizers.  Minimizing the conditional Lagrangian pointwise in
the conditional problem yields the exponential density displayed above; strict convexity gives
uniqueness.  Differentiability in \(\eta\) follows from the conditional implicit-function theorem
applied to the moment equations, whose Jacobian is the conditional covariance matrix of the active
constraints under the tilted density.
\end{proof}

\begin{remark}[Residual envelopes and canonical selectors]
The residual envelope is the admissible price interval.  A linear basis selector chooses one point
in that interval.  The entropy-minimal selector supplies a canonical no-arbitrage choice once
public-equivalence constraints and residual target exposures are fixed.
\end{remark}

\begin{definition}[Terminal tradeable payoff span]\label{def:terminal-tradeable-payoff-span}
For a horizon \(T>t\), let \(\cS_{t,T}^{\mathrm{pay}}\subset L^2(\QQ,\cF_T)\) be the closed
subspace of discounted terminal payoffs generated by admissible public self-financing strategies
and by the listed or idealized public instruments admitted for trading between \(t\) and \(T\).
Orthogonal hedgeability statements are taken relative to this terminal payoff span.
\end{definition}

\begin{proposition}[Recovery does not imply tradeable completion]\label{prop:recovery-not-completion}
Fix \(T>t\).  Suppose the option manifold recovers the canonical state \(\Xi_t^\QQ\).  If
\((\cS_{t,T}^{\mathrm{pay}})^\perp\) contains a nonzero residual direction correlated with a claim
\(X\), and this direction can be chosen to annihilate the admitted public payoff span, then
recovering \(\Xi_t^\QQ\) does not make \(X\) tradeably replicable from
\(\cS_{t,T}^{\mathrm{pay}}\).  Residual kernel selection remains a pricing input.
\end{proposition}

\begin{proof}
Recovery makes the canonical coordinates measurable functions of the option manifold.  It does
not change the terminal tradeable payoff span \(\cS_{t,T}^{\mathrm{pay}}\).  If the orthogonal
residual of \(X\) relative to this span is nonzero, the Hilbert projection theorem leaves a strict
residual norm after the best tradeable projection.  If the residual direction also annihilates the
admitted public payoff span, Lemma~\ref{lem:residual-annihilator} constructs public-equivalent
density families that price that direction without changing public prices.  A selector is
therefore still needed.
\end{proof}

\begin{definition}[Dynamic residual selector]\label{def:residual-selector}
Let
\[
        \cR_{t,T}=\{R^1,\dots,R^m\}\subset L^\infty(\QQ,\cK_t)
\]
be a finite horizon-specific family of bounded date-\(t\) residual directions satisfying
\begin{gather*}
        \EE^\QQ[R^r\mid\cH_t]=0,
        \qquad r=1,\dots,m,\\
        \EE^\QQ[R^rX_a\mid\cH_t]=0
        \quad\text{for every admitted public payoff }X_a
        \text{ and every }r.
\end{gather*}
A dynamic residual selector is an \(\cH_t\vee\sigma(\Xi_t^\QQ)\)-measurable vector
\[
        \theta_t=(\theta_{1,t},\dots,\theta_{m,t})
\]
such that
\[
        1+\sum_{r=1}^m\theta_{r,t}R^r>0 .
\]
The associated density family is horizon-wise; it need not be a dynamically consistent density
process unless the martingale consistency conditions in Section~\ref{sec:martingale} are imposed.
\end{definition}

\begin{definition}[Dynamic option-manifold selected residual price]\label{def:dynamic-selected-price}
For \(X\in L^2(\QQ,\cF_T)\), define
\[
        \Pi_t^\theta(X)
        :=
        \EE^\QQ[X\mid\cH_t]
        +
        \sum_{r=1}^m
        \theta_{r,t}\,
        \EE^\QQ[L_{t,T}^\QQ(X)R^r\mid\cH_t].
\]
This is the linear price induced by the selected density
\[
        Z_{t,S}^{\theta}=1+\sum_{r=1}^m\theta_{r,t}R^r .
\]
\end{definition}

\begin{definition}[Residual price envelope]\label{def:residual-envelope}
For a conditionally bounded residual family \(\mathfrak R_{t,T}\) and admissible radius
\(\bar\theta_t(R)\norm{R}_\infty<1\), define the upper residual envelope by
\[
        \overline\Pi_t(X)
        :=
        \EE^\QQ[X\mid\cH_t]
        +
        \operatorname*{ess\,sup}_{R\in\mathfrak R_{t,T},\,|\eta|\le\bar\theta_t(R)}
        \eta\,\EE^\QQ[L_{t,T}^\QQ(X)R\mid\cH_t],
\]
whenever the conditional essential supremum is finite.  The lower envelope is defined by the
corresponding essential infimum.
\end{definition}

\begin{theorem}[Selected prices lie in the residual envelope]\label{thm:selected-envelope}
Suppose \(\cR_{t,T}=\{R^1,\dots,R^m\}\subset\mathfrak R_{t,T}\) is conditionally bounded in
\(L^\infty\) and the selector satisfies
\[
        |\theta_{r,t}|\le\bar\theta_t(R^r),
        \qquad
        \left\|\sum_{r=1}^m\theta_{r,t}R^r\right\|_\infty<1.
\]
Then \(\Pi_t^\theta(X)\) lies inside the public-equivalent residual price envelope generated by
the density families \(Z_{t,S}^{\eta,R}=1+\eta R\), \(|\eta|\le\bar\theta_t(R)\),
\(R\in\mathfrak R_{t,T}\), provided the selected aggregate residual
\(\sum_r\theta_{r,t}R^r\) is admitted by the envelope family.
\end{theorem}

\begin{proof}
For each \(R\), the density family \(Z_{t,S}^{\eta,R}=1+\eta R\) is positive and
public-equivalent by Lemma~\ref{lem:residual-annihilator}.  Its conditional price of \(X\) is
\[
        \EE^\QQ[X\mid\cH_t]
        +
        \eta\EE^\QQ[L_{t,T}^\QQ(X)R\mid\cH_t].
\]
The selected price is the conditional expectation under the single positive density family
\[
        1+\sum_r\theta_{r,t}R^r.
\]
If the selected aggregate residual belongs to the admissible envelope family, this linear price
is one member of the support family over which the envelope takes the essential supremum and
infimum.  Hence it lies between the lower and upper residual envelopes.
\end{proof}

\begin{definition}[Residual Gram coordinates]\label{def:residual-gram-coordinates}
For a finite residual basis \(R=(R^1,\dots,R^m)^\top\), define the conditional residual Gram
matrix and payoff-residual vector by
\[
        G_t^R:=\EE^\QQ[RR^\top\mid\cH_t],
        \qquad
        g_t^X:=
        \EE^\QQ[L_{t,T}^\QQ(X)R\mid\cH_t].
\]
For a selector vector \(\theta_t\), the selected residual price can be written
\[
        \Pi_t^\theta(X)
        =
        \EE^\QQ[X\mid\cH_t]+\theta_t^\top g_t^X,
\]
provided the selected density is positive and public-equivalent.
\end{definition}

\begin{theorem}[Ellipsoidal residual envelope]\label{thm:ellipsoidal-residual-envelope}
Assume \(G_t^R\) is positive definite on the active residual span.  For the conditional selector
budget
\[
        \Theta_t^0(\rho):=
        \{\theta:\theta^\top G_t^R\theta\le\rho_t^2\},
\]
the unconstrained linear residual support function is
\[
        \sup_{\theta\in\Theta_t^0(\rho)}
        \theta^\top g_t^X
        =
        \rho_t\sqrt{(g_t^X)^\top(G_t^R)^{-1}g_t^X}.
\]
The corresponding lower endpoint is obtained by changing the sign.  If the selector budget is
intersected with the positivity domain of the selected densities, the same formula remains an
upper bound, and equality holds whenever the Cauchy--Schwarz optimizer lies in that positivity
domain.  In particular, equality holds for every claim in the active residual span under the
small-ball condition
\[
        \rho_t
        \sqrt{
        \operatorname*{ess\,sup}
        R^\top(G_t^R)^{-1}R
        }
        <1 .
\]
Therefore, inside the unconstrained ellipsoidal selector class, and inside the constrained class
whenever the preceding positivity condition applies,
\[
        \Pi_t^\theta(X)
        \in
        \EE^\QQ[X\mid\cH_t]
        \pm
        \rho_t\sqrt{(g_t^X)^\top(G_t^R)^{-1}g_t^X}.
\]
\end{theorem}

\begin{proof}
This is Cauchy--Schwarz in the conditional inner product
\(\ip{a}{b}_{G}:=a^\top G_t^Rb\):
\[
        \theta^\top g
        =
        \ip{\theta}{(G_t^R)^{-1}g}_{G}
        \le
        (\theta^\top G_t^R\theta)^{1/2}
        (g^\top(G_t^R)^{-1}g)^{1/2}.
\]
Equality is achieved in the direction \((G_t^R)^{-1}g_t^X\), unless restricted by the positivity
domain.  Intersecting with the positivity domain can only reduce the support function.  The
small-ball condition implies \(1+\theta^\top R>0\) almost surely for every
\(\theta^\top G_t^R\theta\le\rho_t^2\), again by Cauchy--Schwarz, and hence the unconstrained
optimizer is admissible.  The lower endpoint follows by the opposite direction.
\end{proof}

\begin{corollary}[Box residual envelope]\label{cor:box-residual-envelope}
For the unconstrained box selector budget \(|\theta_r|\le\bar\theta_r\),
\[
        \sup_{|\theta_r|\le\bar\theta_r}\theta^\top g_t^X
        =
        \sum_{r=1}^m \bar\theta_r |g_{t,r}^X|,
\]
and the lower endpoint is obtained by changing the sign.  If the box is intersected with the
positivity domain of the selected densities, the same expression is an upper bound and equality
holds when the coordinatewise sign optimizer is positive.
\end{corollary}

\begin{proof}
The support function of the box \(\prod_r[-\bar\theta_r,\bar\theta_r]\) in direction \(g_t^X\) is
the sum of the coordinate support functions,
\[
        \sum_{r=1}^m \sup_{|\theta_r|\le\bar\theta_r}\theta_rg_{t,r}^X
        =
        \sum_{r=1}^m\bar\theta_r|g_{t,r}^X|.
\]
The lower endpoint is the support function in direction \(-g_t^X\).  Intersecting the box with the
positivity domain only removes admissible selectors, so the expression remains an upper bound and
is attained whenever the coordinatewise sign optimizer lies in that domain.
\end{proof}


\begin{proposition}[Adaptedness of selected prices]\label{prop:selected-adapted}
If \(\Xi_t^\QQ\), \(\cR_{t,T}\), and \(\theta_t\) are
\(\cH_t\vee\sigma(\Xi_t^\QQ)\)-measurable and the recovered canonical state is
\(\cH_t\)-measurable through the option manifold, then \(\Pi_t^\theta(X)\) is \(\cH_t\)-measurable.
\end{proposition}

\begin{proof}
The public price is \(\cH_t\)-measurable.  Under recovery, the canonical coordinates entering
\(\theta_t\) are measurable functions of \(\cM_t^{\mathrm{opt}}\), hence \(\cH_t\)-measurable.
Conditional pairings with residual directions are conditional expectations given \(\cH_t\), hence
\(\cH_t\)-measurable.  The finite selected sum is therefore \(\cH_t\)-measurable under the stated
finiteness assumption.
\end{proof}

\section{Native stress Greeks}\label{sec:greeks}

Classical option Greeks differentiate prices with respect to public surface coordinates or
underlying price variables.  Volterra--Perron Greeks differentiate selected prices with respect to
recovered quotient coordinates.  Once the source quotient is recovered, the primary sensitivity is
a cotangent vector on that quotient.  Public quote Greeks are adjoint representations of the same
object, not separate primitives.

\begin{definition}[Native stress Greeks]\label{def:native-greeks}
Let \(\Pi_t^\theta(X;\Theta)\) be the selected price at state \(\Theta\).  If the indicated
directional derivatives exist, define
\[
        \Delta_\varpi X
        :=
        D\Pi_t^\theta(X;\Theta)[\delta\varpi],
        \qquad
        \Delta_\rho X
        :=
        D\Pi_t^\theta(X;\Theta)[\delta\rho],
\]
\[
        \Delta_{e_i} X
        :=
        D\Pi_t^\theta(X;\Theta)[\delta e_i],
        \qquad
        \Delta_{\mathfrak H} X
        :=
        D\Pi_t^\theta(X;\Theta)[\delta\mathfrak H],
\]
\[
        \Delta_{G_{ij}} X
        :=
        D\Pi_t^\theta(X;\Theta)[\delta G_{ij}],
        \qquad
        \Delta_\theta X
        :=
        D\Pi_t^\theta(X;\Theta)[\delta\theta].
\]
These are, respectively, Perron-stress, spectral-radius, Perron-eigenvector, roughness,
directed-contagion, and residual-selector Greeks.  If carrier weights are separate state
coordinates, their derivatives define additional carrier-weight Greeks.
\end{definition}

\begin{assumption}[Differentiable selected pricing]\label{ass:differentiable-pricing}
In a neighborhood of the recovered state, the map
\[
        \Theta\mapsto \Pi_t^\theta(X;\Theta)
\]
is Fr\'echet differentiable on a chart transverse to gauge directions.  The residual selector is
differentiable in the same chart.
\end{assumption}

\begin{theorem}[Existence and gauge invariance of native Greeks]\label{thm:native-greeks}
Under Assumption~\ref{ass:differentiable-pricing}, native stress Greeks exist as continuous linear
functionals on the canonical quotient tangent space.  If two chart perturbations differ by a gauge
direction, they have the same native Greek.
\end{theorem}

\begin{proof}
Fr\'echet differentiability gives a continuous linear derivative on the chart tangent space.  The
selected price is a function of the gauge class of the state.  Gauge transformations leave latent
loadings, public option prices, and selected residual prices unchanged.  Hence the derivative
vanishes on gauge directions and factors through the quotient tangent space.  The quotient
functional is the native Greek.
\end{proof}

\begin{theorem}[Recovered-manifold Greek]\label{thm:recovered-manifold-greek}
Suppose the recovery theorem holds in a \(C^1\) embedded quotient chart, for instance because the
quotient source chart is finite dimensional with full column rank or because
\(D\bar F\) has closed complemented range and locally constant split rank.  Suppose also that
\(\Pi_t^\theta(X;\Theta)\) is differentiable in the canonical state.  On the recovery
neighborhood, the selected price can be written as a function of the public option manifold,
\[
        \widehat\Pi_t^\theta(X;y)
        :=
        \Pi_t^\theta\!\left(X;\bar F^{-1}(y)\right).
\]
Its first derivative is
\[
        D_y\widehat\Pi_t^\theta(X;y)
        =
        D_\Theta\Pi_t^\theta(X;\Theta)\circ D_y\bar F^{-1}(y),
        \qquad \Theta=\bar F^{-1}(y).
\]
Thus a public surface perturbation has a latent-stress decomposition through the recovered
Volterra--Perron coordinates.
\end{theorem}

\begin{proof}
The local inverse \(\bar F^{-1}\) is differentiable by the recovery theorem and the inverse
function theorem.  The displayed identity is the chain rule applied to the composition
\(\Pi_t^\theta\circ\bar F^{-1}\).
\end{proof}

\begin{definition}[Counterfactual Volterra--Perron shock]\label{def:counterfactual-vp-shock}
At a recovered state \(\Xi_t^\QQ\), a counterfactual Volterra--Perron shock is a
quotient tangent vector
\[
        h\in T_{[\Xi_t^\QQ]}(\mathfrak S_t/\!\sim_g)
\]
that satisfies the linearized no-arbitrage tangent conditions of the public option manifold.  Its
first-order public surface displacement is
\[
        \delta_h\cM_t^{\mathrm{opt}}
        :=
        D\mathfrak P_t([\Xi_t^\QQ])h,
\]
and its selected latent-price displacement for claim \(X\) is
\[
        \mathfrak G_h(X)
        :=
        D_\Xi\Pi_t^\theta(X;\Xi_t^\QQ)[h].
\]
Canonical examples are Perron-amplification shocks, carrier-thinning shocks, roughness shocks,
directed-route shocks, and residual-selector shocks.
\end{definition}

\begin{proposition}[Gauge-invariant stress calculus]\label{prop:gauge-invariant-stress-calculus}
If the public pricing map and selected latent price are well defined on the quotient state, then
both maps
\[
        h\mapsto D\mathfrak P_t([\Xi_t^\QQ])h,
        \qquad
        h\mapsto D_\Xi\Pi_t^\theta(X;\Xi_t^\QQ)[h]
\]
are independent of the representative chosen for \(\Xi_t^\QQ\).  Before quotienting, both
derivatives vanish on pure gauge directions.
\end{proposition}

\begin{proof}
Maps defined on a quotient are constant on gauge orbits.  Their derivatives annihilate tangent
vectors to those orbits and descend to the quotient tangent space.  The two displayed
derivatives are exactly the descended derivatives.
\end{proof}

\begin{definition}[Adjoint quote Greek]\label{def:adjoint-quote-greek}
Let \(F_\tau:U\subset\R^d\to\R^M\) be a finite-panel observation map with full column-rank
Jacobian \(J=DF_\tau(u)\).  Let \(p_X:=D_u\Pi_t^\theta(X;u)\in(\R^d)^*\) be a native Greek row
vector, and let \(W\succ0\) be the finite-panel quote metric.  A quote-space adjoint Greek is any
quote covector \(q_X\in\R^M\) satisfying
\[
        J^\top q_X=p_X^\top .
\]
The minimum \(W\)-norm adjoint Greek, where \(\|q\|_W^2=q^\top Wq\), is
\[
        q_X^{*,W}:=
        W^{-1}J(J^\top W^{-1}J)^{-1}p_X^\top .
\]
\end{definition}

\begin{theorem}[Quote perturbation representation]\label{thm:quote-perturbation-greek}
For any quote perturbation \(\delta y\in\Ran J\),
\[
        D_y\bigl(\Pi_t^\theta\circ F_\tau^{-1}\bigr)[\delta y]
        =
        (q_X^{*,W})^\top\delta y .
\]
If \(\delta y=J\delta u\), then this equals \(p_X\delta u\).
\end{theorem}

\begin{proof}
The formula for \(q_X^{*,W}\) is the Lagrange multiplier solution of minimizing
\(q^\top Wq\) subject to \(J^\top q=p_X^\top\).  Hence
\(J^\top q_X^{*,W}=p_X^\top\).  For tangent perturbations \(\delta y=J\delta u\),
\[
        (q_X^{*,W})^\top\delta y
        =
        (q_X^{*,W})^\top J\delta u
        =
        (J^\top q_X^{*,W})^\top\delta u
        =
        p_X\delta u.
\]
This is exactly the native-state derivative of the selected price.
\end{proof}

\begin{corollary}[Native and public Greeks are adjoint images]\label{cor:native-public-greek-correspondence}
Under the recovery hypotheses of Theorem~\ref{thm:spine}, every differentiable selected claim
price has a canonical native Greek on the source quotient.  On any finite full-rank public panel
with quote metric \(W\), its minimum \(W\)-norm public quote Greek is the adjoint solution
\[
        q_X^{*,W}
        =
        W^{-1}J(J^\top W^{-1}J)^{-1}p_X^\top .
\]
Thus Black--Scholes Greeks and Greeks computed in public quote coordinates are adjoint images of
quotient sensitivities to roughness, directed routes, Perron amplification, carrier absorption,
phase transition scores, and residual selection.
\end{corollary}

\begin{proof}
The native Greek is the quotient derivative from Theorem~\ref{thm:native-greeks}.  On a recovered
finite panel, \(J\) maps quotient source perturbations to public quote perturbations.  The
weighted adjoint equation is exactly the representation identity in
Theorem~\ref{thm:quote-perturbation-greek}.
\end{proof}

\begin{remark}[Finite Greek suites]\label{rem:finite-greek-suites}
In a finite implementation one may report a named suite of native Greek directions by choosing
source-chart tangents for Perron stress, spectral amplification, rough maturity scaling, directed
routes, carrier weights, phase scores, and residual-selector shifts.  The weighted adjoint
equation then turns each declared tangent into a quote-space functional and records any component
outside the admitted quote tangent.  Such a suite is a certificate-layer object.  It measures how
the finite option panel represents selected quotient sensitivities, without changing the quotient
recovery theorem.
\end{remark}

\begin{theorem}[First-order overidentifying restrictions]\label{thm:overidentifying-restrictions}
Let \(F_\tau:\R^d\to\R^M\) be a finite-panel source response with \(M>d\) and full
column-rank Jacobian \(J=DF_\tau(u_0)\).  Let
\[
        \cT_{y_0}:=\Ran J,\qquad
        \cN_{y_0}:=\Ker J^\top,
        \qquad y_0=F_\tau(u_0).
\]
If a perturbed quote panel \(y_\eps=y_0+\eps h+o(\eps)\) remains on the model image to first
order, then
\[
        n^\top h=0
        \qquad\text{for every }n\in\cN_{y_0}.
\]
Equivalently, \((I-P_{\cT})h=0\), where \(P_{\cT}\) is Euclidean projection onto
\(\cT_{y_0}\).  Thus a finite panel with more quote coordinates than source coordinates imposes
\(M-d\) first-order normal-space restrictions at a full-rank point.
\end{theorem}

\begin{proof}
If \(y_\eps=F_\tau(u_0+\eps v+o(\eps))\), then
\[
        h=Jv\in\Ran J.
\]
Membership in \(\Ran J\) is equivalent to orthogonality to \(\Ker J^\top\), and also to
\((I-P_{\cT})h=0\).
\end{proof}

\begin{proposition}[Perron spectral Greek]\label{prop:spectral-greek}
Suppose the screened synchronization operator \(A^\QQ\) has a simple Perron eigenvalue
\(\rho^\QQ\) with normalized left and right eigenvectors \(\ell^\QQ(e^\QQ)=1\).  For a bounded
operator perturbation \(\delta A\),
\[
        \delta\rho^\QQ
        =
        \ell^\QQ(\delta A\,e^\QQ).
\]
Consequently, any selected price differentiable in \(\rho^\QQ\) has spectral sensitivity
\[
        D\Pi_t^\theta[\delta A]
        =
        \partial_{\rho}\Pi_t^\theta\,
        \ell^\QQ(\delta A\,e^\QQ)
        +
        \text{eigenvector and residual-selector terms}.
\]
\end{proposition}

\begin{proof}
Differentiate \(Ae=\rho e\) and apply the normalized left eigenfunctional \(\ell\):
\[
        \ell(\delta A e)+\ell(A\delta e)
        =
        \delta\rho\,\ell(e)+\rho\ell(\delta e).
\]
Since \(\ell A=\rho\ell\) and \(\ell(e)=1\), the \(\ell(\delta e)\) terms cancel and
\(\delta\rho=\ell(\delta A e)\).  Differentiating the selected price through the spectral chart
and applying the chain rule gives the second display.
\end{proof}

\begin{proposition}[Perron eigenvector derivatives]\label{prop:perron-eigenvector-derivatives}
Let \(A e=\rho e\), \(\ell A=\rho\ell\), and \(\ell(e)=1\), with \(\rho\) simple and
isolated.  Let \(S\) be the reduced resolvent of \(A-\rho I\) on the complement of
\(e\ell\).  For a perturbation \(\delta A\) preserving the chosen normalization,
\[
        \delta e=-S(\delta A-\delta\rho I)e,
        \qquad
        \delta\ell=-\ell(\delta A-\delta\rho I)S .
\]
Consequently,
\[
        |\delta\rho|\le\norm{\ell}\,\norm{e}\,\norm{\delta A},
        \qquad
        \norm{\delta e}+\norm{\delta\ell}
        \le
        C\kappa_P\norm{\delta A},
\]
where \(\kappa_P:=\norm{S}\) is the reduced-resolvent norm and \(C\) depends only on the local
normalization.  On a finite-dimensional Perron stratum with uniformly bounded eigenprojection
conditioning, \(\kappa_P\le C'/\gamma_P\); without such conditioning, \(\kappa_P\) is the correct
stability constant.
\end{proposition}

\begin{proof}
Differentiating \(Ae=\rho e\) gives
\[
        (A-\rho I)\delta e=-(\delta A-\delta\rho I)e .
\]
The normalization fixes the component of \(\delta e\) in the Perron direction.  Inverting
\(A-\rho I\) on the complement of \(e\ell\) gives the displayed formula for \(\delta e\).  The
left-eigenvector identity follows in the same way from differentiating \(\ell A=\rho\ell\) and
imposing \(\delta\ell(e)+\ell(\delta e)=0\).  The norm bound is controlled by the reduced-resolvent
norm \(\kappa_P=\norm{S}\); the gap estimate is a conditioned finite-dimensional corollary, not the
primitive constant.
\end{proof}

\section{Synchronization relative value and public-neutral hedging}\label{sec:rv}

Let \(I\) be an index with weights \(\omega_i\).  The index variance decomposition separates
single-name variance from synchronized covariance.  In recovered coordinates, the same split
becomes a selected pricing identity.

\begin{definition}[Recoverability, tradeability, and hedgeability]\label{def:recoverability-tradeability-hedgeability}
Fix \(t<T\).  The state is recoverable when the public response is locally invertible on the
Volterra--Perron source quotient.  A claim is tradeable if its payoff lies in the admitted
listed or idealized public instrument span at horizon \(T\).  A claim is hedgeable if its centered
payoff lies in the closed gain space generated by admissible self-financing strategies between
\(t\) and \(T\).
\end{definition}

\begin{theorem}[Recovery does not imply tradeable hedgeability]\label{thm:recovery-not-trading}
Assume the option manifold recovers the canonical state \(\Xi_t^\QQ\).  Then the following
hold.
\begin{enumerate}[leftmargin=2em]
\item recovery of \(\Xi_t^\QQ\) does not imply that a Perron stress claim or residual-stress claim
belongs to the tradeable public span;
\item residual-kernel selection can price latent-stress claims while leaving them unspanned by
tradeable instruments;
\item a public-neutral book can have zero exposure to selected public surface directions and
nonzero exposure to a recovered synchronization direction, but this exposure is hedgeable only
under an additional traded-spanning assumption.
\end{enumerate}
Consequently,
\[
        \text{recoverability}
        \ne
        \text{tradeability}
        \ne
        \text{hedgeability}.
\]
The first is a quotient-inversion property, the second is an instrument-spanning property, and the
third is a dynamic gain-space property.
\end{theorem}

\begin{proof}
The recovery theorem identifies a canonical source coordinate from public option prices.  It does
not enlarge the closed terminal payoff span generated by tradeable instruments, nor does it alter
the dynamic gain space of self-financing strategies.  If a latent-stress payoff has a nonzero
orthogonal component relative to the terminal tradeable span, the Hilbert projection theorem leaves
a strictly positive hedge residual after the best tradeable projection.  Residual selectors change
prices through public-annihilating directions, but public equivalence only preserves admitted
public prices; it does not construct a gain representation.  Recoverability is therefore quotient
inversion, tradeability is payoff-spanning, and hedgeability is membership in the dynamic gain
space.
\end{proof}

\begin{remark}[Carrier resolution]\label{rem:carrier-resolution}
In finite option manifolds, tradeability also depends on the resolution of the admitted carrier
span.  A coarse underlier carrier library, an index/sector carrier library, a single-name
option-coordinate carrier library, and a regularized coordinate book need not span the same
synchronization directions.  Carrier absorption is a finite-span property of the chosen public
instrument system, not a consequence of latent-state recoverability alone.
\end{remark}

\begin{definition}[Carrier-replicated value]\label{def:carrier-replicated-value}
For an index option or variance claim \(X_I\), let \(X_{\mathrm{car}}\) be the carrier-replicated
single-name claim formed from the admitted carrier weights and matching maturity coordinates.  Let
\(\varpi_{t,T}^\QQ\) be the horizon-lifted synchronized stress direction.  In the scalar case define
\[
        X_{\mathrm{sync}}
        :=
        \frac{
        \EE^\QQ[(X_I-X_{\mathrm{car}})\varpi_{t,T}^\QQ\mid\cH_t]
        }{
        \EE^\QQ[(\varpi_{t,T}^\QQ)^2\mid\cH_t]
        }
        \varpi_{t,T}^\QQ,
\]
with the convention that the coefficient is zero if the conditional denominator vanishes.  For a
multidimensional stress subspace, replace the scalar ratio by the corresponding conditional Gram
inverse.  The tradeable basis error is
\[
        B_t(X_I):=
        X_I-X_{\mathrm{car}}-X_{\mathrm{sync}},
\]
so the synchronized component is canonical once the carrier book and stress subspace are fixed.
\end{definition}

\begin{definition}[Synchronization premium and residual charge]\label{def:sync-premium-charge}
The synchronization premium is
\[
        \mathrm{SynPrem}_t(X_I)
        :=
        \EE^\QQ[X_{\mathrm{sync}}\mid\cH_t].
\]
The residual compression charge selected by \(\theta\) is
\[
        \mathrm{ResCharge}_t^\theta(X_I)
        :=
        \Pi_t^\theta(X_{\mathrm{sync}})
        -
        \EE^\QQ[X_{\mathrm{sync}}\mid\cH_t].
\]
\end{definition}

\begin{remark}[Public shadow versus pure residual synchronization]\label{rem:sync-public-shadow}
The stress direction \(\varpi_{t,T}^\QQ\) is centered by convention, but it need not be
orthogonal to \(\cH_t\).  Its public compression may therefore carry a nonzero synchronization
premium.  If one chooses instead the purely public-residual stress direction
\[
        \varpi_{t,T}^\QQ-\EE^\QQ[\varpi_{t,T}^\QQ\mid\cH_t],
\]
then the public synchronization premium of that residual component vanishes and the selected
price contribution appears through the residual charge.  The identity below allows both
conventions; it separates the public shadow of synchronized stress from the residual
public-equivalent pricing selection.
\end{remark}

\begin{theorem}[No-arbitrage synchronization relative-value identity]\label{thm:rv-identity}
Suppose the selected price on the span of
\(\{X_I,X_{\mathrm{car}},X_{\mathrm{sync}},B_t(X_I)\}\) is induced by a fixed positive
public-equivalent density family \(Z^\theta_{t,\cdot}\).  Then the selected valuation is a
no-arbitrage extension of the public option manifold on this span.  For every square-integrable index claim
\(X_I\) with decomposition
\[
        X_I=X_{\mathrm{car}}+X_{\mathrm{sync}}+B_t(X_I),
\]
the selected price satisfies
\[
        \Pi_t^\theta(X_I)-\Pi_t^\theta(X_{\mathrm{car}})
        =
        \mathrm{SynPrem}_t(X_I)
        +
        \mathrm{ResCharge}_t^\theta(X_I)
        +
        \Pi_t^\theta(B_t(X_I)).
\]
\end{theorem}

\begin{proof}
Because \(Z^\theta_{t,\cdot}\) is positive and public-equivalent, conditional expectation under
the horizon component \(Z^\theta_{t,T}\) is a positive linear extension of the public price system
and preserves all date-\(t\) public option prices.  The extension is therefore arbitrage-free on
the finite span under consideration.  Applying this linear functional to
\(X_I=X_{\mathrm{car}}+X_{\mathrm{sync}}+B_t(X_I)\) gives the decomposition of
\(\Pi_t^\theta(X_I)\).  Subtracting \(\Pi_t^\theta(X_{\mathrm{car}})\) leaves the selected price of
\(X_{\mathrm{sync}}\) and the basis term.  Splitting \(\Pi_t^\theta(X_{\mathrm{sync}})\) into its
public conditional expectation and selected residual increment gives \(\mathrm{SynPrem}_t\) and
\(\mathrm{ResCharge}_t^\theta\).
\end{proof}

\begin{definition}[Model public-neutral book]\label{def:model-public-neutral-book}
Let \(\mathsf H_t^{\mathrm{pub}}\subseteq L^2(\QQ,\cH_t)\) be the finite public surface span used
for neutrality.  A listed or idealized book \(Y_t\in L^2(\QQ)\) is model public-neutral if
\[
        C_{\mathsf H_t^{\mathrm{pub}}}^\QQ Y_t=0.
\]
It is synchronization-exposed if
\[
        \ip{Y_t}{\varpi_t^\QQ}_{L^2(\QQ)}\ne0
\]
after projecting \(\varpi_t^\QQ\) into the claim horizon.
\end{definition}

\begin{theorem}[Public-neutral residual hedge]\label{thm:public-neutral-hedge}
Let \(X_I\) be an index claim and let
\[
        \cB_t=Y_0+\cV_t
\]
be a nonempty closed affine class of books satisfying the public-neutrality constraint, with
\(\cV_t\) a closed linear subspace of \(L^2(\QQ)\).  Then the minimum-residual book
\[
        Y_t^\star
        =
        \argmin_{Y\in\cB_t}
        \norm{X_{\mathrm{sync}}-Y}_{2,\QQ}
\]
exists and is unique.  The residual
\[
        X_{\mathrm{sync}}-Y_t^\star
\]
is orthogonal to feasible public-neutral perturbations of the book.
If additional exposure caps make \(\cB_t\) closed and convex rather than affine, the projection
still exists and is unique, and the orthogonality condition is replaced by the usual variational
inequality.
\end{theorem}

\begin{proof}
The Hilbert projection theorem gives a unique metric projection of \(X_{\mathrm{sync}}\) onto the
nonempty closed affine subspace \(\cB_t\).  For an affine feasible class, first-order optimality
is exactly
\[
        X_{\mathrm{sync}}-Y_t^\star\perp \cV_t .
\]
For a closed convex feasible class, the same projection theorem gives existence and uniqueness,
and first-order optimality becomes the variational inequality
\[
        \ip{X_{\mathrm{sync}}-Y_t^\star}{Y-Y_t^\star}_{L^2(\QQ)}\le0,
        \qquad Y\in\cB_t .
\]
\end{proof}

\section{Equivalent-martingale-measure specialization}\label{sec:emm}

The preceding theory works with adapted arbitrage-free option manifolds.  This section records one
equivalent-martingale-measure specialization in a continuous-time risk-neutral stock-volatility
system.

\begin{assumption}[Reduced risk-neutral stock-volatility system]\label{ass:emm}
For \(i=1,\dots,N\), under \(\QQ\),
\[
        \frac{\dd S_i(t)}{S_i(t)}
        =
        r_t\,\dd t+\sqrt{V_i(t)}\,\dd W_i^\QQ(t),
\]
where \(V_i(t)=\chi_i(Y_i(t))\) for a locally Lipschitz map
\(\chi_i:\R\to[0,\infty)\), such as an exponential or softplus transform.  The stress factor
vector satisfies the Volterra equation
\[
        Y_i(t)
        =
        y_i^0(t)
        +
        \sum_{a=1}^m\int_0^t K_{ia}^\QQ(t-s)\,\dd B_a^\QQ(s)
        +
        \sum_{j\ne i}\int_0^t G_{ij}^\QQ(t-s)\phi_j(Y_j(s))\,\dd s.
\]
The kernels are square-integrable on compact horizons, the functions \(\phi_j\) are Lipschitz
with linear growth, and \(Y\) admits a weak solution with
\(\EE^\QQ\int_0^T V_i(t)\dd t<\infty\).  The Brownian vectors \((W^\QQ,B^\QQ)\) have a specified
joint covariance; in the reduced examples below this covariance is fixed as part of the source
chart.
\end{assumption}

\begin{theorem}[Risk-neutral option manifold under an EMM]\label{thm:emm-specialization}
Under Assumption~\ref{ass:emm}, discounted asset prices are local \(\QQ\)-martingales.  If they
are true martingales on \([0,T]\), then European option prices
\[
        C_t(a)
        =
        \EE^\QQ\!\left[
        e^{-\int_t^{T(a)}r_s\dd s}X_a
        \,\middle|\,\cH_t
        \right]
\]
form an arbitrage-free public option manifold.  The directed Volterra kernel \(G^\QQ\) and the
screened Perron operator derived from it define a Volterra--Perron state in the sense of
Definition~\ref{def:q-state}.
\end{theorem}

\begin{proof}
The discounted stock dynamics have zero drift under \(\QQ\), so discounted prices are local
martingales.  Under the true-martingale assumption, conditional-expectation pricing gives an
arbitrage-free price system on the admitted European payoff span.  Square-integrability of the
kernels and the weak solution assumptions place the directed Volterra operator in the Hilbert
space setting used throughout the analysis.  Screening and Perron projection then produce the spectral triple
whenever the positive screened operator has a simple Perron eigenvalue.
\end{proof}

\begin{assumption}[Reduced normal-forward Volterra system]\label{ass:normal-forward-emm}
For each underlier \(i\) and maturity \(T\), let \(F_i(s,T)\), \(0\le s\le T\), be the
\(T\)-forward price under the \(T\)-forward measure \(\QQ^T\).  On the pricing window,
\[
        \dd F_i(s,T)
        =
        \sum_{a=1}^m \sigma_{ia}^{\mathrm N}(s,T,Y_s)\,\dd W_a^{T}(s),
        \qquad 0\le s\le T,
\]
where the normal-forward volatility coefficients are locally square-integrable and adapted.
Their latent stress state is generated by the same directed Volterra source family as in
Assumption~\ref{ass:emm}, possibly after replacing the positive variance transform by a
normal-volatility transform.  The money-market account is strictly positive and the discount
curve is public.  In the finite benchmark, the normal covariance is
\[
        \int_t^T
        \Sigma^{\mathrm N}(s,T,Y_s)
        \Sigma^{\mathrm N}(s,T,Y_s)^\top\,\dd s
        =
        \Gamma_t(T-t;U_t)
\]
after conditioning on \(\cH_t\) and freezing the recovered state at date \(t\), or exactly when
the displayed integrated covariance is \(\cH_t\)-measurable.
\end{assumption}

\begin{proposition}[Normal-forward option manifold under an EMM]\label{prop:normal-forward-emm}
Under Assumption~\ref{ass:normal-forward-emm}, a \(T\)-settled European payoff \(X_T\) has the
arbitrage-free representation
\[
        C_t(X_T)
        =
        P(t,T)\EE^{\QQ^T}\!\left[X_T\mid\cH_t\right].
\]
If a one-period conditional Gaussian approximation is used in the recovered state chart, the
local price coordinates are Bachelier coordinates with normal total variance
\[
        v_{u,t}^{\mathrm N}(T)
        =
        \Var^{\QQ^T}(F_u(T,T)-F_u(t,T)\mid\cH_t),
\]
and the response ranks are the chart-invariant ranks of
Theorem~\ref{thm:pricing-chart-rank-invariance}.
\end{proposition}

\begin{proof}
The forward dynamics have zero drift under \(\QQ^T\), so forward prices are local martingales in
the \(T\)-numeraire.  The change-of-numeraire formula gives
\(C_t=P(t,T)\EE^{\QQ^T}[X_T\mid\cH_t]\) for admitted \(T\)-settled claims under the usual
true-martingale integrability condition.  Conditional Gaussian one-period coordinates reduce
these prices to Definition~\ref{def:black-bachelier-charts} in the normal chart.  The rank
statement follows because the Bachelier chart has strictly positive variance vega away from zero
normal variance.
\end{proof}

\section{Reduced examples}\label{sec:examples}

\subsection{Two-name Gaussian Volterra recovery}

Consider two names with Gaussian variance-stress factors
\[
        V_1(t)=\int_0^t K_1(t-s)\dd B_1(s),\qquad
        V_2(t)=\int_0^t K_2(t-s)\dd B_2(s)+
        \gamma\int_0^t G_{21}(t-s)\dd B_1(s).
\]
The positive variance entering the stock model may be a transform of these factors as in
Assumption~\ref{ass:emm}.  The cross-kernel coefficient \(\gamma\) is latent.  A single
marginal surface for name \(2\)
contains the law of \(V_2\) after public compression, but it does not identify whether the
additional variance comes from an idiosyncratic route, a directed route from name \(1\), or a
public-equivalent residual tilt.  This is the non-identification theorem in its smallest form.

Now impose a Gaussian--Volterra source family in which the singular profile of \(G_{21}\) is known
up to the scalar coefficient and the common-channel sign convention is fixed.  Let
\[
        R_{12}(\tau)
        :=
        \Cov^\QQ\!\left(
        \int_t^{t+\tau}V_1(s)\dd s,
        \int_t^{t+\tau}V_2(s)\dd s
        \,\middle|\,\cH_t
        \right)
\]
be the cross-maturity response extracted from the index/single-name quadratic option block.  In
the source family,
\[
        R_{12}(\tau)
        =
        \gamma\,q_{21}(\tau)+r_{\mathrm{low}}(\tau),
\]
where \(q_{21}\) is the leading Volterra singular-profile response and
\(r_{\mathrm{low}}\) is lower order uniformly on the maturity interval.  If \(q_{21}\) is not
identically zero, the map \(\gamma\mapsto R_{12}\) has nonzero derivative in the quotient chart.
An open maturity interval, or a finite analytic maturity panel as in
Theorem~\ref{thm:finite-panel-recovery}, therefore recovers \(\gamma\) up to the fixed gauge
convention.  Public compression alone fails; source-family maturity structure supplies the
required rank.

\subsection{Index plus carrier manifold}

Let an index have weights \(\omega_i\).  The implied index total variance can be written in the
schematic form
\[
        v_I^{\mathrm{imp}}(\tau)
        =
        \sum_i \omega_i^2 v_i^{\mathrm{imp}}(\tau)
        +
        D_\omega(v^{\mathrm{imp}}(\tau))\,
        \rho_{\mathrm{sync}}^{\mathrm{imp}}(\tau),
\]
where \(D_\omega\) is the public synchronization denominator.  If the carrier surface span is
nondegenerate, with Gram matrix \(\Gamma_{\mathrm{car}}\succeq\kappa_{\mathrm{car}}I\), the
observed family
\[
        \tau\mapsto
        \bigl(v_I^{\mathrm{imp}}(\tau),v_i^{\mathrm{imp}}(\tau)\bigr)
\]
over an open maturity interval identifies the implied synchronization curve
\(\rho_{\mathrm{sync}}^{\mathrm{imp}}(\tau)\).  A Volterra--Perron source family then represents
that curve as the public image of \((\rho^\QQ,e^\QQ,\ell^\QQ,P^\QQ)\).  The carrier Gram lower
bound is what prevents the single-name side from becoming an ill-conditioned explanation of the
index response.

The remaining price difference between the index claim and its carrier-replicated single-name
claim decomposes as
\[
        \Pi_t^\theta(X_I)-\Pi_t^\theta(X_{\mathrm{car}})
        =
        \mathrm{SynPrem}_t(X_I)
        +
        \mathrm{ResCharge}_t^\theta(X_I)
        +
        \Pi_t^\theta(B_t(X_I)).
\]
The first term is the public synchronized-stress premium read from the option manifold.  The
second is the selected residual compression charge attached to unspanned latent stress.  The
third is the tradeable basis error left by the carrier span.  Thus the same manifold that
recovers the canonical synchronization state also tells the model where recovery stops and
pricing selection begins.

\subsection{Three-name carrier-rank example}

For \(N=3\), write the active unordered edge vector as
\[
        V_{\cE}(\tau)
        =
        \bigl(V_{12,t}^\QQ(\tau),V_{13,t}^\QQ(\tau),V_{23,t}^\QQ(\tau)\bigr)^\top .
\]
One broad index with weights \(w=(w_1,w_2,w_3)\) encodes only
\[
        B_I(\tau)
        =
        2w_1w_2V_{12,t}^\QQ(\tau)
        +2w_1w_3V_{13,t}^\QQ(\tau)
        +2w_2w_3V_{23,t}^\QQ(\tau),
\]
which is one equation for three cross routes.  Adding three pair carriers with weights
\[
        b^{12}=(1,1,0),\qquad
        b^{13}=(1,0,1),\qquad
        b^{23}=(0,1,1)
\]
gives a basket design matrix proportional to the identity on the three active pairs.  The three
cross routes are then identified at each maturity.  Carrier rank is the algebraic condition that
separates synchronization routes rather than only their weighted aggregate.

\subsection{Directed-route seam example}

With two names, suppose the ordered routes enter a public cross response as
\[
        R_{12}(\tau)
        =
        \gamma_{1\leftarrow2}f(\tau)
        +
        \gamma_{2\leftarrow1}f(\tau).
\]
Then only the sum \(\gamma_{1\leftarrow2}+\gamma_{2\leftarrow1}\) is identified.  If instead
\[
        R_{12}(\tau)
        =
        \gamma_{1\leftarrow2}f_{12}(\tau)
        +
        \gamma_{2\leftarrow1}f_{21}(\tau),
\]
and \(f_{12}\) and \(f_{21}\) are linearly independent analytic maturity functions, an open
interval or a generic two-point panel separates the two directions.  This is the smallest
illustration of the direction-resolving block in
Theorem~\ref{thm:orientation-separation}.

\subsection{Residual-envelope calculation}

Let \(R^1,R^2\) be conditionally orthonormal residual directions, so \(G_t^R=I_2\), and suppose
\[
        g_t^X=(g_1,g_2)^\top .
\]
Under the ellipsoidal residual budget \(\theta_1^2+\theta_2^2\le\rho^2\),
Theorem~\ref{thm:ellipsoidal-residual-envelope} gives
\[
        \left[
        \EE^\QQ[X\mid\cH_t]-\rho\sqrt{g_1^2+g_2^2},\,
        \EE^\QQ[X\mid\cH_t]+\rho\sqrt{g_1^2+g_2^2}
        \right].
\]
Under the box budget \(|\theta_i|\le\bar\theta_i\), Corollary~\ref{cor:box-residual-envelope}
gives
\[
        \left[
        \EE^\QQ[X\mid\cH_t]-\bar\theta_1|g_1|-\bar\theta_2|g_2|,\,
        \EE^\QQ[X\mid\cH_t]+\bar\theta_1|g_1|+\bar\theta_2|g_2|
        \right].
\]
The residual compression charge is the support function of the admissible selector set evaluated
at the payoff-residual vector.  The calculation is a pricing-fiber computation, not a recovery
claim.

\section{Conclusion}\label{sec:conclusion}

The option manifold identifies a quotient, not the full latent risk-neutral state.  The governing
map is
\[
        \text{public option prices}\rightsquigarrow\mathfrak X_t^{\mathrm{pub}},
        \qquad
        \text{state recovery}=\mathfrak S_t/\!\sim_g\hookrightarrow
        \mathfrak X_t^{\mathrm{pub}} .
\]
Stable recovery is positivity of the quotient observability modulus.  Residual prices collapse to
singletons exactly at public payoff-module completion.  Dynamic admissibility adds convex-order
and recovered martingale drift restrictions.  Recovery, completion, and hedgeability remain separate.

The same quotient records the failure modes.  Marginal surfaces leave synchronization unresolved.
Finite basket panels leave elliptope uncertainty unless their design spans the active covariance
routes.  Carrier Gram collapse makes public-neutral dispersion unstable.  Perron reduced-resolvent
blow-up, including gap collapse, makes the dominant stress mode non-smooth.  Maturity-rank
collapse prevents Volterra source identification.  Public payoff incompletion leaves residual
pricing fibers.  Convex-order or martingale-drift failure rules out a dynamically admissible
risk-neutral model.  These are local non-identification mechanisms in the option manifold, not
estimation artifacts.

Once recovered and paired with an admissible residual selector, the state supports selected prices
beyond the public surface.  Native Greeks measure sensitivity to Perron stress, spectral radius,
carrier direction, roughness, directed contagion, and residual selection.  The synchronization
relative-value identity expresses the gap between index options and carrier-replicated single-name
options as synchronized-stress premium, residual compression charge, and basis error.  The finite
Bachelier--Volterra--Perron benchmark gives a minimal pricing equation for this geometry: a
recovered finite state, a closed option formula, native state Greeks, and a computable public hedge
residual.  The benchmark is the finite chart in which the quotient recovery theory becomes a
tractable no-arbitrage pricing object, rather than a full option-manifold model.

\bibliographystyle{unsrt}
\bibliography{references}

\end{document}
